\chapter{Bloch-Kato Selmer groups, Tamagawa factors and the adjoint motive}\label{ch:blochkato}
\markboth{\MakeUppercase{\chaptername\ \thechapter.\ Bloch-Kato Selmer groups and Tamagawa factors}}{}

Theorem thm:main of the paper says that $v_\ell(\degphi)$ is the length of a Bloch-Kato Selmer group plus Tamagawa factors plus omitted Euler factors. This chapter
defines the first two, computes the second at every bad prime, and proves the Selmer half of the theorem: the comparison of the Bloch-Kato Selmer group with the
$\Sigma$-imprimitive one of Diamond, Flach and Guo, with every local term identified. It proves the paper's Lemmas lem:fpr, lem:tamfactor and lem:tamadd and its
Corollary cor:explicit, restates Lemma lem:dfg21 and Lemma lem:h0 in the language of Selmer groups, and places the whole in the framework of the Tamagawa number
conjecture of Bloch and Kato.

\section*{Core results at a glance}

Throughout, $E/\Q$ has conductor $N$, $\ell$ is an odd prime, $T=T_\ell E$, $V=T\otimes\Ql$, $A=\ad T$, $A_f=\ad V$, $W=A\otimes\Qlmod$, $T^D=\Hom(A,\Zl(1))$, $V^D=T^D\otimes\Ql$, $W^D=V^D/T^D$, and $\Sigma$ is
the set of primes dividing $N$; $p\ne\ell$ in all local statements unless said otherwise.

\begin{enumerate}[label=\textbf{\thechapter.\arabic*}]
\item \textbf{Conventions.} Linear dual $M^\vee$, Kummer dual $M^D=\Hom(M,\Zl(1))=M^\vee(1)$, Pontryagin dual $X^*$; $\Frob_p$ arithmetic, $L_p(M,s)=\det(1-\Frob_p^{-1}p^{-s}\mid M^{I_p})^{-1}$
with the geometric Frobenius; $\Ql(1)$ has Hodge-Tate weight $+1$; Tate data $q_E$, $\nu=-v_q(j_E)$, $u=q_Eq^{-\nu}$, $c(u)$ the class of $\bar u$ in $\F_q^\times/(\F_q^\times)^{\ell^{v_\ell(q-1)}}$, $v_\ell(c(u))$ the
largest $k\le v_\ell(q-1)$ with $\bar u$ an $\ell^k$-th power. \emph{(convention)} \Cref{conv:blochkato:setup}.
\item \textbf{$\ell$-adic representations.} A continuous $G_K\to\GL(V)$ on a finite-dimensional $\Ql$-space; a stable lattice $T$; $W=V/T$. \emph{(definition)}
\Cref{def:blochkato:l-adic-rep}.
\item \textbf{Lattices.} Stable lattices exist; $W[\ell^k]\simeq T/\ell^kT$; lattices are commensurable, and all stable lattices are homothetic if $T/\ell T$ is irreducible. \emph{(proved
here)} \Cref{prop:blochkato:lattices}.
\item \textbf{Tate twists.} $\Zl(1)=\varprojlim\mu_{\ell^k}$, $M(n)=M\otimes\Zl(1)^{\otimes n}$. \emph{(definition)} \Cref{def:blochkato:twists}.
\item \textbf{Twists at $p\ne\ell$.} $M(n)^{I_p}=M^{I_p}(n)$ and Frobenius eigenvalues are multiplied by $p^n$ (arithmetic) or $p^{-n}$ (geometric); $\wedge^2T_\ell E\simeq\Zl(1)$, $(T_\ell E)^\vee\simeq T_\ell E(-1)$.
\emph{(proved here, given \Cref{ch:ellcurves})} \Cref{prop:blochkato:twists}.
\item \textbf{Kummer duals.} $V^D=V^\vee(1)$, $T^D=T^\vee(1)$, $W^D=V^D/T^D$; $X^*=\Hom_{\mathrm{cont}}(X,\Qlmod)$. \emph{(definition)} \Cref{def:blochkato:duals}.
\item \textbf{Kummer duality.} $W^D\simeq\Hom(T,\mu_{\ell^\infty})$, $T^D\simeq\Hom(W,\mu_{\ell^\infty})$, perfect equivariant pairings; $(T^D)^D=T$; $X^*$ has the length of $X$ for finite $X$.
\emph{(proved here, given \Cref{ch:modules})} \Cref{prop:blochkato:duals}.
\item \textbf{The adjoint is self-dual.} For $\ell$ odd the trace form is perfect on $A$ (Gram determinant $-2$), so $T^D\simeq A(1)$, $V^D\simeq A_f(1)$, $W^D\simeq\ad T(1)\otimes\Qlmod$;
$\ad$ is unchanged by twisting by a character, by homothety and by dualising; so $\ad M_{f,\lambda}\simeq\ad T_\ell E$, and $\ad T$ is the same $G_{\Q_q}$-module for $E$ and its twists.
\emph{(proved here, given \Cref{ch:reps,ch:moddeg})} \Cref{prop:blochkato:adjoint-self-dual}.
\item \textbf{Geometric representations.} Unramified almost everywhere and de Rham at $\ell$. \emph{(definition)} \Cref{def:blochkato:geometric}.
\item \textbf{Purity.} $V_\ell E$ and its tensor constructions are geometric \emph{(stated; reference)}; at $p\nmid N\ell$ the Frobenius eigenvalues on $V_\ell E$, $A_f$, $V^D$ have absolute values
$p^{1/2}$, $1$, $p$ \emph{(proved here, given \Cref{ch:ellcurves})}. \Cref{prop:blochkato:purity}.
\item \sloppy \textbf{Global invariants.} If $\bar\rho$ is absolutely irreducible, $H^0(\Q,A_f)=H^0(\Q,W)=0$; $H^0(\Q,V^D)=0$ for every $E$; $A_f$, $A_f(1)$, $A_f(-1)$ are pairwise non-isomorphic.
\emph{(proved here)} \Cref{prop:blochkato:global-invariants}.
\item \textbf{The Frobenius dictionary.} $L_p(M(1),s)=L_p(M,s+1)$; $(F-1)X=(F^{-1}-1)X$; $\#X/(F-1)X=\ell^{v_\ell\det(1-F^{-1}\mid X)}$ for a lattice $X$; $\det(1-\Frob_p^{-1}p^{-s}\mid V_\ell E)^{-1}=L_p(E,s+1)$.
\emph{(proved here)} \Cref{prop:blochkato:frobenius-dictionary}.
\item \textbf{The adjoint at $p\ne\ell$.} $A_f^{I_p}$, the geometric Frobenius on it and $L_p(A_f,1)^{-1}=\det(1-\Frob_p^{-1}\mid(V^D)^{I_p})$ by reduction type: $1-p^{-2}$ at potentially multiplicative primes,
$(1-\frac{\alpha}{\beta p})(1-\frac1p)(1-\frac{\beta}{\alpha p})$ at good primes, $1-cp^{-1}$ for cyclic inertia of order $3,4,6$, $\frac{p-1}p\cdot\frac{(p+1)^2-a^2}{p^2}$ for $\rho(I_p)=\{\pm1\}$, $1$ for non-abelian
inertia. \emph{(proved in \Cref{ch:galrep,ch:lfunctions}; recalled here with the reasons)} \Cref{prop:blochkato:local-adjoint}.
\item \textbf{Local invariants.} $H^0(\Q_p,V^D)=0$ and $L_p(A_f,1)^{-1}\in\Q^\times$; $H^0(\Q_p,A_f)$ is $0$ or a line by type, nonzero at some bad primes. \emph{(proved here)}
\Cref{lem:blochkato:no-local-invariants}.
\item \textbf{Local finiteness.} $H^1(\Q_p,T^D)$ has rank $\dim H^0(\Q_p,A_f)$; it is finite when that is $0$; its torsion is always finite. \emph{(proved here, given \Cref{ch:galcoh})}
\Cref{cor:blochkato:local-finiteness}.
\item \textbf{Integrality.} $L_p(A_f,1)^{-1}=\det(1-\Frob_p^{-1}\mid(T^D)^{I_p})\in\Z_{(\ell)}\cap\Q^\times$, $\#(T^D)^{I_p}/(\Frob_p-1)=\ell^{v_\ell(L_p(A_f,1)^{-1})}$; the valuation is $v_\ell(p^2-1)$ at potentially
multiplicative primes and $e_\ell(E,p)$ at additive primes. \emph{(the lattice statements proved here; the value and valuation proved in \Cref{ch:lfunctions})} \Cref{lem:blochkato:euler-integral}.
\item \textbf{Period rings.} $\BdR\supset\Bst\supset\Bcris$, $D_\bullet(V)=(B_\bullet\otimes V)^{G_{\Q_\ell}}$; de Rham, semistable, crystalline representations; Hodge-Tate weights. \emph{(stated; reference)}
\Cref{def:blochkato:period-rings}.
\item \textbf{Good reduction.} $V_\ell B$ is crystalline iff the abelian variety $B/\Q_\ell$ has good reduction; for $E$ good at $\ell$, $\Dcris((V_\ell E)^\vee)$ has Frobenius polynomial
$X^2-a_\ell X+\ell$. \emph{(stated; reference)} \Cref{thm:blochkato:crystalline-good-reduction}.
\item \textbf{Stability.} Crystalline and de Rham representations are stable under the tensor operations; for $\ell\nmid N$, $A_f$, $V^D$, $V^\vee$ are crystalline with weights $\{-1,0,1\}$,
$\{0,1,2\}$, $\{-1,0\}$. \emph{(stated; reference; the consequences proved here)} \Cref{prop:blochkato:crystalline-closure}.
\item \textbf{Fontaine-Laffaille.} An exact equivalence between strongly divisible filtered modules with weights in an interval of length at most $\ell-2$ and lattices in
crystalline representations. \emph{(stated; reference)} \Cref{thm:blochkato:fontaine-laffaille}.
\item \sloppy \textbf{Local dimensions at $\ell$.} $\dim H^1_f(\Q_\ell,V)=\dim H^0(\Q_\ell,V)+\#\{\text{positive Hodge-Tate weights}\}$; for $\ell\nmid N$, $\dim H^1_f(\Q_\ell,A_f)=\dim H^0(\Q_\ell,A_f)+1$ and
$\dim H^1_f(\Q_\ell,V^D)=2$. \emph{(proof sketched here; reference)} \Cref{prop:blochkato:dimension-formula}.
\item \textbf{Bloch-Kato local conditions.} $H^1_f(\Q_p,V)=H^1_{\ur}(\Q_p,V)$ for $p\ne\ell$; $H^1_f(\Q_\ell,V)=\ker(H^1(\Q_\ell,V)\to H^1(\Q_\ell,\Bcris\otimes V))$, with $H^1_e\subseteq H^1_f\subseteq H^1_g$;
$H^1_f(\R,V)=0$. \emph{(definition)} \Cref{def:blochkato:local-conditions}.
\item \textbf{Propagation.} $H^1_f(\Q_p,T)=\iota^{-1}H^1_f(\Q_p,V)$, $H^1_f(\Q_p,W)=\pi H^1_f(\Q_p,V)$. \emph{(definition)} \Cref{def:blochkato:propagation}.
\item \sloppy \textbf{Propagated conditions.} $H^1_f(\Q_p,T)\supseteq H^1(\Q_p,T)_{\mathrm{tor}}=\ker\iota$; $H^1_f(\Q_p,W)\simeq(\Qlmod)^{\dim H^1_f(\Q_p,V)}$. \emph{(proved here, given \Cref{ch:galcoh})}
\Cref{prop:blochkato:propagation}.
\item \textbf{The conditions at $p\ne\ell$.} $H^1_f(\Q_p,V)=V^{I_p}/(\Frob_p-1)$; $0\to H^1_{\ur}(\Q_p,T)\to H^1(\Q_p,T)\to H^1(I_p,T)^{G_{\Q_p}}\to0$; $H^1_f(\Q_p,T)/H^1_{\ur}(\Q_p,T)\simeq(H^1(I_p,T)_{\mathrm{tor}})^{G_{\Q_p}}$;
$H^1_f(\Q_p,W)$ is the maximal divisible subgroup of $H^1_{\ur}(\Q_p,W)$; if $H^0(\Q_p,V)=0$ then $H^1_f(\Q_p,T)=H^1(\Q_p,T)_{\mathrm{tor}}$ is finite. \emph{(proved here, given
\Cref{ch:cohomology,ch:galcoh})} \Cref{prop:blochkato:local-unramified}.
\item \textbf{Duality.} $H^1_f(\Q_v,V)$ and $H^1_f(\Q_v,V^D)$ are exact annihilators; $H^1(\Q_v,W)/H^1_f(\Q_v,W)\simeq H^1_f(\Q_v,T^D)^*$. \emph{(proved here for $v\ne\ell$, given \Cref{ch:galcoh};
stated with reference at $v=\ell$)} \Cref{prop:blochkato:local-duality-f}.
\item \textbf{Bloch-Kato Selmer groups.} $H^1_f(\Q,M)$: classes whose localisations lie in $H^1_f(\Q_v,M)$ at every place. \emph{(definition)} \Cref{def:blochkato:selmer}.
\item \textbf{Finite level.} $H^1_f=H^1_{\ur}$ at unramified primes, and $H^1_f(\Q,M)\subseteq H^1(G_{\Q,S},M)$ is cut out by the conditions at $S$. \emph{(proved here)} \Cref{lem:blochkato:selmer-S}.
\item \textbf{Relations.} $H^1_f(\Q,T)=\iota^{-1}H^1_f(\Q,V)$ with torsion the image of $H^0(\Q,W)$; $H^1_f(\Q,W)$ is cofinitely generated with maximal divisible subgroup $\pi H^1_f(\Q,V)$, so it is finite
iff $H^1_f(\Q,V)=0$. \emph{(proved here, given the finiteness theorems of \Cref{ch:galcoh})} \Cref{prop:blochkato:selmer-relations}.
\item \textbf{The $\Sigma$-imprimitive Selmer group.} $H^1_\Sigma(\Q,W)$: Bloch-Kato at $\ell$, unramified outside $\Sigma\cup\{\ell\}$, nothing on $\Sigma$. \emph{(definition)}
\Cref{def:blochkato:sigma-selmer}.
\item \textbf{Finiteness of $H^1_\Sigma$.} If $\Sigma$ ($\ell\notin\Sigma$) contains the primes $p\ne\ell$ at which $W$ is ramified, $H^1_f(\Q,V)=0$ and $H^0(\Q_p,V^D)=0$ for $p\in\Sigma$, then $0\to H^1_f(\Q,W)\to H^1_\Sigma(\Q,W)\to\bigoplus_{p\in\Sigma}H^1_f(\Q_p,T^D)^*$ is exact with finite outer
terms. \emph{(proved here)} \Cref{prop:blochkato:sigma-finite}.
\item \textbf{The paper's Selmer groups.} $H^1_f(\Q,\ad T\otimes\Qlmod)$ and $H^1_\Sigma(\Q,\ad T\otimes\Qlmod)$. \emph{(definition)} \Cref{def:blochkato:adjoint-selmer}.
\item \textbf{Euler characteristic.} $\dim H^1_f(\Q,V)-\dim H^1_f(\Q,V^D)=\dim H^0(\Q,V)-\dim H^0(\Q,V^D)+\sum_v(\dim H^1_f(\Q_v,V)-\dim H^0(\Q_v,V))$ \emph{(proof sketched here; reference)}; for $A_f$ the
right side is $0$ when $\ell\nmid N$ and $\bar\rho$ is absolutely irreducible, so then $H^1_f(\Q,A_f)=0\iff H^1_f(\Q,V^D)=0$ \emph{(proved here, given the formula)}. \Cref{prop:blochkato:euler-characteristic}.
\item \textbf{Tamagawa factor.} $\Tam_p(M)=\Fitt_{\Zl}\bigl((H^1(I_p,M)_{\mathrm{tor}})^{G_{\Q_p}}\bigr)$, the Fitting ideal of the invariants. \emph{(definition)} \Cref{def:blochkato:tamagawa}.
\item \textbf{Index description.} $\ell^{v_\ell(\Tam_p(M))}=[H^1_f(\Q_p,M):H^1_{\ur}(\Q_p,M)]$. \emph{(proved here)} \Cref{prop:blochkato:tam-index}.
\item \textbf{Basic properties.} $\Tam_p(M)=\Zl$ if $M$ is unramified or inertia acts through a finite group of order prime to $\ell$; $\Tam_p(\ad T)=\Tam_p(\ad(T\otimes\chi))$; $\Tam_p(\ad T_\ell E)=\Zl$ at good
primes and when $\rho(I_p)=\{\pm1\}$. \emph{(proved here)} \Cref{prop:blochkato:tam-basic}.
\item \textbf{Fontaine and Perrin-Riou (Lemma lem:fpr).} If $H^0(\Q_p,X\otimes\Ql)=0$ then $0\to X^{I_p}/(\Frob_p-1)\to H^1(\Q_p,X)_{\mathrm{tor}}\to(H^1(I_p,X)_{\mathrm{tor}})^{G_{\Q_p}}\to0$ is exact and
$\Fitt H^1_f(\Q_p,X)=\det(1-\Frob_p^{-1}\mid X^{I_p})\Tam_p(X)$; for $X=T^D$, $\Fitt H^1_f(\Q_p,T^D)=L_p(A_f,1)^{-1}\Tam_p(T^D)$ and $\len H^1_f(\Q_p,T^D)=\len H^0(\Q_p,W^D)$. \emph{(proved here, given
\Cref{ch:galcoh})} \Cref{thm:blochkato:fpr-local}.
\item \textbf{$\Tam_p(T^D)=\Tam_p(A)$.} For the adjoint of $E$ at every $p\ne\ell$; at potentially multiplicative primes both invariant modules are $\Z/\ell^a\oplus\Z/\ell^m$, with $a=v_\ell(\nu)$ and $m=\min(a,v_\ell(p-1),v_\ell(c(u)))$.
\emph{(proved here, given \Cref{ch:cohomology,ch:galcoh})}
\Cref{prop:blochkato:tam-dual}.
\item \textbf{Tamagawa numbers.} $v_\ell(\Tam_p(T_\ell E))=v_\ell(c_p)$. \emph{(proved here, given \Cref{ch:reduction})} \Cref{prop:blochkato:tam-of-E}.
\item \textbf{Tate data.} The Tate shape $\sm{\chi_\ell\varepsilon}{\varepsilon\kappa}{0}{\varepsilon}$, $\kappa|_I=n\,t_\ell$, $\Frob\,t\,\Frob^{-1}=t^q$, $\kappa=n\kappa_q+\kappa_u$, $[\kappa_u]\leftrightarrow\bar u$. \emph{(convention;
recollection from \Cref{ch:local,ch:galrep,ch:galcoh})} \Cref{conv:blochkato:tate-data}.
\item \textbf{Lemma lem:tamfactor.} At a multiplicative prime $q\ne\ell$: $H^1(I_q,\ad T)_{\mathrm{tor}}\simeq(\Z/\ell^{v_\ell(n)})^2$ and
\[
v_\ell(\Tam_q(\ad T))=v_\ell(n)+\min\bigl(v_\ell(n),v_\ell(q-1),v_\ell(c(u))\bigr)\ \ge\ t_\ell(E,q).
\]
\emph{(proved here, given \Cref{ch:galrep,ch:cohomology,ch:galcoh})} \Cref{lem:blochkato:tamfactor}.
\item \textbf{Consequences.} The excess $\min(\dots)$ is positive iff $\ell\mid n$, $q\equiv1\pmod\ell$ and $\bar u\in(\F_q^\times)^\ell$; it is $0$ at $q=2$ and independent of split or nonsplit. \emph{(proved
here)} \Cref{cor:blochkato:tamfactor-consequences}.
\item \sloppy \textbf{Potentially multiplicative additive primes.} At $q\ne\ell$, $q=2$ included: $v_\ell(\Tam_q(\ad T))=v_\ell(\nu)+\min(v_\ell(\nu),v_\ell(q-1),v_\ell(c(u)))$. \emph{(proved here, given
\Cref{ch:reduction,ch:galrep})} \Cref{prop:blochkato:tam-potmult}.
\item \textbf{Examples.} Tate data and $v_\ell(\Tam_q(\ad T))$ for \curve{11a1}, \curve{14a1}, \curve{539c1}, \curve{791a1}, \curve{1184f1}, \curve{1748e1}. \emph{(computed)}
\Cref{ex:blochkato:tamfactor-examples}.
\item \textbf{Finite inertia image.} If $I_q$ acts through a finite $\Delta$, $H^1(I_q,M)_{\mathrm{tor}}=H^1(\Delta,M)$. \emph{(proved here)} \Cref{lem:blochkato:finite-inertia-image}.
\item \sloppy \textbf{$\ad T$ is free over $\Z_3[C_3]$.} For $\sigma\in\GL_2(\Z_3)$ of order $3$, $\ad T\simeq\Z_3[\langle\sigma\rangle]$, so $H^1(\langle\sigma\rangle,\ad T)=0$, while $H^1(\langle\sigma\rangle,T)\simeq\F_3$. \emph{(proved here)}
\Cref{prop:blochkato:ad-free-c3}.
\item \textbf{Additive primes (Lemma lem:tamadd).} For $\ell$ odd and $q\ne\ell$ of potentially good additive reduction, $H^1(I_q,\ad T)_{\mathrm{tor}}=0$ and $\Tam_q(\ad T)=\Tam_q(\ad T(1))=\Zl$.
\emph{(proved here, given \Cref{ch:galrep})} \Cref{thm:blochkato:tam-additive}.
\item \sloppy \textbf{The explicit identity (Corollary cor:explicit).} $\sum_{q\mid N}v_\ell(\Tam_q(\ad T))=\sum_{v_q(j)<0}\bigl(t_\ell+\min(t_\ell,v_\ell(q-1),v_\ell(c(u_q)))\bigr)$ \emph{(proved here)}; hence, under the
hypothesis of Theorem thm:main, $v_\ell(\degphi)=\len H^1_f(\Q,W)+\sum_qe_\ell(E,q)+\sum_{v_q(j)<0}\bigl(t_\ell+\min(\dots)\bigr)$ \emph{(given Theorem thm:main, \Cref{ch:congruence})}.
\Cref{cor:blochkato:explicit-identity}.
\item \textbf{The excess on the tables.} Below conductor $10{,}000$ the excess is never negative and is $0$ for $8{,}688$ of $13{,}203$, $19{,}019$ of $23{,}361$, $25{,}133$ of $28{,}134$ curves at
$\ell=3,5,7$. \emph{(computed)} \Cref{rem:blochkato:open-question-3}.
\item \textbf{The exceptional set.} $S_f=\{\ell:\ell\mid2N$ or $\bar\rho|_{\Gal(\Qbar/\Q(\sqrt{\ell^*}))}$ not absolutely irreducible$\}$. \emph{(definition)} \Cref{def:blochkato:exceptional-set}.
\item \textbf{Theorem thm:dfg27.} If $\ell\notin S_f$ then $H^0(\Q,V^D)=H^1_f(\Q,V^D)=0=H^1_f(\Q,A_f)$. \emph{(stated; reference; a derivation from Theorem thm:dfg37 sketched)}
\Cref{thm:blochkato:dfg27}.
\item \textbf{Lemma lem:h0 in Selmer form.} If $H^0(\Q,V^D)=H^1_f(\Q,V^D)=0$, $H^0(\Q,W^D)\simeq H^1_f(\Q,T^D)=H^1(\Q,T^D)_{\mathrm{tor}}$ \emph{(proved here)}; $H^0(\Q,W^D)\simeq\Z/\ell^{h_\ell(E)}$, $h_\ell=0$ for
$\ell\ge5$, $h_3\le1$ with equality iff the image is $N_s(\F_3)$ \emph{(proved in \Cref{ch:galrep})}. \Cref{prop:blochkato:h0-selmer}.
\item \textbf{Lemma lem:dfg21.} If $H^0(\Q,V^D)=H^1_f(\Q,V^D)=0$ and $\Sigma\ne\emptyset$, $\ell\notin\Sigma$, contains the ramified primes $p\ne\ell$:
\[
0\to H^1_f(\Q,W)\to H^1_\Sigma(\Q,W)\to\bigoplus_{p\in\Sigma\cup\{\infty\}}H^1_f(\Q_p,T^D)^*\to H^0(\Q,W^D)^*\to0.
\]
\emph{(proved here, given \Cref{ch:galcoh})} \Cref{lem:blochkato:dfg21}.
\item \textbf{The length identity.} Under these hypotheses,
$\len H^1_\Sigma(\Q,W)=\len H^1_f(\Q,W)+\sum_{p\in\Sigma}\bigl(v_\ell(L_p(A_f,1)^{-1})+v_\ell(\Tam_p(\ad T))\bigr)-h_\ell(E)$. \emph{(proved here, given \Cref{prop:blochkato:euler-characteristic}
and \Cref{ch:galcoh,ch:lfunctions})} \Cref{cor:blochkato:length-identity}.
\item \textbf{The Selmer identity.} For $\ell\notin S_f$, $H^1_f(\Q,W)$ and $H^1_\Sigma(\Q,W)$ are finite and
\begin{multline*}
\len H^1_\Sigma(\Q,W)=\len H^1_f(\Q,W)+\sum_{q\mid N}v_\ell\bigl(L_q(A_f,1)^{-1}\bigr)\\
+\sum_{v_q(j)<0}\Bigl(t_\ell(E,q)+\min\bigl(t_\ell(E,q),v_\ell(q-1),v_\ell(c(u_q))\bigr)\Bigr).
\end{multline*}
\emph{(proved here, given Theorem thm:dfg27)} \Cref{thm:blochkato:selmer-identity}.
\item \textbf{Remark rem:ns.} For $\ell=3$, image $N_s(\F_3)$, $3\nmid N$: if $H^1_f(\Q,V^D)=0$, the length identity holds with $-h_3(E)=-1$. \emph{(proved here, given \Cref{prop:blochkato:euler-characteristic}; conditional)}
\Cref{prop:blochkato:rem-ns}.
\item \textbf{Section sec:tam.} Below conductor $10{,}000$ the refined inequality has no exception; equality rises from $596$ to $714$ of $1{,}050$ curves at $\ell=3$ and from $203$ to $220$ of $252$ at
$\ell=5$. \emph{(computed)} \Cref{ex:blochkato:sec-tam}.
\item \textbf{The adjoint motive.} Realisations $A_f$, Hodge numbers $h^{1,-1}=h^{0,0}=h^{-1,1}=1$, $L(\ad f,s)=L(\Sym f,s+1)$; $s=1$ is critical and $L(\ad f,1)\ne0$. \emph{(definition; properties
stated with reference, the nonvanishing given \Cref{ch:lfunctions})} \Cref{def:blochkato:adjoint-motive}.
\item \textbf{The Tamagawa number conjecture.} Schematically,
\[
\frac{L(M,0)}{\Omega(M)}\sim_\ell\frac{\prod_{p\ne\ell}\#(\Zl/\Tam_p(T))\cdot\#H^1_f(\Q,W_M^D)}{\#H^0(\Q,W_M)\cdot\#H^0(\Q,W_M^D)}.
\] \emph{(stated in outline;
reference)} \Cref{conj:blochkato:tnc}.
\item \textbf{Diamond, Flach and Guo, Theorem 0.1.} For $\lambda\notin S_f$ the $\lambda$-part of the conjecture holds for the adjoint motive at $s=1$. \emph{(stated; reference)} \Cref{thm:blochkato:dfg-01}.
\item \textbf{Diamond, Flach and Guo, Theorem 0.2 (thm:dfg37).} For $\ell\notin S_f$, $H^1_\Sigma(\Q,W)$ is finite of length $v_\lambda(\eta^\Sigma_{f,\lambda})$. \emph{(stated; reference; see \Cref{ch:congruence})}
\Cref{thm:blochkato:dfg-02}.
\item \textbf{Flach.} For $\ell$ outside an explicit finite set the Bloch-Kato Selmer group of $\Sym T_\ell E\otimes\Qlmod$ ($=W^D$) is finite, annihilated by an integer built from $\degphi$. \emph{(stated in
outline; reference)} \Cref{thm:blochkato:flach}.
\end{enumerate}

\section*{Where the paper uses this chapter}

\begin{itemize}
\item \textbf{Section 1.2} (the definition of $\Tam_p(M)$; ``at a prime of good reduction it is $0$''; the Bloch-Kato Selmer group): \Cref{def:blochkato:tamagawa},
\Cref{prop:blochkato:tam-basic}, \Cref{def:blochkato:selmer}, \Cref{def:blochkato:adjoint-selmer}.
\item \textbf{Section 1.1} (the geometric Frobenius in $L_p(A_f,s)$): \Cref{prop:blochkato:frobenius-dictionary}, \Cref{prop:blochkato:local-adjoint}.
\item \textbf{Section 4.2} ($A$, $W$, $T^D\simeq A(1)$ via the trace form, $M_{f,\lambda}$ and its adjoint, $H^1_f$, $H^1_\Sigma$, $S_f$): \Cref{prop:blochkato:adjoint-self-dual},
\Cref{def:blochkato:sigma-selmer}, \Cref{def:blochkato:exceptional-set}.
\item \textbf{Theorem thm:dfg27}: \Cref{thm:blochkato:dfg27}, \Cref{prop:blochkato:euler-characteristic}, \Cref{prop:blochkato:global-invariants}.
\item \textbf{Lemma lem:dfg21 and the paragraph after it}: \Cref{lem:blochkato:dfg21}, \Cref{prop:blochkato:local-unramified}, \Cref{prop:blochkato:local-duality-f},
\Cref{lem:blochkato:no-local-invariants}, \Cref{cor:blochkato:local-finiteness}.
\item \textbf{Lemma lem:fpr}: \Cref{thm:blochkato:fpr-local}, \Cref{prop:blochkato:tam-dual}, \Cref{lem:blochkato:euler-integral}, \Cref{prop:blochkato:frobenius-dictionary}.
\item \textbf{Lemma lem:h0} (as it is used in Selmer groups): \Cref{prop:blochkato:h0-selmer}.
\item \textbf{Lemma lem:tamfactor} (all of it, including the additive potentially multiplicative case and $q=2$): \Cref{conv:blochkato:tate-data}, \Cref{lem:blochkato:tamfactor},
\Cref{prop:blochkato:tam-potmult}, \Cref{prop:blochkato:tam-of-E}.
\item \textbf{Lemma lem:tamadd and Corollary cor:explicit}: \Cref{lem:blochkato:finite-inertia-image}, \Cref{prop:blochkato:ad-free-c3}, \Cref{thm:blochkato:tam-additive},
\Cref{cor:blochkato:explicit-identity}.
\item \textbf{The proof of Theorem thm:main} (the Selmer bookkeeping from Theorem thm:dfg27 to the local terms; the potentially multiplicative primes): \Cref{cor:blochkato:length-identity},
\Cref{thm:blochkato:selmer-identity}, \Cref{rem:blochkato:assembly}.
\item \textbf{Remark rem:ko(1)} and the abstract's sentence on the Tamagawa factor: \Cref{cor:blochkato:tamfactor-consequences}.
\item \textbf{Remark rem:ns, Section 1.3, Conjecture conj:euler} (the term $h_\ell(E)$), \textbf{Open question 2}: \Cref{prop:blochkato:rem-ns}, \Cref{rem:blochkato:conj-euler-h},
\Cref{rem:blochkato:open-question-2}.
\item \textbf{Section sec:tam} and \textbf{Section sec:sharp} (the identity (eq:identity) ``says what the excess is''): \Cref{ex:blochkato:sec-tam}, \Cref{rem:blochkato:what-sec-tam-tests},
\Cref{ex:blochkato:selmer-examples}.
\item \textbf{Open question 3}: \Cref{thm:blochkato:tam-additive}, \Cref{rem:blochkato:open-question-3}.
\item \textbf{Section 1} (the Tamagawa number conjecture of Diamond, Flach and Guo; Flach's theorem; the Petersson norm): \Cref{thm:blochkato:dfg-01}, \Cref{thm:blochkato:flach},
\Cref{rem:blochkato:degree-as-l-value}.
\end{itemize}

\section{Overview and conventions}\label{sec:blochkato:overview}

The paper's Theorem thm:main is an identity between $v_\ell(\degphi)$ and three kinds of terms: the length of a Bloch-Kato Selmer
group, Tamagawa factors, and omitted Euler factors. This chapter owns the first two and the bookkeeping that joins them to the
third. It defines the objects of the paper's Section 1.2 ($\ell$-adic representations, lattices, the three duals, the Bloch-Kato local
conditions, the Selmer groups $H^1_f$ and $H^1_\Sigma$, the Tamagawa factor $\Tam_p$); proves Lemma lem:fpr at every $p\ne\ell$, Lemma
lem:tamfactor in every step, and Lemma lem:tamadd; specialises the abstract comparison of Selmer structures of \Cref{ch:galcoh} to the
paper's Lemma lem:dfg21; and proves the Selmer half of Theorem thm:main (\Cref{thm:blochkato:selmer-identity}), which
\Cref{ch:congruence} combines with the congruence ideals. It ends with Remark rem:ns, Section sec:tam, Open questions 2 and 3, and the
Tamagawa number conjecture of Bloch and Kato, of which (eq:identity) is a reading through the modular degree. The group cohomology
is \Cref{ch:cohomology}'s, the Galois cohomology (Tate local duality, the Euler characteristic formulas, unramified classes,
Poitou-Tate) is \Cref{ch:galcoh}'s, and the local representations of $E$ are \Cref{ch:galrep}'s; this chapter cites them and does not
prove them again. Two facts are imported from later chapters and are used only in remarks and examples: the values $e_\ell(E,q)$
(\Cref{ch:lfunctions}) and the congruence-ideal theorems of Diamond, Flach and Guo (\Cref{ch:congruence}).

\begin{convention}[notation of this chapter]\label{conv:blochkato:setup}
\begin{enumerate}[label=(\alph*)]
\item \emph{Objects.} $E/\Q$ is given by a global minimal model and has conductor $N$; $\ell$ is an odd prime unless said otherwise;
$p$ is any prime and $q$ usually a bad prime of $E$. $T=T_\ell E$, $V=T\otimes\Ql$, $\rho=\rho_{E,\ell}$, $\bar\rho$ its reduction. The paper's
objects are $A=\ad T$ (trace zero), $A_f=\ad V=A\otimes\Ql$, $W=A\otimes\Qlmod$, and their Kummer duals $T^D$, $V^D$, $W^D$ (item (b));
$\Sigma$ is the set of primes dividing $N$, and $\delta_p=\dim V^{I_p}$. In the general statements of Sections \ref{sec:blochkato:representations}
to \ref{sec:blochkato:selmer} the letters $V$, $T$, $W=V/T$ denote an arbitrary $\ell$-adic representation, a stable lattice and the
quotient; when such a statement is applied to the paper, the triple is $(A_f,A,W)$, with Kummer dual $(V^D,T^D,W^D)$. This is the
paper's own usage (Section 4.2), and every application below names the triple.
\item \emph{Three duals} (Appendix B, convention B-C4). For a $\Zl$-module $M$: the \emph{linear dual} $M^\vee=\Hom_{\Zl}(M,\Zl)$; the
\emph{Kummer} (or Cartier) \emph{dual} $M^D=\Hom_{\Zl}(M,\Zl(1))=M^\vee(1)$; for a compact or discrete $\ell$-primary module $X$, the
\emph{Pontryagin dual} $X^*=\Hom_{\mathrm{cont}}(X,\Qlmod)$. The symbol ${}^*$ never means a linear dual here. For the paper's objects
$T^D=\Hom(A,\Zl(1))$, not $\Hom(T_\ell E,\Zl(1))$.
\item \emph{Frobenius} (B-C2, B-C3). $\Frob_p$ is an \emph{arithmetic} Frobenius at $p$, and $\Frob_p^{-1}$ the \emph{geometric} one. Group
theory and cohomology use $\Frob_p$: conjugation on $H^1(I_p,-)$, unramified classes $M^{I_p}/(\Frob_p-1)M^{I_p}$, the Tate matrices of Lemma
lem:tamfactor, the tame relation $\Frob\,t\,\Frob^{-1}=t^q$. Local $L$-factors use the geometric Frobenius: for $p\ne\ell$ and a
representation $M$ of $G_{\Q_p}$ on a finite-dimensional $\Ql$-space,
\[
L_p(M,s)=\det\bigl(1-\Frob_p^{-1}p^{-s}\bigm|M^{I_p}\bigr)^{-1}.
\]
We do not write $\Phi_p$ for the geometric Frobenius, because component groups occur in this chapter; $c_p=[E(\Q_p):E_0(\Q_p)]$ is the
Tamagawa number and $\Phig_p$ the geometric component group.
\item \emph{Twists and Hodge-Tate weights} (B-C5). $\Zl(1)=\varprojlim\mu_{\ell^k}$, $M(n)=M\otimes\Zl(1)^{\otimes n}$. The Hodge-Tate weight of $\Ql(1)$
is $+1$; so $V_\ell E$ has weights $\{0,1\}$, $A_f$ has $\{-1,0,1\}$ and $V^D=A_f(1)$ has $\{0,1,2\}$. Every formula below that depends on this
sign says so.
\item \emph{Cohomology} (\Cref{ch:cohomology}). Continuous cochain cohomology; $H^i(K,M)=H^i(G_K,M)$, $H^i(\R,M)=H^i(\Gal(\C/\R),M)$;
$G_{\Q,S}$ is the Galois group of the maximal extension of $\Q$ unramified outside a finite set $S$ of places containing $\ell$ and
$\infty$. ``Length'' is $\Zl$-length; $\Fitt_{\Zl}(X)=\ell^{\len X}\Zl$ for a finite $\Zl$-module $X$, and $v_\ell(\ell^a\Zl)=a$ (B-C10).
\item \emph{Tate data} (B-C14). At a prime $q\ne\ell$ with $v_q(j_E)<0$: $q_E\in\Q_q^\times$ is the Tate parameter (never $u$), $\nu=v_q(q_E)=-v_q(j_E)$
(equal to $n=n_q=v_q(\Delta_E)$ at a multiplicative prime), $u=q_Eq^{-\nu}\in\Z_q^\times$ its unit part, $\bar u\in\F_q^\times$ its reduction, $c(u)$ the
class of $\bar u$ in $\F_q^\times/(\F_q^\times)^{\ell^{v_\ell(q-1)}}$, and $v_\ell(c(u))$ the largest $k\le v_\ell(q-1)$ with $\bar u\in(\F_q^\times)^{\ell^k}$ (so the paper's
``$v_\ell(0)$'' reads $v_\ell(q-1)$). The paper writes $c$ for this class in Lemma lem:tamfactor and $c(u)$ in Section sec:tam; $c_q$ is always the
Tamagawa number. $\varepsilon$ is the unramified character of $G_{\Q_q}$ with $\varepsilon(\Frob_q)=a_q(E)$ at a multiplicative prime.
\item \emph{Characters} (B-C6). $\chi_\ell$ is the $\ell$-adic cyclotomic character and $\bar\chi_\ell$ its reduction; $\chi_d$ is the quadratic character
of $\Q(\sqrt d)$ for a fundamental discriminant $d$ ($\chi_{-3}$, $\chi_{-4}$, and $\chi_{q^*}$ of conductor $q$); $\bar\chi_3=\chi_{-3}$.
\end{enumerate}
\end{convention}

\begin{warning}[three duals]\label{warn:blochkato:three-duals}
$M^\vee=\Hom(M,\Zl)$ is linear (Bloch-Kato and Fontaine-Perrin-Riou write $V^*$ for it); $M^D=\Hom(M,\Zl(1))$ is the paper's ${}^D$; $X^*=\Hom(X,\Qlmod)$
is the paper's ${}^*$. Tate local duality pairs $H^1(\Q_p,W)$ with $H^1(\Q_p,T^D)$, and the only Pontryagin duals in Lemma lem:dfg21 are
duals of finite groups, which have the same length (\Cref{ch:modules}).
\end{warning}

\begin{warning}[two Frobenii]\label{warn:blochkato:two-frobenii}
In cohomology the choice of Frobenius is invisible up to $\ell$-adic units: $(\Frob-1)X=(\Frob^{-1}-1)X$ for any $\Frob$-stable $X$, and
$\det(1-\Frob\mid X)$, $\det(1-\Frob^{-1}\mid X)$ differ by the unit $\det(-\Frob^{-1}\mid X)$ on a lattice $X$. As rational numbers they differ: at a
multiplicative prime $\det(1-\Frob_p^{-1}\mid(V^D)^{I_p})=1-p^{-2}$, while $\det(1-\Frob_p\mid(V^D)^{I_p})=1-p^2$. Only the first is $L_p(A_f,1)^{-1}$.
\end{warning}

The paper's sentences and the results of this chapter that prove them are matched in the following table; the section
\emph{Where the paper uses this chapter} at the head of the chapter lists the same correspondence by result.

\begin{center}\small
\begin{tabular}{p{0.52\textwidth}p{0.4\textwidth}}
\toprule
paper & this chapter\\
\midrule
Section 1.2: the definition of $\Tam_p(M)$, ``at a prime of good reduction it is $0$'' & \Cref{def:blochkato:tamagawa}, \Cref{prop:blochkato:tam-basic}\\
Section 1.2, Section 4.2: the Bloch-Kato Selmer group, $H^1_\Sigma$ & \Cref{def:blochkato:selmer}, \Cref{def:blochkato:sigma-selmer}, \Cref{def:blochkato:adjoint-selmer}\\
Section 4.2: $T^D=\Hom(A,\Zl(1))\simeq A(1)$; $\ad M_{f,\lambda}=\ad T$ & \Cref{prop:blochkato:adjoint-self-dual}\\
Theorem thm:dfg27 & \Cref{thm:blochkato:dfg27}\\
Lemma lem:dfg21 and the paragraph after it & \Cref{lem:blochkato:dfg21}, \Cref{prop:blochkato:local-unramified}, \Cref{lem:blochkato:no-local-invariants}\\
Lemma lem:fpr & \Cref{thm:blochkato:fpr-local}, \Cref{prop:blochkato:tam-dual}, \Cref{lem:blochkato:euler-integral}\\
Lemma lem:h0 (Selmer form) & \Cref{prop:blochkato:h0-selmer}\\
Lemma lem:tamfactor & \Cref{lem:blochkato:tamfactor}, \Cref{prop:blochkato:tam-potmult}\\
Lemma lem:tamadd, Corollary cor:explicit & \Cref{thm:blochkato:tam-additive}, \Cref{cor:blochkato:explicit-identity}\\
proof of Theorem thm:main (the Selmer bookkeeping) & \Cref{cor:blochkato:length-identity}, \Cref{thm:blochkato:selmer-identity}\\
Remark rem:ko(1) & \Cref{cor:blochkato:tamfactor-consequences}\\
Remark rem:ns, Open question 2, the term $h_\ell(E)$ & \Cref{prop:blochkato:rem-ns}, \Cref{rem:blochkato:conj-euler-h}, \Cref{rem:blochkato:open-question-2}\\
Section sec:tam & \Cref{ex:blochkato:sec-tam}, \Cref{rem:blochkato:what-sec-tam-tests}\\
Open question 3 & \Cref{thm:blochkato:tam-additive}, \Cref{rem:blochkato:open-question-3}\\
\bottomrule
\end{tabular}
\end{center}

\section[\texorpdfstring{$\ell$}{l}-adic representations and duals]{\texorpdfstring{$\ell$}{l}-adic representations, lattices, twists and duals}\label{sec:blochkato:representations}

This section sets up the coefficient modules of all the cohomology in the paper: an $\ell$-adic representation $V$, a stable lattice
$T$, the discrete module $W=V/T$, their Tate twists and their three duals, and, for the adjoint, the identification $T^D\simeq A(1)$ that the
paper makes in Section 4.2 ``via the trace form''. It ends with the global invariants of $A_f$, $W$ and $V^D$, which enter Theorem
thm:dfg27 and Lemma lem:h0.

\begin{definition}[$\ell$-adic representations]\label{def:blochkato:l-adic-rep}
Let $K$ be a field with absolute Galois group $G_K$. An \emph{$\ell$-adic representation} of $G_K$ is a finite-dimensional $\Ql$-vector space $V$
with a continuous homomorphism $\rho_V\colon G_K\to\GL(V)$. A \emph{lattice} in $V$ is a finitely generated $\Zl$-submodule $T$ that spans $V$
(so $T\simeq\Zl^{\dim V}$); it is \emph{stable} if $\rho_V(G_K)T=T$. Given a stable lattice we put $W=V/T=T\otimes\Qlmod$, a discrete
$\ell$-primary $G_K$-module.
\end{definition}

\begin{proposition}[lattices]\label{prop:blochkato:lattices}
Let $V$ be an $\ell$-adic representation of $G_K$.
\begin{enumerate}[label=(\roman*)]
\item $V$ has a stable lattice.
\item For a stable lattice $T$ the sequence $0\to T\to V\to W\to0$ is exact, $W\simeq(\Qlmod)^{\dim V}$ as a group, and $W[\ell^k]=\ell^{-k}T/T\simeq T/\ell^kT$ as
$G_K$-modules; the projection $V\to W$ has a continuous set-theoretic section.
\item Any two lattices $T,T'$ are commensurable: $\ell^aT\subseteq T'\subseteq\ell^{-a}T$ for some $a\ge0$. If $T$ is stable and $T/\ell T$ is an
irreducible $\Fl[G_K]$-module, every stable lattice is $\ell^bT$ for some $b\in\Z$.
\end{enumerate}
\end{proposition}

\begin{proof}
(i) Let $T_0$ be any lattice. $\GL(T_0)$ is open in $\GL(V)$, so $H=\rho_V^{-1}(\GL(T_0))$ is an open subgroup of the compact group $G_K$ and has
finite index. If $g_1,\dots,g_r$ are coset representatives of $G_K/H$, then $T=\sum_i\rho_V(g_i)T_0$ is a finite sum of lattices, hence a lattice,
and it is stable because $gg_iH=g_{\pi(i)}H$ for a permutation $\pi$. (ii) Exactness is the definition of $W$; $T\simeq\Zl^d$ gives $W\simeq(\Qlmod)^d$, and
$x\mapsto\ell^{-k}x$ identifies $T/\ell^kT$ with $W[\ell^k]=\ell^{-k}T/T$ equivariantly. A section is given by choosing, in a basis of $T$, the representatives
of $\Qlmod$ of the form $\sum_{i<0}a_i\ell^i$ with $0\le a_i<\ell$; it is continuous because $W$ is discrete (see \Cref{ch:cohomology} for the use of
such sections in the long exact sequence of continuous cohomology). (iii) Both lattices are finitely generated and span $V$, so each is
contained in a multiple of the other. If $T'$ is stable, replace it by $\ell^bT'$ with $b$ minimal such that $\ell^bT'\subseteq T$; then $\ell^bT'\not\subseteq\ell T$,
so its image in $T/\ell T$ is a nonzero stable subspace, hence all of $T/\ell T$, and $\ell^bT'+\ell T=T$ gives $\ell^bT'=T$ by Nakayama's lemma
(\Cref{ch:modules}).
\end{proof}

\begin{definition}[Tate twists]\label{def:blochkato:twists}
$\Zl(1)=\varprojlim_k\mu_{\ell^k}$ (transition maps $\zeta\mapsto\zeta^\ell$), on which $G_K$ acts through the cyclotomic character $\chi_\ell$; $\Ql(1)=\Zl(1)\otimes\Ql$;
$\Zl(n)=\Zl(1)^{\otimes n}$ for $n\ge0$ and $\Zl(-1)=\Hom(\Zl(1),\Zl)$, $\Zl(n)=\Zl(-1)^{\otimes(-n)}$ for $n<0$. For a $\Zl[G_K]$-module $M$,
$M(n)=M\otimes_{\Zl}\Zl(n)$, so $g$ acts on $M(n)$ as $\chi_\ell(g)^n$ times its action on $M$.
\end{definition}

\begin{proposition}[twists]\label{prop:blochkato:twists}
\begin{enumerate}[label=(\roman*)]
\item \sloppy $M(n)(m)=M(n+m)$ and $\Hom(M(n),N)\simeq\Hom(M,N)(-n)$, equivariantly.
\item Let $K=\Q_p$ with $p\ne\ell$. The inertia group $I_p$ acts trivially on $\Zl(1)$ and $\Frob_p$ acts on it by $p$. Hence $M(n)^{I_p}=M^{I_p}(n)$, and the
eigenvalues of $\Frob_p$ (resp. $\Frob_p^{-1}$) on $M(n)^{I_p}\otimes\Ql$ are those on $M^{I_p}\otimes\Ql$ multiplied by $p^n$ (resp. $p^{-n}$).
\item The Weil pairing gives $\wedge^2T_\ell E\simeq\Zl(1)$, hence $\det\rho=\chi_\ell$, $(T_\ell E)^\vee\simeq T_\ell E(-1)$ and $(T_\ell E)^D\simeq T_\ell E$.
\end{enumerate}
\end{proposition}

\begin{proof}
(i) Both are the formulas for $\chi_\ell^n\chi_\ell^m=\chi_\ell^{n+m}$ and for $(g\phi)(x)=g\phi(g^{-1}x)$. (ii) $\Q_p(\mu_{\ell^\infty})/\Q_p$ is unramified for $p\ne\ell$, and the
arithmetic Frobenius raises roots of unity of order prime to $p$ to the $p$-th power, so $\chi_\ell(\Frob_p)=p$ (\Cref{ch:local}). As $I_p$ acts trivially on
$\Zl(n)$, $M(n)^{I_p}=M^{I_p}\otimes\Zl(n)$, and $\Frob_p$ acts as $p^n$ times its action on $M^{I_p}$. (iii) The Weil pairing $T\times T\to\Zl(1)$ is perfect,
alternating and equivariant (\Cref{ch:ellcurves}); it identifies $\wedge^2T$ with $\Zl(1)$, so $\det\rho=\chi_\ell$, and $x\mapsto e(x,\cdot)$ is an isomorphism
$T\to\Hom(T,\Zl(1))=T^D=T^\vee(1)$; twisting by $-1$ gives $T^\vee\simeq T(-1)$.
\end{proof}

A twist rescales every Frobenius eigenvalue at $p$ by the same power of $p$, in the way that multiplying a field by a character of the
scaling group changes its weight without changing its shape.

\begin{definition}[Kummer duals]\label{def:blochkato:duals}
For a lattice $T$ in an $\ell$-adic representation $V$ put $V^D=\Hom_{\Ql}(V,\Ql(1))=V^\vee(1)$, $T^D=\Hom_{\Zl}(T,\Zl(1))=T^\vee(1)$, a stable lattice in $V^D$,
and $W^D=V^D/T^D$. For a compact or discrete $\ell$-primary module $X$, $X^*=\Hom_{\mathrm{cont}}(X,\Qlmod)$ is its Pontryagin dual.
\end{definition}

\begin{proposition}[Kummer duality]\label{prop:blochkato:duals}
Let $T$ be a stable lattice in $V$ and $W=V/T$.
\begin{enumerate}[label=(\roman*)]
\item $W^D\simeq\Hom_{\Zl}(T,\Qlmod(1))=\Hom(T,\mu_{\ell^\infty})$ and $T^D\simeq\Hom_{\mathrm{cont}}(W,\mu_{\ell^\infty})$; the evaluation pairings $T\times W^D\to\mu_{\ell^\infty}$ and
$W\times T^D\to\mu_{\ell^\infty}$ are perfect and $G_K$-equivariant.
\item $(T^D)^D=T$, $(V^D)^D=V$ and $(W^D)^D=W$ (the last in the sense $V^D{}^D/T^D{}^D$).
\item For a finite $\Zl$-module $X$, $X^*$ is finite of the same length as $X$.
\item Bloch and Kato and Fontaine and Perrin-Riou write $V^*(1)$ for $V^D$; Diamond, Flach and Guo write $T^D$, $V^D$, $W^D$ exactly as the paper.
\end{enumerate}
\end{proposition}

\begin{proof}
\sloppy (i) $T$ is free, so $\Hom(T,-)$ is exact: applying it to $0\to\Zl(1)\to\Ql(1)\to\Qlmod(1)\to0$ gives $W^D=\Hom(T,\Ql(1))/\Hom(T,\Zl(1))\simeq\Hom(T,\Qlmod(1))$. Next,
$\Hom_{\mathrm{cont}}(T\otimes\Qlmod,\mu_{\ell^\infty})=\Hom(T,\Hom(\Qlmod,\mu_{\ell^\infty}))=\Hom(T,\Zl(1))$, as $\Hom(\Qlmod,\Qlmod)=\Zl$. Perfectness is
Pontryagin duality between $T$ and $T\otimes\Qlmod$ (\Cref{ch:modules}), twisted by $\Zl(1)$; equivariance holds because all the maps are built from
evaluation. The general statement is in \Cref{ch:galcoh}. (ii) $\Hom(\Hom(T,\Zl(1)),\Zl(1))=T$ for $T$ free of finite rank. (iii) is \Cref{ch:modules}. (iv) is
a matter of notation \cite[Section 3]{BlochKato}, \cite{FPR94}, \cite[Section 2.1]{DFG04}.
\end{proof}

\begin{proposition}[the adjoint is self-dual]\label{prop:blochkato:adjoint-self-dual}
Let $T$ be a free $\Zl$-module of rank $2$ with a continuous action of a group $G$, $A=\ad T$ the trace-zero endomorphisms with $g\cdot x=\rho(g)x\rho(g)^{-1}$.
\begin{enumerate}[label=(\roman*)]
\item If $\ell$ is odd, the trace form $(x,y)\mapsto\tr(xy)$ is a perfect $G$-invariant pairing on $A$ (Gram matrix with determinant $-2$ in the basis $e,h,f$), so
$A\simeq A^\vee$ and
\[
T^D=\Hom(A,\Zl(1))\simeq A(1),\qquad V^D\simeq A_f(1),\qquad W^D\simeq\ad T(1)\otimes\Qlmod .
\]
For $\ell=2$ the image of $A$ in $A^\vee$ has index $2$.
\item $\ad(T\otimes\chi)=\ad T$ for a character $\chi\colon G\to\Zl^\times$, $\ad(\ell^bT)=\ad T$ inside $\ad V$, and $x\mapsto x^\vee$ (the dual map) is a $G$-isomorphism
$\ad T\to\ad(T^\vee)$. In particular $\ad(T(n))=\ad T$ for $n\in\Z$.
\item Consequently: $\ad M_{f,\lambda}\simeq\ad T_\ell E$ for the lattice $M_{f,\lambda}=\varphi_E^*H^1(E,\Zl)\simeq(T_\ell E)^\vee\simeq T_\ell E(-1)$ of Diamond, Flach and Guo; and for a
quadratic twist $E'$ of $E$ (or any twist of $E$ over $\Q_q$ by a character of $G_{\Q_q}$) $\ad T_\ell E'\simeq\ad T_\ell E$ as $G_{\Q_q}$-modules.
\end{enumerate}
\end{proposition}

\begin{proof}
(i) is \Cref{ch:reps} (the trace form and $\Hom(A,\Zl(\chi))\simeq A(\chi)$ for $\ell$ odd); the Gram matrix in the basis $e=\sm0100$, $h=\sm100{-1}$, $f=\sm0010$ is
\[
\begin{pmatrix}0&0&1\\0&2&0\\1&0&0\end{pmatrix},
\]
of determinant $-2$, a unit of $\Zl$ exactly when $\ell$ is odd, and its cokernel has order $2$ for $\ell=2$. Twisting the isomorphism $A\to A^\vee$ by $\Zl(1)$
gives $A(1)\simeq A^\vee(1)=T^D$; tensoring with $\Ql$ and $\Qlmod$ gives the other two. (ii) $\chi(g)$ is a scalar, so conjugation by $\chi(g)\rho(g)$ equals
conjugation by $\rho(g)$; $\End(\ell^bT)=\End(T)$ inside $\End(V)$; and for $x\in\End(T)$ the dual map $x^\vee\in\End(T^\vee)$ satisfies
$(\rho(g)x\rho(g)^{-1})^\vee=\rho^\vee(g)x^\vee\rho^\vee(g)^{-1}$ with $\rho^\vee(g)=(\rho(g)^{-1})^\vee$, and $\tr x^\vee=\tr x$ (\Cref{ch:reps}, the invariance of the adjoint).
(iii) The identification $M_{f,\lambda}=\varphi_E^*H^1(E,\Zl)$ is the paper's Proposition prop:pullback (\Cref{ch:moddeg}), and $H^1(E,\Zl)=(T_\ell E)^\vee\simeq T_\ell E(-1)$ by
\Cref{prop:blochkato:twists}(iii); apply (ii) twice. A twist $\rho\otimes\chi$ has the same adjoint by (ii).
\end{proof}

\begin{paperlink}
Section 4.2 writes ``$T^D=\operatorname{Hom}(A,\Z_\ell(1))\simeq A(1)$ (via the trace form)'': this is (i), and it is one of the places where $\ell$ odd is
used. In the same paragraph the paper says that, by Proposition prop:pullback, $\varphi_E^*$ identifies the lattice $M_{f,\lambda}$ of Diamond,
Flach and Guo with $H^1(E,\Zl)\simeq T_\ell E(-1)$, so that its adjoint is $\ad T$: this is (iii). The original version said instead that $M_{f,\lambda}$ and $T_\ell E$ are homothetic;
they are not (one is the dual of the other), but their adjoints agree, which is all that is used. The step at the primes of type
$\kod{I}_n^*$ in the proof of Theorem thm:main and the statement of Lemma lem:tamfactor at additive potentially multiplicative primes use the
last assertion of (iii).
\end{paperlink}

\begin{definition}[geometric representations]\label{def:blochkato:geometric}
An $\ell$-adic representation of $G_\Q$ is \emph{geometric} (Fontaine and Mazur) if it is unramified outside a finite set of primes and its restriction
to $G_{\Q_\ell}$ is de Rham (\Cref{def:blochkato:period-rings}).
\end{definition}

\begin{proposition}[geometricity and purity]\label{prop:blochkato:purity}
\begin{enumerate}[label=(\roman*)]
\item $V_\ell E$, $\Ql(n)$, and every representation obtained from them by direct sums, tensor products, duals and subquotients ($A_f$, $V^D=A_f(1)$,
$\Sym V$) are geometric.
\item For $p\nmid N\ell$ the eigenvalues of $\Frob_p$ on $V_\ell E$ have complex absolute value $p^{1/2}$, those on $A_f$ absolute value $1$, and those on
$V^D=A_f(1)$ absolute value $p$ (and $p^{-1}$ for $\Frob_p^{-1}$).
\end{enumerate}
\end{proposition}

\begin{proof}
(i) $V_\ell E$ is unramified outside $N\ell$ (\Cref{ch:galrep}) and de Rham at $\ell$ (a theorem of Faltings; for good reduction at $\ell$ it is crystalline,
\Cref{thm:blochkato:crystalline-good-reduction}); $\Ql(n)$ is unramified outside $\ell$ and crystalline; and de Rham representations are stable under the
listed operations (\Cref{prop:blochkato:crystalline-closure} and \cite{Berger}). We do not prove the theorem of Faltings; see \cite{Berger}. (ii) The
eigenvalues of $\Frob_p$ on $V_\ell E$ are the roots $\alpha,\beta$ of $X^2-a_pX+p$, and $|\alpha|=|\beta|=p^{1/2}$ by Hasse's bound (\Cref{ch:ellcurves}). On $A_f$ they
are $\alpha/\beta,1,\beta/\alpha$ (\Cref{ch:galrep}), of absolute value $1$, and \Cref{prop:blochkato:twists}(ii) multiplies them by $p$ on $A_f(1)$.
\end{proof}

\begin{proposition}[global invariants]\label{prop:blochkato:global-invariants}
Let $\ell$ be odd.
\begin{enumerate}[label=(\roman*)]
\item If $\bar\rho$ is absolutely irreducible, $H^0(\Q,A_f)=0$ and $H^0(\Q,W)=0$.
\item For every $E$, $H^0(\Q,V^D)=0$.
\item If $\bar\rho$ is irreducible, $H^0(\Q,W^D)$ is finite and cyclic, of length $h_\ell(E)$; $h_\ell(E)=0$ for $\ell\ge5$ and $h_3(E)\le1$
(\Cref{prop:blochkato:h0-selmer}).
\item The $G_\Q$-modules $A_f$, $A_f(1)=V^D$ and $A_f(-1)$ are pairwise non-isomorphic.
\end{enumerate}
\end{proposition}

\begin{proof}
(i) $W[\ell]=A/\ell A=\ad(T/\ell T)$ (trace-zero matrices modulo $\ell$, as $\ell$ is odd). By Schur's lemma an endomorphism of $T/\ell T$ commuting with the
absolutely irreducible $\bar\rho(G_\Q)$ is a scalar $a$, and its trace $2a$ vanishes only for $a=0$. So $H^0(\Q,W)[\ell]=0$, and $H^0(\Q,W)=0$ because a nonzero
element of an $\ell$-primary group has a multiple of order $\ell$. If $0\ne x\in H^0(\Q,A_f)$, scale it into $A\smallsetminus\ell A$; its image in $A/\ell A$ is a nonzero
invariant, a contradiction. (ii) $V^D$ is unramified at a prime $p\nmid N\ell$, and an invariant vector would be fixed by $\Frob_p$; but no eigenvalue of $\Frob_p$ on
$V^D$ equals $1$, by \Cref{prop:blochkato:purity}(ii). (iii) is the paper's Lemma lem:h0, proved in \Cref{ch:galrep} and restated in
\Cref{prop:blochkato:h0-selmer}. (iv) At a prime $p\nmid N\ell$ the eigenvalues of $\Frob_p$ on the three modules have absolute values $1$, $p$ and $p^{-1}$, so their
characteristic polynomials differ.
\end{proof}

By (iv), only $A_f(1)$ carries the local data (Tamagawa factors at $T^D$, Euler factors at $s=1$) that enter Lemma lem:fpr; this is why
the paper writes $T^D$ and not $A$ in the local terms of Lemma lem:dfg21, and why the identity $\Tam_p(T^D)=\Tam_p(A)$ needs a proof
(\Cref{prop:blochkato:tam-dual}).

\section[Local behaviour at \texorpdfstring{$p\ne\ell$}{p not l}]{Local behaviour at \texorpdfstring{$p\ne\ell$}{p not l}: invariants, Euler factors, weights}\label{sec:blochkato:local}

This section collects what the Selmer bookkeeping needs about the restriction of $A_f$ and $V^D$ to $G_{\Q_p}$ for a prime $p\ne\ell$: the Frobenius
dictionary for $L$-factors, the table of inertia invariants and Euler factors by reduction type, the vanishing of $H^0(\Q_p,V^D)$, and the integrality
of $L_p(A_f,1)^{-1}$. The paper uses these in the paragraph after Lemma lem:dfg21 and in Lemma lem:fpr. The local representations themselves are
\Cref{ch:local}'s and \Cref{ch:galrep}'s; the values $e_\ell(E,p)$ are \Cref{ch:lfunctions}'s.

\begin{proposition}[the Frobenius dictionary]\label{prop:blochkato:frobenius-dictionary}
Let $p\ne\ell$, $F=\Frob_p$, and let $M$ be a finite-dimensional representation of $G_{\Q_p}$ over $\Ql$.
\begin{enumerate}[label=(\alph*)]
\item $L_p(M(1),s)=L_p(M,s+1)$. If $L$-factors were defined with $F$ instead of $F^{-1}$, one would get $L_p(M,s-1)$ on the right.
\item For any $\Zl[F]$-module $X$ on which $F$ is invertible, $(F-1)X=(F^{-1}-1)X$.
\item Let $X$ be a free $\Zl$-module of finite rank $d$ with an automorphism $F$ such that $\det(1-F^{-1}\mid X)\ne0$. Then $F-1$ is injective on $X$ and
\[
\#\,X/(F-1)X=\ell^{v_\ell(\det(1-F^{-1}\mid X))}=\ell^{v_\ell(\det(1-F\mid X))}.
\]
\item On $V_\ell E$ at a good prime $p\ne\ell$, $\det(1-F^{-1}p^{-s}\mid V_\ell E)^{-1}=L_p(E,s+1)$, while $\det(1-F^{-1}p^{-s}\mid(V_\ell E)^\vee)^{-1}=L_p(E,s)$.
\end{enumerate}
\end{proposition}

\begin{proof}
(a) $M(1)^{I_p}=M^{I_p}(1)$ and $F^{-1}$ acts on it as $p^{-1}F^{-1}$ (\Cref{prop:blochkato:twists}(ii)), so $\det(1-F^{-1}p^{-s}\mid M(1)^{I_p})=\det(1-F^{-1}p^{-(s+1)}\mid M^{I_p})$.
With $F$ the factor $p$ would appear instead: $\det(1-Fp^{-s}\mid M(1)^{I_p})=\det(1-Fp^{-(s-1)}\mid M^{I_p})$. (b) $F^{-1}-1=-F^{-1}(F-1)$ and $F^{-1}$ is an
automorphism of $X$. (c) $F-1=F(1-F^{-1})$, so $\det(F-1)=\det F\cdot\det(1-F^{-1})\ne0$, and $\det F\in\Zl^\times$ because $F$ is an automorphism of $X$. The
cokernel of an injective endomorphism $\phi$ of $\Zl^d$ has order $\ell^{v_\ell(\det\phi)}$ (Smith normal form, \Cref{ch:modules}). Similarly $1-F=-F(1-F^{-1})$. (d) The
eigenvalues of $F$ on $V_\ell E$ are $\alpha,\beta$ with $\alpha+\beta=a_p$, $\alpha\beta=p$, so those of $F^{-1}$ are $\alpha^{-1},\beta^{-1}$, with sum $a_p/p$ and product
$1/p$, and $(1-\alpha^{-1}X)(1-\beta^{-1}X)=1-a_pp^{-1}X+p^{-1}X^2$; at $X=p^{-s}$ this is $1-a_pp^{-(s+1)}+p\cdot p^{-2(s+1)}=L_p(E,s+1)^{-1}$. On the dual the
geometric Frobenius has eigenvalues $\alpha,\beta$ and gives $1-a_pp^{-s}+p^{1-2s}$ (\Cref{ch:local}, Appendix B, B-C3).
\end{proof}

\begin{paperlink}
The paper's Section 1.1 now defines ``$L_p(A_f,s)=\det(1-\Frob_p^{-1}p^{-s}\mid A_f^{I_p})^{-1}$ is the true local factor, defined with the geometric
Frobenius as in \cite[\S1.1.1]{DFG04}'', and Lemma lem:fpr writes ``$L_p(A_f,1)^{-1}=L_p(A_f(1),0)^{-1}$'', which is (a) and holds only with the
geometric Frobenius. The paragraph after Lemma lem:dfg21 writes $(V^D)^{I_p}/(\Frob_p-1)$, an arithmetic expression, next to
$\det(1-\Frob_p^{-1}\mid(V^D)^{I_p})=L_p(A_f,1)^{-1}$, a geometric one; (b) and (c) show that the two are compatible. The original version wrote
$\Frob_p$ in the definition of $L_p(A_f,s)$; with the arithmetic Frobenius the multiplicative factor at $s=1$ would be $1-p^2$, not $1-p^{-2}$.
\end{paperlink}

\begin{proposition}[the adjoint at $p\ne\ell$]\label{prop:blochkato:local-adjoint}
Let $p\ne\ell$. The inertia invariants of $A_f$ and of $V^D=A_f(1)$, the eigenvalues of the geometric Frobenius $\Frob_p^{-1}$ on them, and
$L_p(A_f,1)^{-1}=\det(1-\Frob_p^{-1}\mid(V^D)^{I_p})$ are as follows ($\alpha,\beta$ the roots of $X^2-a_pX+p$ at a good prime; $a=a_p(E')$ for a quadratic twist $E'$ of
$E$ with good reduction at $p$ in the row $\kod{I}_0^*$):
\begin{center}\small\setlength{\tabcolsep}{3pt}
\begin{tabular}{llllll}
\toprule
reduction at $p$ & $\delta_p$ & $\dim A_f^{I_p}$ & $\Frob_p^{-1}$ on $A_f^{I_p}$ & $\Frob_p^{-1}$ on $(V^D)^{I_p}$ & $L_p(A_f,1)^{-1}$\\
\midrule
good & $2$ & $3$ & $\frac\alpha\beta,1,\frac\beta\alpha$ & $\frac{\alpha}{\beta p},\frac1p,\frac{\beta}{\alpha p}$ & $(1-\frac{\alpha}{\beta p})(1-\frac1p)(1-\frac{\beta}{\alpha p})$\\
multiplicative & $1$ & $1$ & $p^{-1}$ & $p^{-2}$ & $1-p^{-2}$\\
additive, $v_p(j)<0$ & $0$ & $1$ & $p^{-1}$ & $p^{-2}$ & $1-p^{-2}$\\
$\rho(I_p)=\{\pm1\}$ & $0$ & $3$ & as for $E'$ & as for $E'$ & $\frac{p-1}p\cdot\frac{(p+1)^2-a^2}{p^2}$\\
$\rho(I_p)$ cyclic of order $3,4,6$ & $0$ & $1$ & $c$ & $c\,p^{-1}$ & $1-c\,p^{-1}$\\
$\rho(I_p)$ non-abelian & $0$ & $0$ & none & none & $1$\\
\bottomrule
\end{tabular}
\end{center}
Here $c=\pm1$ is the eigenvalue of Frobenius (either one) on the line $(A_f)^{I_p}$. For $p\ge5$: the row $\{\pm1\}$ is the type $\kod{I}_0^*$; the cyclic row
is the types $\kod{II},\kod{II}^*,\kod{IV},\kod{IV}^*$ with $c=\chi_{-3}(p)$ and the types $\kod{III},\kod{III}^*$ with $c=\chi_{-4}(p)$; the non-abelian row does not
occur. For $p=2$ the types $\kod{IV},\kod{IV}^*$ are in the cyclic row with $c=\chi_{-3}(2)=-1$. At $p=2,3$ the Kodaira type does not determine the row in
general.
\end{proposition}

\begin{proof}
This is the table of \Cref{ch:local} (the adjoint $L$-factors of the local shapes) and of the paper's Lemma lem:euler and Proposition prop:shapes, whose
proofs are in \Cref{ch:galrep} (the images of inertia) and \Cref{ch:lfunctions} (the factors; the value of $L_p(A_f,1)^{-1}$ is the $L$-part of Lemma lem:fpr proved
there); we recall the reasons. At a good prime $A_f$ is unramified
and \Cref{prop:blochkato:purity}(ii) gives the eigenvalues; the product formula is
$(1-\frac\alpha{\beta p})(1-\frac\beta{\alpha p})=1-\frac{\alpha^2+\beta^2}{p^2}+\frac1{p^2}=\frac{(p+1)^2-a_p^2}{p^2}$. At a potentially multiplicative prime, over $\Q_p$, $\rho$ is a
twist by a character (unramified, $\varepsilon$, if the reduction is multiplicative; ramified quadratic if it is additive) of the representation of a Tate curve,
$\sm{\chi_\ell}{\kappa}{0}{1}$, so by \Cref{prop:blochkato:adjoint-self-dual}(ii) $A_f$ is the adjoint of the Tate representation; inertia acts through
$\sm1{*}01$ with $*\ne0$, the invariants are the line of $e=e_1\otimes e_2^\vee$, and $\Frob_p$ acts on $e$ by $\chi_\ell(\Frob_p)=p$ (the upper-left entry divided by the
lower-right one), so $\Frob_p^{-1}$ acts by $p^{-1}$, and by $p^{-2}$ on $e\otimes\Zl(1)$. In the potentially good cases $\rho(I_p)$ is finite (\Cref{ch:galrep}); for
$\rho(I_p)=\{\pm1\}$ the adjoint is unramified and equal to that of the twist $E'$ of good reduction; for a cyclic image generated by $g$ with eigenvalues
$\zeta^{\pm1}$, $\zeta\ne\zeta^{-1}$, the trace-zero commutant of $g$ is the line spanned by $g-g^{-1}$, and Frobenius, which normalises $\langle g\rangle$, conjugates $g$
to $g^{\pm1}$, so it acts on the line by a sign $c$, the same for $\Frob_p$ and $\Frob_p^{-1}$; for $p\ge5$ the tame relation identifies $c$ with $\chi_{-3}(p)$ or
$\chi_{-4}(p)$ (\Cref{ch:galrep}); for a non-abelian image the faithful representation $\rho|_{I_p}$ is absolutely irreducible and the commutant consists of scalars. In each
case the factor is $\prod(1-\lambda p^{-1})$ over the eigenvalues $\lambda$ of $\Frob_p^{-1}$ on $A_f^{I_p}$, that is, $\det(1-\Frob_p^{-1}\mid(V^D)^{I_p})$ by
\Cref{prop:blochkato:frobenius-dictionary}(a).
\end{proof}

\begin{lemma}[local invariants]\label{lem:blochkato:no-local-invariants}
Let $p\ne\ell$.
\begin{enumerate}[label=(\roman*)]
\item $H^0(\Q_p,V^D)=0$, and $\det(1-\Frob_p^{-1}\mid(V^D)^{I_p})=L_p(A_f,1)^{-1}$ is a nonzero rational number.
\item $H^0(\Q_p,A_f)$ is the subspace of $A_f^{I_p}$ fixed by Frobenius. It is $0$ at multiplicative primes, at additive potentially multiplicative primes, in
the cyclic row with $c=-1$ and in the non-abelian row; it is a line at good primes and in the cyclic row with $c=+1$; in the row $\{\pm1\}$ it contains the line
of the eigenvalue $1$.
\end{enumerate}
\end{lemma}

\begin{proof}
(i) By the table every eigenvalue of $\Frob_p^{-1}$ on $(V^D)^{I_p}$ has complex absolute value $p^{-1}$ or $p^{-2}$, so none is $1$; hence $\Frob_p$ fixes no nonzero
vector of $(V^D)^{I_p}$, that is, $H^0(\Q_p,V^D)=((V^D)^{I_p})^{\Frob_p=1}=0$, and the determinant is nonzero. It is rational by the last column of the table
($\alpha\beta=p$, $\alpha+\beta=a_p\in\Z$). (ii) $H^0(\Q_p,A_f)=(A_f^{I_p})^{\Frob_p=1}$; read off the table. At a good prime $\alpha/\beta\ne1$, since $\alpha=\beta$ would give $a_p^2=4p$,
which is impossible for an integer $a_p$; so the eigenvalue $1$ has multiplicity one.
\end{proof}

\begin{warning}\label{warn:blochkato:local-h0-af}
$H^0(\Q_p,A_f)\ne0$ does occur at bad primes: at every prime with $\rho(I_p)=\{\pm1\}$ (type $\kod{I}_0^*$ for $p\ge5$), and at a prime of type $\kod{III}$ with
$p\equiv1\pmod4$. There $H^1(\Q_p,T^D)$ is infinite (\Cref{cor:blochkato:local-finiteness}). Only (i) enters the paper's argument.
\end{warning}

\begin{corollary}[local finiteness]\label{cor:blochkato:local-finiteness}
Let $p\ne\ell$.
\begin{enumerate}[label=(\roman*)]
\item $H^1(\Q_p,T^D)$ is a finitely generated $\Zl$-module of rank $\dim H^0(\Q_p,A_f)$.
\item If $H^0(\Q_p,A_f)=0$ (multiplicative and additive potentially multiplicative primes, the cyclic row with $c=-1$, the non-abelian row), then $H^1(\Q_p,T^D)$
is finite, and $H^1(\Q_p,W)\simeq H^1(\Q_p,T^D)^*$ is finite.
\item In all cases $H^1(\Q_p,T^D)_{\mathrm{tor}}$ is finite.
\end{enumerate}
\end{corollary}

\begin{proof}
$H^1(\Q_p,T^D)$ is finitely generated over $\Zl$ and $H^1(\Q_p,T^D)\otimes\Ql=H^1(\Q_p,V^D)$ (\Cref{ch:galcoh}, the cohomology of $T$, $V$, $W$). For $p\ne\ell$ the local Euler
characteristic of a $\Ql$-representation is $0$, so $\dim H^1(\Q_p,V^D)=\dim H^0(\Q_p,V^D)+\dim H^2(\Q_p,V^D)$, and Tate local duality gives
$H^2(\Q_p,V^D)\simeq H^0(\Q_p,(V^D)^D)^\vee=H^0(\Q_p,A_f)^\vee$ (\Cref{ch:galcoh}). With \Cref{lem:blochkato:no-local-invariants}(i) this is (i). In (ii) a finitely
generated $\Zl$-module of rank $0$ is finite, and Tate local duality for the pair $(W,T^D)$ (\Cref{ch:galcoh}) gives the second statement. (iii) The torsion
of a finitely generated $\Zl$-module is finite.
\end{proof}

\begin{lemma}[integrality of the Euler factor]\label{lem:blochkato:euler-integral}
Let $p\ne\ell$ and $X=T^D$. Then $X^{I_p}=X\cap(V^D)^{I_p}$ is a $\Frob_p$-stable lattice in $(V^D)^{I_p}$,
\[
L_p(A_f,1)^{-1}=\det(1-\Frob_p^{-1}\mid X^{I_p})\in\Z_{(\ell)}\cap\Q^\times,\qquad\#\,X^{I_p}/(\Frob_p-1)X^{I_p}=\ell^{\,v_\ell(L_p(A_f,1)^{-1})}.
\]
The valuation is $v_\ell(p^2-1)$ at multiplicative and at additive potentially multiplicative primes, and $e_\ell(E,p)$ at every additive prime.
\end{lemma}

\begin{proof}
The value of $L_p(A_f,1)^{-1}$, its rationality and nonvanishing, and its valuation by reduction type are the $L$-part of Lemma lem:fpr, proved in
\Cref{ch:lfunctions}; we do not repeat that proof. What this chapter needs in addition is the lattice. $X^{I_p}$ is the intersection of a lattice with a subspace, hence
a lattice in $(V^D)^{I_p}$ (it is saturated in $X$ and spans the subspace, since every vector has a multiple in $X$); it is $\Frob_p$-stable because $I_p$ is normal in
$G_{\Q_p}$. So $\Frob_p^{-1}$ has a characteristic polynomial with coefficients in $\Zl$ on it, and $\det(1-\Frob_p^{-1}\mid X^{I_p})=\det(1-\Frob_p^{-1}\mid(V^D)^{I_p})$ lies
in $\Zl$; being rational and nonzero (\Cref{ch:lfunctions}, or \Cref{lem:blochkato:no-local-invariants}(i)), it lies in $\Zl\cap\Q^\times=\Z_{(\ell)}\cap\Q^\times$. The order of
the cokernel is \Cref{prop:blochkato:frobenius-dictionary}(c).
\end{proof}

\begin{paperlink}
Lemma lem:fpr: ``$L_p(A_f,1)^{-1}=L_p(A_f(1),0)^{-1}$ is a nonzero rational number which is an $\ell$-adic integer, of valuation $e_\ell(E,p)$ at an additive
prime and $v_\ell(1-p^{-2})$ at a multiplicative one''. This $L$-part of the lemma is proved in \Cref{ch:lfunctions}; the first equality is also
\Cref{prop:blochkato:frobenius-dictionary}(a), and \Cref{lem:blochkato:euler-integral} adds the lattice $(T^D)^{I_p}$ and the order of its Frobenius cokernel, which the
Fitting-ideal half (\Cref{thm:blochkato:fpr-local}) uses. The paragraph after Lemma lem:dfg21 says ``for the same reason
$H^0(\Q_p,V^D)=0$'' (\Cref{lem:blochkato:no-local-invariants}(i)) and ``So $H^1_f(\Q_p,T^D)$ is the torsion of $H^1(\Q_p,T^D)$, a finite group''
(\Cref{cor:blochkato:local-finiteness}(iii) and \Cref{prop:blochkato:local-unramified}(v) below). It does not say that $H^1(\Q_p,T^D)$ is finite, which would be
false at the primes of \Cref{warn:blochkato:local-h0-af}.
\end{paperlink}

\section[The representation at \texorpdfstring{$\ell$}{l}]{The representation at \texorpdfstring{$\ell$}{l}: \texorpdfstring{$p$}{p}-adic Hodge theory as a black box}\label{sec:blochkato:hodge}

Nothing in the paper's own proofs uses $p$-adic Hodge theory directly. It enters in two places only: in the definition of the Bloch-Kato local
condition at $\ell$, which is part of the definition of $H^1_f(\Q,W)$ and of $H^1_\Sigma(\Q,W)$ in Section 1.2 and Section 4.2, and in the hypotheses
and proofs of the theorems of Diamond, Flach and Guo that the paper imports. This section states what is needed, with references; every item is
stated, not proved, and the reader may treat $H^1_f(\Q_\ell,-)$ as a black box with the properties of \Cref{prop:blochkato:dimension-formula} and
\Cref{prop:blochkato:local-duality-f}.

\begin{definition}[period rings and admissible representations]\label{def:blochkato:period-rings}
Fontaine constructs topological $\Ql$-algebras $\BdR\supset\Bst\supset\Bcris$ with continuous actions of $G_{\Q_\ell}$, with the following structure:
$\BdR$ is the fraction field of a complete discrete valuation ring with residue field $\C_\ell$ and a uniformiser $t$ on which $G_{\Q_\ell}$ acts by $\chi_\ell$, and it
carries the filtration $\mathrm{Fil}^i\BdR=t^i\BdR^+$ with $\mathrm{gr}^i\BdR\simeq\C_\ell(i)$; $\Bcris$ and $\Bst$ carry a Frobenius $\varphi$, and $\Bst$ a monodromy operator
$N$ with $N\varphi=\ell\varphi N$ and $\Bcris=\Bst^{N=0}$; and $\BdR^{G_{\Q_\ell}}=\Bst^{G_{\Q_\ell}}=\Bcris^{G_{\Q_\ell}}=\Ql$. For an $\ell$-adic representation $V$ of $G_{\Q_\ell}$ and
$\bullet\in\{\mathrm{dR},\mathrm{st},\mathrm{cris}\}$ put $D_\bullet(V)=(B_\bullet\otimes_{\Ql}V)^{G_{\Q_\ell}}$; then $\dim_{\Ql}D_\bullet(V)\le\dim V$, and $V$ is called \emph{de Rham},
\emph{semistable} or \emph{crystalline} when equality holds for $\bullet=\mathrm{dR}$, $\mathrm{st}$, $\mathrm{cris}$. Crystalline implies semistable implies de Rham. $\DdR(V)$ is a
filtered vector space; the \emph{Hodge-Tate weights} of a de Rham $V$ are the integers $-i$ with $\mathrm{gr}^i\DdR(V)\ne0$, counted with multiplicity, so that
$\Ql(1)$, with $\DdR(\Ql(1))=t^{-1}\Ql$, has the weight $+1$ (\Cref{conv:blochkato:setup}(d)). We do not prove these facts; see Fontaine \cite{Fontaine} and
Berger's survey \cite{Berger}.
\end{definition}

\begin{theorem}[good reduction and crystalline representations]\label{thm:blochkato:crystalline-good-reduction}
Let $B$ be an abelian variety over $\Q_\ell$. Then $V_\ell B$ is crystalline if and only if $B$ has good reduction, and semistable if and only if $B$ has
semistable reduction. If $E/\Q$ has good reduction at $\ell$, then $V_\ell E$ is crystalline with Hodge-Tate weights $0,1$, and the crystalline Frobenius on
$\Dcris((V_\ell E)^\vee)$ has characteristic polynomial $X^2-a_\ell X+\ell$.
\end{theorem}

\begin{proof}[Reference]
We do not prove this. The direct implications are theorems of Fontaine and of Faltings, and the converse for crystalline representations is due to
Coleman and Iovita \cite{ColemanIovita99}; see Berger \cite{Berger} for an introduction. The statement for $E$ is the one recorded in \Cref{ch:galrep}
(good reduction at $\ell$) and in \Cref{ch:modgalrep}.
\end{proof}

\begin{proposition}[stability]\label{prop:blochkato:crystalline-closure}
Crystalline (resp. de Rham) representations of $G_{\Q_\ell}$ are stable under direct sums, tensor products, duals, subquotients and Tate twists, and $\Dcris$
(resp. $\DdR$) is compatible with these operations. Consequently, if $\ell\nmid N$, the representations $A_f$, $V^D=A_f(1)$ and $(V_\ell E)^\vee$ are crystalline,
with Hodge-Tate weights $\{-1,0,1\}$, $\{0,1,2\}$ and $\{-1,0\}$.
\end{proposition}

\begin{proof}
The formal properties are stated with references (\cite{Fontaine}, \cite[Section 2]{Berger}); we do not prove them. The consequences: $V_\ell E$ is crystalline
with weights $\{0,1\}$ (\Cref{thm:blochkato:crystalline-good-reduction}), so its dual has weights $\{-1,0\}$ and $V_\ell E\otimes(V_\ell E)^\vee=\End(V_\ell E)$ has weights
$\{-1,0,0,1\}$; $A_f$ is the complement of the scalars $\Ql\cdot1$ (weight $0$), a direct summand because $2$ is invertible ($x=\tfrac12\tr(x)+(x-\tfrac12\tr(x))$),
so it has weights $\{-1,0,1\}$; and twisting by $\Ql(1)$ adds $1$.
\end{proof}

\begin{theorem}[Fontaine and Laffaille]\label{thm:blochkato:fontaine-laffaille}
Let $\ell$ be odd and $a\le b$ integers with $b-a\le\ell-2$. The functor of Fontaine and Laffaille is an exact, fully faithful functor from the category of
strongly divisible filtered Dieudonn\'e modules over $\Zl$ with filtration jumps in $[a,b]$ to the category of $\Zl$-lattices in crystalline representations of $G_{\Q_\ell}$,
whose essential image consists of the stable lattices in crystalline representations with Hodge-Tate weights in $[-b,-a]$ (up to the sign convention); it is
compatible with subquotients and with reduction modulo $\ell^k$. For $V_\ell E$, $\ell\nmid N$, the weights $\{0,1\}$ lie in an interval of length $1\le\ell-2$ for every odd
$\ell$, and the theory describes $T_\ell E$ and the finite flat group schemes $\mathcal E[\ell^k]$ of \Cref{ch:galrep}.
\end{theorem}

\begin{proof}[Reference]
We do not prove this; see Fontaine and Laffaille \cite{FontaineLaffaille}, and \cite[Section 4]{BlochKato} and \cite[Sections 1 and 2]{DFG04} for the use made of
it in the Bloch-Kato and Diamond-Flach-Guo theories.
\end{proof}

\begin{remark}[the case $\ell=3$]\label{rem:blochkato:ell-three}
The adjoint has weights $\{-1,0,1\}$, an interval of length $2$, which exceeds $\ell-2=1$ when $\ell=3$. So \Cref{thm:blochkato:fontaine-laffaille} cannot be
applied to $A$ directly at $\ell=3$. The paper states, after Theorem thm:dfg37, that the theorem of Diamond, Flach and Guo ``covers $\ell=3$ in weight $2$:
$S_f$ excludes only the primes dividing $2N$ and the induced case''. We have not verified in \cite{DFG04} how the integral local condition at $\ell$ is set up for
$\ell=3$; the paper's sentence is reported here as an unverified pointer. The book follows the paper here; the hypotheses of Diamond, Flach and Guo are discussed in \Cref{ch:congruence}.
\end{remark}

\begin{proposition}[local dimensions at $\ell$]\label{prop:blochkato:dimension-formula}
Let $V$ be a de Rham representation of $G_{\Q_\ell}$. With the sign convention of \Cref{conv:blochkato:setup}(d),
\[
\dim_{\Ql}H^1_f(\Q_\ell,V)=\dim H^0(\Q_\ell,V)+\dim\DdR(V)/\mathrm{Fil}^0\DdR(V),
\]
and $\dim\DdR(V)/\mathrm{Fil}^0\DdR(V)$ is the number of positive Hodge-Tate weights of $V$. Examples: $H^1_f(\Q_\ell,\Ql)=H^1_{\ur}(\Q_\ell,\Ql)$ and
$H^1_f(\Q_\ell,\Ql(1))$ (the $\Ql$-span of the Kummer classes of units of $\Z_\ell$) both have dimension $1$. If $\ell\nmid N$, then $H^0(\Q_\ell,V^D)=0$,
\[
\dim H^1_f(\Q_\ell,A_f)=\dim H^0(\Q_\ell,A_f)+1,\qquad\dim H^1_f(\Q_\ell,V^D)=2,
\]
and the sum of the two is $\dim H^1(\Q_\ell,A_f)$, as the duality of \Cref{prop:blochkato:local-duality-f} requires.
\end{proposition}

\begin{proof}[Sketch of proof]
The dimension formula is due to Bloch and Kato \cite[Section 3]{BlochKato}; see also \cite{Berger} and Rubin \cite[Section 1.3]{Rubin}. We do not prove it. For
$\ell\nmid N$: $A_f$ has exactly one positive weight, and $V^D$ has two. An invariant vector of $V^D$ would give a line in $\Dcris(V^D)$ on which $\varphi=1$;
but the eigenvalues of $\varphi$ on $\Dcris(V^D)=\Dcris(A_f)\otimes\Dcris(\Ql(1))$ are $\ell^{-1}$ times those on $\Dcris(A_f)$, which are of the form
$\alpha/\beta,1,\beta/\alpha$ with $\alpha\beta=\ell$ and $|\alpha|=|\beta|=\ell^{1/2}$ (\Cref{thm:blochkato:crystalline-good-reduction}), so none is $1$. Finally Tate's local
Euler characteristic formula at $\ell$ gives $\dim H^1(\Q_\ell,A_f)=\dim H^0(\Q_\ell,A_f)+\dim H^2(\Q_\ell,A_f)+3$, and $H^2(\Q_\ell,A_f)$ is dual to $H^0(\Q_\ell,V^D)=0$
(\Cref{ch:galcoh}); so $\dim H^1(\Q_\ell,A_f)=\dim H^0(\Q_\ell,A_f)+3$.
\end{proof}

For $A_f$ the weights are symmetric about $0$, so the count ``one positive weight'' does not depend on the sign convention; for $V^D$ it does, and
with the opposite convention the formula is stated with negative weights. This is the only place where the sign matters below.

\section{Bloch-Kato local conditions}\label{sec:blochkato:conditions}

This section defines the local conditions $H^1_f(\Q_p,-)$ of Bloch and Kato on $V$, $T$ and $W$, computes them at $p\ne\ell$, and records their duality.
The paper uses them in Section 1.2 (``the Bloch-Kato Selmer group''), in the statement of Lemma lem:dfg21 (``$H^1_f(\Q_p,T^D)=\iota^{-1}H^1_f(\Q_p,V^D)$'') and
in the paragraph after it; the difference between the Bloch-Kato and the unramified conditions on $W$, worked out in \Cref{ex:blochkato:ur-versus-f}, is the
origin of Tamagawa factors.

\begin{definition}[Bloch-Kato local conditions]\label{def:blochkato:local-conditions}
Let $V$ be an $\ell$-adic representation of $G_{\Q_p}$.
\begin{enumerate}[label=(\roman*)]
\item For $p\ne\ell$: $H^1_f(\Q_p,V)=H^1_{\ur}(\Q_p,V)=\ker\bigl(H^1(\Q_p,V)\to H^1(I_p,V)\bigr)$.
\item \sloppy For $p=\ell$: $H^1_e(\Q_\ell,V)=\ker\bigl(H^1(\Q_\ell,V)\to H^1(\Q_\ell,\Bcris^{\varphi=1}\otimes V)\bigr)$, $H^1_f(\Q_\ell,V)=\ker\bigl(H^1(\Q_\ell,V)\to H^1(\Q_\ell,\Bcris\otimes V)\bigr)$ and
$H^1_g(\Q_\ell,V)=\ker\bigl(H^1(\Q_\ell,V)\to H^1(\Q_\ell,\BdR\otimes V)\bigr)$; so $H^1_e\subseteq H^1_f\subseteq H^1_g$, and $H^1_f(\Q_\ell,V)$ classifies the extensions
$0\to V\to X\to\Ql\to0$ with $X$ crystalline when $V$ is crystalline.
\item At the real place: $H^1_f(\R,V)=0$, which is all of $H^1(\R,V)$ (a group killed by $2$ and a $\Ql$-vector space).
\end{enumerate}
Only $H^1_f$ is used in the paper; $H^1_e$ and $H^1_g$ are recorded for orientation \cite[Section 3]{BlochKato}.
\end{definition}

\begin{definition}[propagation]\label{def:blochkato:propagation}
\sloppy Let $T$ be a $G_{\Q_p}$-stable lattice in $V$, $W=V/T$, and $\iota\colon H^1(\Q_p,T)\to H^1(\Q_p,V)$, $\pi\colon H^1(\Q_p,V)\to H^1(\Q_p,W)$ the maps induced by
$T\subset V\to W$. Put
\[
H^1_f(\Q_p,T)=\iota^{-1}\bigl(H^1_f(\Q_p,V)\bigr),\qquad H^1_f(\Q_p,W)=\pi\bigl(H^1_f(\Q_p,V)\bigr),
\]
and likewise $H^1_f(\Q_p,T^D)$, $H^1_f(\Q_p,W^D)$ from $H^1_f(\Q_p,V^D)$.
\end{definition}

\begin{proposition}[propagated conditions]\label{prop:blochkato:propagation}
In the situation of \Cref{def:blochkato:propagation}:
\begin{enumerate}[label=(\alph*)]
\item $\ker\iota=H^1(\Q_p,T)_{\mathrm{tor}}$, it is the image of $H^0(\Q_p,W)$ under the connecting map, and $H^1_f(\Q_p,T)\supseteq H^1(\Q_p,T)_{\mathrm{tor}}$;
\item \sloppy $\iota(H^1_f(\Q_p,T))$ is a lattice in $H^1_f(\Q_p,V)$, so $H^1_f(\Q_p,T)/H^1(\Q_p,T)_{\mathrm{tor}}$ is free of rank $\dim H^1_f(\Q_p,V)$;
\item $H^1_f(\Q_p,W)\simeq H^1_f(\Q_p,V)/\iota H^1_f(\Q_p,T)\simeq(\Qlmod)^{\dim H^1_f(\Q_p,V)}$ is divisible, and it is $0$ exactly when $H^1_f(\Q_p,V)=0$.
\end{enumerate}
\end{proposition}

\begin{proof}
The long exact sequence of $0\to T\to V\to W\to0$ (\Cref{ch:galcoh}) gives $H^0(\Q_p,W)\to H^1(\Q_p,T)\xrightarrow{\iota}H^1(\Q_p,V)\xrightarrow{\pi}H^1(\Q_p,W)$, and $H^1(\Q_p,T)$ is
finitely generated with $H^1(\Q_p,T)\otimes\Ql=H^1(\Q_p,V)$. (a) The kernel of $\iota$ is the image of the torsion group $H^0(\Q_p,W)$, so it consists of torsion
elements; conversely a torsion element maps to a torsion element of a $\Ql$-vector space, which is $0$. The inclusion holds since $0\in H^1_f(\Q_p,V)$. (b) $\iota H^1(\Q_p,T)$ is a
lattice $L$ in $H^1(\Q_p,V)$; for a subspace $U$, $L\cap U$ is finitely generated and spans $U$ (every $u\in U$ has a multiple in $L$), so it is a lattice in $U$, and
$\iota H^1_f(\Q_p,T)=L\cap H^1_f(\Q_p,V)$. (c) The kernel of $\pi$ is $\iota H^1(\Q_p,T)=L$, so $\pi$ restricted to $U=H^1_f(\Q_p,V)$ has kernel $L\cap U$, a lattice, and
$U/(L\cap U)\simeq(\Qlmod)^{\dim U}$.
\end{proof}

\begin{proposition}[the conditions at $p\ne\ell$]\label{prop:blochkato:local-unramified}
Let $p\ne\ell$, $G=G_{\Q_p}$, $I=I_p$, $F=\Frob_p$, and let $T$ be a stable lattice in $V$, $W=V/T$.
\begin{enumerate}[label=(\roman*)]
\item $H^1_f(\Q_p,V)\simeq V^I/(F-1)V^I$, of dimension $\dim H^0(\Q_p,V)$.
\item \sloppy The sequence $0\to H^1_{\ur}(\Q_p,T)\to H^1(\Q_p,T)\xrightarrow{\operatorname{res}}H^1(I,T)^{G/I}\to0$ is exact, with $H^1_{\ur}(\Q_p,T)\simeq T^I/(F-1)T^I$; the same holds
for $W$.
\item $H^1_f(\Q_p,T)=\operatorname{res}^{-1}\bigl(H^1(I,T)_{\mathrm{tor}}\bigr)$, and
\[
H^1_f(\Q_p,T)/H^1_{\ur}(\Q_p,T)\;\simeq\;\bigl(H^1(I,T)_{\mathrm{tor}}\bigr)^{G/I}.
\]
\item $H^1_f(\Q_p,W)$ is the maximal divisible subgroup of $H^1_{\ur}(\Q_p,W)=W^I/(F-1)W^I$.
\item If $H^0(\Q_p,V)=0$, then $H^1_f(\Q_p,V)=0$, $H^1_f(\Q_p,W)=0$ and $H^1_f(\Q_p,T)=H^1(\Q_p,T)_{\mathrm{tor}}$, a finite group. If moreover $H^0(\Q_p,V^D)=0$, then
$H^1(\Q_p,T)$ is finite and $H^1_f(\Q_p,T)=H^1(\Q_p,T)$.
\end{enumerate}
\end{proposition}

\begin{proof}
(i) and (ii) are \Cref{ch:galcoh} (unramified cohomology): inflation identifies $H^1_{\ur}(\Q_p,M)$ with $H^1(G/I,M^I)=M^I/(F-1)M^I$, $G/I\simeq\widehat\Z$ being
procyclic, and the five-term inflation-restriction sequence gives the exactness of (ii) up to the last map. For the discrete module $W$ the sequence is
\Cref{ch:cohomology}'s, and $\operatorname{res}$ is onto $H^1(I,W)^{G/I}$ because the next term $H^2(G/I,W^I)$ vanishes ($\widehat\Z$ has cohomological dimension $1$ on torsion
modules). For the compact module $T$ and for $V$ we use the extension of the five-term sequence to compact and $\Ql$ coefficients when the quotient is isomorphic
to $\widehat\Z$, proved in \Cref{ch:cohomology}, and the surjectivity of $H^1(\Q_p,M)\to H^1(I,M)^{G/I}$ for such $M$ at $p\ne\ell$, proved in \Cref{ch:galcoh}; this is the only
place where it enters, and it is used only at $p\ne\ell$. For $V$, $F-1$ is an endomorphism of the finite-dimensional space $V^I$, so its kernel
$H^0(\Q_p,V)=(V^I)^{F=1}$ and its cokernel have the same dimension.

(iii) Let $\iota_I\colon H^1(I,T)\to H^1(I,V)$. Restriction commutes with $\iota$, so $c\in H^1(\Q_p,T)$ lies in $H^1_f(\Q_p,T)$ if and only if $\iota_I(\operatorname{res}c)=0$. As in
\Cref{prop:blochkato:propagation}(a), the kernel of $\iota_I$ is the image of $W^I$ under the connecting map, a torsion group, and it contains every torsion
element, $H^1(I,V)$ being a $\Ql$-vector space; so $\ker\iota_I=H^1(I,T)_{\mathrm{tor}}$. By (ii) $\operatorname{res}$ maps $H^1(\Q_p,T)$ onto $H^1(I,T)^{G/I}$ with kernel $H^1_{\ur}(\Q_p,T)$,
so $H^1_f(\Q_p,T)/H^1_{\ur}(\Q_p,T)\simeq H^1(I,T)^{G/I}\cap H^1(I,T)_{\mathrm{tor}}=(H^1(I,T)_{\mathrm{tor}})^{G/I}$.

(iv) $\pi$ commutes with inflation, so $H^1_f(\Q_p,W)=\pi(H^1_{\ur}(\Q_p,V))$ is the image of $V^I$ in $W^I/(F-1)W^I$. Let $D\subseteq W^I$ be the image of $V^I$, that is,
$V^I/T^I$ (note $T^I=T\cap V^I$), a divisible group of corank $\dim V^I$. It is the maximal divisible subgroup of $W^I$: $W^I$ is cofinitely generated (its $\ell$-torsion lies
in the finite $W[\ell]$), so $(W^I)_{\mathrm{div}}\simeq(\Qlmod)^r$, and its Tate module is an $I$-invariant submodule of the Tate module $T$ of $W$, hence lies in $T^I$, so
$r\le\operatorname{rank}T^I=\dim V^I$; thus $(W^I)_{\mathrm{div}}=D$ and $W^I/D$ is finite. The image of $D$ in $H^1_{\ur}(\Q_p,W)=W^I/(F-1)W^I$ is divisible with finite cokernel (a quotient
of $W^I/D$), so it is the maximal divisible subgroup: a divisible subgroup maps to a divisible subgroup of a finite group, which is $0$.

(v) By (i) $H^1_f(\Q_p,V)=0$; then $H^1_f(\Q_p,W)=\pi(0)=0$ and $H^1_f(\Q_p,T)=\ker\iota=H^1(\Q_p,T)_{\mathrm{tor}}$, finite because $H^1(\Q_p,T)$ is finitely generated. If also
$H^0(\Q_p,V^D)=0$, then $\dim H^1(\Q_p,V)=\dim H^0(\Q_p,V)+\dim H^0(\Q_p,V^D)=0$ (local Euler characteristic and duality, as in the proof of
\Cref{cor:blochkato:local-finiteness}), so $H^1(\Q_p,T)$ is torsion, hence finite and equal to $H^1_f(\Q_p,T)$.
\end{proof}

\begin{paperlink}
The paragraph after Lemma lem:dfg21: ``For $p\in\Sigma$ the group $H^1_f(\Q_p,V^D)=H^1_{\mathrm{ur}}(\Q_p,V^D)=(V^D)^{I_p}/(\Frob_p-1)$ vanishes, because
$\det(1-\Frob_p^{-1}\mid(V^D)^{I_p})=L_p(A_f,1)^{-1}\ne0$''. This is (i) applied to $V^D$, with \Cref{lem:blochkato:no-local-invariants}(i) and
\Cref{prop:blochkato:frobenius-dictionary}(c) for the nonvanishing. Then (v), for the lattice $T^D$, is ``$H^1_f(\Q_p,T^D)$ is the torsion of $H^1(\Q_p,T^D)$''.
\end{paperlink}

\begin{proposition}[duality of the Bloch-Kato conditions]\label{prop:blochkato:local-duality-f}
Let $V$ be an $\ell$-adic representation of $G_{\Q_v}$, $T\subset V$ a stable lattice, $W=V/T$. Under the cup-product pairing
$H^1(\Q_v,V)\times H^1(\Q_v,V^D)\to H^2(\Q_v,\Ql(1))=\Ql$ the subspaces $H^1_f(\Q_v,V)$ and $H^1_f(\Q_v,V^D)$ are exact annihilators of each other, for every place $v$
(for $v=\ell$ when $V$ is de Rham). Under the perfect pairing $H^1(\Q_v,W)\times H^1(\Q_v,T^D)\to\Qlmod$ the subgroups $H^1_f(\Q_v,W)$ and $H^1_f(\Q_v,T^D)$ are exact
annihilators of each other; hence
\[
H^1(\Q_v,W)/H^1_f(\Q_v,W)\;\simeq\;H^1_f(\Q_v,T^D)^*.
\]
\end{proposition}

\begin{proof}
For $v=p\ne\ell$ both conditions are the unramified ones, and the exact orthogonality of $H^1_{\ur}(\Q_p,V)$ and $H^1_{\ur}(\Q_p,V^D)$ is \Cref{ch:galcoh} (the cup product
of two inflated classes is inflated from $H^2(\widehat\Z,\Ql(1))=0$, and $\dim H^1_{\ur}(V)+\dim H^1_{\ur}(V^D)=\dim H^0(V)+\dim H^0(V^D)=\dim H^1(\Q_p,V)$). For $v=\ell$ this is
\cite[Section 3]{BlochKato}, which we do not prove. For $v=\infty$ all the groups vanish ($\ell$ odd). The integral statement follows from the rational one for
propagated conditions (\Cref{ch:galcoh}, propagation of local conditions): if $y\in H^1_f(\Q_v,T^D)$ and $x\in H^1_f(\Q_v,V)$, then
$\langle\pi x,y\rangle=\langle x,\iota y\rangle\bmod\Zl=0$; conversely, if $y$ annihilates $\pi H^1_f(\Q_v,V)$, then $\langle\ell^{-k}x,\iota y\rangle\in\Zl$ for all $k$, so $\langle x,\iota y\rangle=0$ for
all $x\in H^1_f(\Q_v,V)$ and $\iota y\in H^1_f(\Q_v,V^D)$; the pairing between $H^1(\Q_v,W)$ and $H^1(\Q_v,T^D)$ is perfect (Tate local duality, \Cref{ch:galcoh}), and for a perfect
pairing between a discrete and a compact group the quotient by a subgroup is dual to its annihilator.
\end{proof}

\begin{example}[unramified is not $f$ for $W$]\label{ex:blochkato:ur-versus-f}
Let $E$ have split multiplicative reduction at $p\ne\ell$ with $\ell\mid n=n_p$, and take $V=V_\ell E$, $T=T_\ell E$, $W=E[\ell^\infty]$. In the basis of the Tate shape
(\Cref{ch:galrep}) inertia acts through $t$ with $te_1=e_1$, $te_2=e_2+n'e_1$, $v_\ell(n')=v_\ell(n)$, and $\Frob_p$ by $\sm{p}{c'}{0}{1}$. Then $V^{I_p}=\Ql e_1$, on which $\Frob_p$ acts by
$p\ne1$, so $H^0(\Q_p,V)=0$, $H^1_f(\Q_p,V)=0$ and $H^1_f(\Q_p,W)=0$ (\Cref{prop:blochkato:local-unramified}(v)). On the other hand $x e_1+ye_2\bmod T$ is fixed by $t$ iff
$yn'e_1\in T$, so
\[
W^{I_p}=(\Qlmod)e_1\oplus(\ell^{-v_\ell(n)}\Zl/\Zl)e_2,\qquad(\Frob_p-1)(xe_1+ye_2)=\bigl((p-1)x+c'y\bigr)e_1,
\]
and since $(p-1)\Qlmod=\Qlmod$, $(\Frob_p-1)W^{I_p}=(\Qlmod)e_1$ and $H^1_{\ur}(\Q_p,W)\simeq\Z/\ell^{v_\ell(n)}$. This is the $\ell$-part of the component group $\Z/n$, on which
Frobenius acts trivially in the split case. The Kummer condition at $p\ne\ell$ in the classical Selmer group is $E(\Q_p)\otimes\Qlmod=0$ (\Cref{ex:blochkato:classical-selmer}),
which is the Bloch-Kato condition, not the unramified one; the difference is the Tamagawa phenomenon of \Cref{sec:blochkato:tamagawa}. For the adjoint at a
multiplicative prime $q$, $H^1_f(\Q_q,W)=0$ always (\Cref{lem:blochkato:no-local-invariants}(ii) gives $H^0(\Q_q,A_f)=0$, and \Cref{prop:blochkato:local-unramified}(v) then
gives $H^1_f(\Q_q,W)=0$), while $H^1_{\ur}(\Q_q,W)=0$ when $\ell\nmid n_q$ (\Cref{ex:blochkato:ex-mult}).
\end{example}

\section[Global Selmer groups]{Global Bloch-Kato Selmer groups and the \texorpdfstring{$\Sigma$}{Sigma}-imprimitive Selmer group}\label{sec:blochkato:selmer}

This section defines the Selmer groups of Theorem thm:main: the Bloch-Kato Selmer group $H^1_f(\Q,W)$ (Section 1.2 of the paper) and the
$\Sigma$-imprimitive group $H^1_\Sigma(\Q,W)$ of Diamond, Flach and Guo (Section 4.2), proves their basic finiteness properties, and records the Euler
characteristic count that makes the two vanishing statements of Theorem thm:dfg27 equivalent.

\begin{definition}[Bloch-Kato Selmer groups]\label{def:blochkato:selmer}
Let $V$ be a geometric representation of $G_\Q$, $T$ a stable lattice, $W=V/T$, and let $M$ be one of $V$, $T$, $W$, $V^D$, $T^D$, $W^D$. The \emph{Bloch-Kato Selmer group}
is
\[
H^1_f(\Q,M)=\bigl\{c\in H^1(\Q,M):\ \operatorname{loc}_v(c)\in H^1_f(\Q_v,M)\text{ for every place }v\bigr\},
\]
with the local conditions of \Cref{def:blochkato:local-conditions} and \Cref{def:blochkato:propagation}.
\end{definition}

\begin{lemma}[finite level]\label{lem:blochkato:selmer-S}
Let $S$ be a finite set of places containing $\ell$, $\infty$ and the primes at which $V$ is ramified. For $p\notin S$, $H^1_f(\Q_p,M)=H^1_{\ur}(\Q_p,M)$ for each of the six
modules, and $H^1_f(\Q,M)$ is the subgroup of $H^1(G_{\Q,S},M)$ defined by the conditions at the places of $S$.
\end{lemma}

\begin{proof}
For $V$ this is the definition. If $T$ is unramified at $p$, then $H^1(I_p,T)=\Hom_{\mathrm{cont}}(I_p,T)$ is torsion-free (a homomorphism into the pro-$\ell$ group $T$ factors
through the maximal pro-$\ell$ quotient $\Zl(1)$ of $I_p$, and $\Hom(\Zl(1),T)\simeq T(-1)$), so \Cref{prop:blochkato:local-unramified}(iii) gives $H^1_f=H^1_{\ur}$. If $W$ is
unramified at $p$, $H^1_{\ur}(\Q_p,W)=W/(F-1)W$ is a quotient of a divisible group, hence divisible, and equals $H^1_f(\Q_p,W)$ by
\Cref{prop:blochkato:local-unramified}(iv). Finally, a class of $H^1(\Q,M)$ that is unramified at every $p\notin S$ is inflated from $G_{\Q,S}$: by inflation-restriction
its restriction to the kernel $H$ of $G_\Q\to G_{\Q,S}$ is a homomorphism $H\to M$ that kills every inertia group at a prime outside $S$, and these inertia groups
generate $H$ as a closed normal subgroup (the fixed field of the subgroup they generate is unramified outside $S$).
\end{proof}

\begin{proposition}[relations]\label{prop:blochkato:selmer-relations}
In the situation of \Cref{def:blochkato:selmer}:
\begin{enumerate}[label=(\roman*)]
\item \sloppy $H^1_f(\Q,T)$ is the preimage of $H^1_f(\Q,V)$ under $\iota\colon H^1(G_{\Q,S},T)\to H^1(G_{\Q,S},V)$, and $H^1_f(\Q,T)_{\mathrm{tor}}=H^1(G_{\Q,S},T)_{\mathrm{tor}}=\ker\iota$ is the
image of $H^0(\Q,W)$.
\item $H^1_f(\Q,W)$ is a cofinitely generated $\Zl$-module.
\item The maximal divisible subgroup of $H^1_f(\Q,W)$ is $\pi\bigl(H^1_f(\Q,V)\bigr)$; so the corank of $H^1_f(\Q,W)$ is $\dim H^1_f(\Q,V)$, and $H^1_f(\Q,W)$ is finite if
and only if $H^1_f(\Q,V)=0$.
\end{enumerate}
\end{proposition}

\begin{proof}
We use from \Cref{ch:galcoh} that $H^i(G_{\Q,S},M)$ is finite for finite $M$ and finitely generated over $\Zl$ for $M=T$, and the same for $\Q_p$. (i) The local
condition on $T$ is the preimage of the one on $V$, place by place; and the argument of \Cref{prop:blochkato:propagation}(a) applies to $G_{\Q,S}$; torsion classes
satisfy every local condition because they lie in $\ker\iota$. (ii) $H^1_f(\Q,W)[\ell]$ lies in $H^1(G_{\Q,S},W)[\ell]$, which is a quotient of the finite group
$H^1(G_{\Q,S},W[\ell])$ (from $0\to W[\ell]\to W\xrightarrow{\ell}W\to0$, $W$ being divisible); a discrete $\ell$-primary module with finite $\ell$-torsion is cofinitely generated
(\Cref{ch:modules}). (iii) $D_0=\pi(H^1_f(\Q,V))$ lies in $H^1_f(\Q,W)$ (the local conditions on $W$ are the images of those on $V$) and is divisible of corank
$\dim H^1_f(\Q,V)$ (its kernel $\iota H^1_f(\Q,T)$ is a lattice, as in \Cref{prop:blochkato:propagation}). Conversely let $D$ be the maximal divisible subgroup of
$X=H^1_f(\Q,W)$. The cokernel of $\pi\colon H^1(G_{\Q,S},V)\to H^1(G_{\Q,S},W)$ injects into the finitely generated $\Zl$-module $H^2(G_{\Q,S},T)$, which has no nonzero divisible
subgroup; so $D\subseteq\pi(H^1(G_{\Q,S},V))$. Let $Y=\{y\in H^1(G_{\Q,S},V):\pi(y)\in X\}$ and $\lambda\colon H^1(G_{\Q,S},V)\to Q=\bigoplus_{p\in S}H^1(\Q_p,V)/H^1_f(\Q_p,V)$ the
localisation. For $y\in Y$ and $p\in S$, $\pi_p(\operatorname{loc}_py)\in H^1_f(\Q_p,W)=\pi_p(H^1_f(\Q_p,V))$, so $\operatorname{loc}_py\in H^1_f(\Q_p,V)+\iota_pH^1(\Q_p,T)$. Hence $\lambda(Y)$ lies in
the image $\Lambda$ of $\bigoplus_p\iota_pH^1(\Q_p,T)$, a finitely generated $\Zl$-submodule of $Q$. Choose a complement $C$ of $\ker\lambda=H^1_f(\Q,V)$ in $H^1(G_{\Q,S},V)$;
$\lambda|_C$ is injective, so $M_0=C\cap\lambda^{-1}(\Lambda)$ is finitely generated, and $Y\subseteq H^1_f(\Q,V)+M_0$. Then $X\cap\pi(H^1(G_{\Q,S},V))=\pi(Y)\subseteq D_0+\pi(M_0)$, and $\pi(M_0)$ is
finitely generated and torsion, hence finite. So $D\subseteq D_0+\pi(M_0)$, and the image of $D$ in the finite group $(D_0+\pi(M_0))/D_0$ is divisible, hence $0$: $D\subseteq D_0$.
\end{proof}

\begin{definition}[the $\Sigma$-imprimitive Selmer group]\label{def:blochkato:sigma-selmer}
For a finite set $\Sigma$ of primes not containing $\ell$, and $\ell$ odd,
\[
H^1_\Sigma(\Q,W)=\bigl\{c\in H^1(\Q,W):\ \operatorname{loc}_\ell c\in H^1_f(\Q_\ell,W),\ \operatorname{loc}_pc\in H^1_{\ur}(\Q_p,W)\text{ for }p\notin\Sigma\cup\{\ell\}\bigr\},
\]
with no condition at the primes of $\Sigma$ (none is needed at $\infty$, where $H^1(\R,W)=0$). If $\Sigma$ contains every prime $p\ne\ell$ at which $W$ is ramified, then
$H^1_{\ur}(\Q_p,W)=H^1_f(\Q_p,W)$ for $p\notin\Sigma\cup\{\ell\}$ (\Cref{lem:blochkato:selmer-S}), so $H^1_\Sigma(\Q,W)\supseteq H^1_f(\Q,W)$ is ``Bloch-Kato outside $\Sigma$, nothing on
$\Sigma$''.
\end{definition}

For the paper $\Sigma$ is the set of primes dividing $N$. It may contain primes at which $W$ is unramified, namely the primes with $\rho(I_p)=\{\pm1\}$ (type
$\kod{I}_0^*$ for $p\ge5$); dropping the condition there is harmless.

\begin{proposition}[finiteness of $H^1_\Sigma$]\label{prop:blochkato:sigma-finite}
Let $\Sigma$ be as in \Cref{def:blochkato:sigma-selmer}, containing the primes $p\ne\ell$ at which $W$ is ramified, and suppose $H^1_f(\Q,V)=0$ and $H^0(\Q_p,V^D)=0$ for every
$p\in\Sigma$. Then $H^1_\Sigma(\Q,W)$ is finite: the sequence
\[
0\to H^1_f(\Q,W)\to H^1_\Sigma(\Q,W)\to\bigoplus_{p\in\Sigma}H^1(\Q_p,W)/H^1_f(\Q_p,W)\simeq\bigoplus_{p\in\Sigma}H^1_f(\Q_p,T^D)^*
\]
is exact and its outer terms are finite.
\end{proposition}

\begin{proof}
The kernel of the localisation map consists of the classes satisfying the Bloch-Kato condition at every place (at $p\notin\Sigma\cup\{\ell\}$ by \Cref{lem:blochkato:selmer-S},
at $\infty$ trivially), that is, $H^1_f(\Q,W)$, finite by \Cref{prop:blochkato:selmer-relations}(iii). The isomorphism is \Cref{prop:blochkato:local-duality-f}. For $p\in\Sigma$,
$H^1_f(\Q_p,V^D)=0$ by \Cref{prop:blochkato:local-unramified}(i), so $H^1_f(\Q_p,T^D)=H^1(\Q_p,T^D)_{\mathrm{tor}}$ is finite (\Cref{prop:blochkato:local-unramified}(v)).
\end{proof}

\begin{example}[the classical Selmer group]\label{ex:blochkato:classical-selmer}
Take $V=V_\ell E$, $T=T_\ell E$, $W=E[\ell^\infty]$, $\ell$ odd. At $p\ne\ell$: $E(\Q_p)$ has an open subgroup of finite index isomorphic to $\Z_p$, which is $\ell$-divisible, so
$E(\Q_p)\otimes\Qlmod=0$; and $H^0(\Q_p,V)=0$ because $E(\Q_p)[\ell^\infty]$ is finite, so $H^1_f(\Q_p,W)=0$ by \Cref{prop:blochkato:local-unramified}(v). At $\ell$ the image of the
Kummer map $E(\Q_\ell)\otimes\Qlmod\to H^1(\Q_\ell,W)$ is $H^1_f(\Q_\ell,W)$, by a theorem of Bloch and Kato \cite[Section 3]{BlochKato}, which we do not prove. Hence
$H^1_f(\Q,E[\ell^\infty])=\Sel_{\ell^\infty}(E/\Q)$, the classical Selmer group of \Cref{ch:galcoh}. The Tamagawa numbers $c_p$ do not appear in the Selmer group itself:
they appear in the comparison with relaxed conditions (\Cref{ex:blochkato:ur-versus-f}, \Cref{sec:blochkato:comparison}).
\end{example}

\begin{definition}[the paper's Selmer groups]\label{def:blochkato:adjoint-selmer}
Let $\ell$ be odd. The \emph{adjoint Bloch-Kato Selmer group} is $H^1_f(\Q,W)$ for $W=\ad T\otimes\Qlmod$ (\Cref{def:blochkato:selmer} for the triple $(A_f,A,W)$), and
$H^1_\Sigma(\Q,W)$ is \Cref{def:blochkato:sigma-selmer} with $\Sigma$ the set of primes dividing $N$. Their duals involve $W^D\simeq\ad T(1)\otimes\Qlmod$
(\Cref{prop:blochkato:adjoint-self-dual}), $T^D$ and $V^D$.
\end{definition}

The Selmer group with $\Ql$ coefficients $H^1_f(\Q,A_f)$ vanishes under the hypothesis of Theorem thm:main (\Cref{thm:blochkato:dfg27}), so the length in
(eq:identity) is that of $H^1_f(\Q,W)$ with discrete coefficients; the revised abstract writes $\len_{\Zl}H^1_f(\Q,\ad T_\ell E\otimes\Q_\ell/\Z_\ell)$, while the original
spoke of the Selmer group ``of $\ad V_\ell E$''.

\begin{proposition}[Euler characteristic]\label{prop:blochkato:euler-characteristic}
Let $V$ be a geometric representation of $G_\Q$. Then
\begin{multline*}
\dim H^1_f(\Q,V)-\dim H^1_f(\Q,V^D)\\
=\dim H^0(\Q,V)-\dim H^0(\Q,V^D)+\sum_v\bigl(\dim H^1_f(\Q_v,V)-\dim H^0(\Q_v,V)\bigr),
\end{multline*}
the sum over all places, almost all terms being $0$. For $V=A_f$ with $\ell\nmid N$ and $\bar\rho$ absolutely irreducible the terms are:
\begin{center}\small
\begin{tabular}{lll}
\toprule
term & value & reason\\
\midrule
$\dim H^0(\Q,A_f)-\dim H^0(\Q,V^D)$ & $0-0$ & \Cref{prop:blochkato:global-invariants} ($\bar\rho$ abs.\ irreducible)\\
$v=p\ne\ell$ & $0$ & \Cref{prop:blochkato:local-unramified}(i)\\
$v=\ell$ & $+1$ & \Cref{prop:blochkato:dimension-formula}\\
$v=\infty$ & $-\dim A_f^{c=1}=-1$ & $c$ has eigenvalues $1,-1,-1$ on $A_f$\\
\bottomrule
\end{tabular}
\end{center}
So, under these hypotheses, $\dim H^1_f(\Q,A_f)=\dim H^1_f(\Q,V^D)$; in particular $H^1_f(\Q,A_f)=0$ if and only if $H^1_f(\Q,V^D)=0$.
\end{proposition}

\begin{proof}[Sketch of proof]
The general formula is the $\Ql$ form of the formula of Greenberg and Wiles, which \Cref{ch:galcoh} states for finite modules (with references to
\cite{Wiles95} and \cite{DDT}; it is not proved in this book): apply it to $T/\ell^kT$ with the
conditions induced from $H^1_f(\Q_v,V)$ (preimages and images, which are mutually dual by \Cref{prop:blochkato:local-duality-f}) and compare the growth in $k$ of the
two sides, using \Cref{prop:blochkato:propagation} and \Cref{prop:blochkato:selmer-relations}; see \cite{FPR94} and \cite[Section 1.7]{Rubin}. For the table: complex
conjugation $c$ acts on $V_\ell E$ with eigenvalues $1,-1$ (\Cref{ch:galrep}), so in a basis where $\rho(c)=\operatorname{diag}(1,-1)$ it fixes $h$ and negates $e$ and $f$; and at
$\infty$, $H^1_f(\R,A_f)=0$ and $H^0(\R,A_f)=A_f^{c=1}$. The irreducibility of $\bar\rho$ enters only through $H^0(\Q,A_f)=0$; without it one may use instead
Faltings's theorem $\End_{\Ql[G_\Q]}V_\ell E=\Ql$ for $E/\Q$ (which gives $H^0(\Q,A_f)=0$ directly), which we do not prove.
\end{proof}

\begin{paperlink}
Theorem thm:dfg27 states ``$H^0(\Q,V^D)=H^1_f(\Q,V^D)=0$ (and the same for $A\otimes\Q_\ell$)''. The parenthesis is the other half of the equivalence just proved,
and Remark rem:ns uses the same equivalence: ``The vanishing of $H^1_f(\Q,V^D)$, equivalently of $H^1_f(\Q,\ad V_3E)$, is not known for these curves''.
\end{paperlink}

\section{Tamagawa factors}\label{sec:blochkato:tamagawa}

This section defines the Tamagawa factor of Fontaine and Perrin-Riou at a prime $p\ne\ell$, proves its index description and its basic
properties, and proves the paper's Lemma lem:fpr: the local term $H^1_f(\Q_p,T^D)$ of Lemma lem:dfg21 has Fitting ideal $L_p(A_f,1)^{-1}\Tam_p(T^D)$, and
$\Tam_p(T^D)=\Tam_p(A)$. It also checks that for $T_\ell E$ the Tamagawa factor is the $\ell$-part of the Tamagawa number $c_p$, as the name promises. The
paper uses the definition in Section 1.2 and the lemma in the proof of Theorem thm:main and in Remark rem:ns. A Tamagawa factor is the size of a local
defect counted by torsion in a first cohomology group, much as the index of a sublattice measures a ratio of volumes.

\begin{definition}[Tamagawa factor]\label{def:blochkato:tamagawa}
Let $p\ne\ell$ and let $M$ be a free $\Zl$-module of finite rank with a continuous action of $G_{\Q_p}$. The group $G_{\Q_p}/I_p$ acts on $H^1(I_p,M)$ by
$(\sigma\cdot c)(i)=\sigma\,c(\sigma^{-1}i\sigma)$ (\Cref{ch:cohomology}), and preserves its torsion subgroup. The \emph{Tamagawa factor} of $M$ at $p$ is
\[
\Tam_p(M)=\Fitt_{\Zl}\Bigl(\bigl(H^1(I_p,M)_{\mathrm{tor}}\bigr)^{G_{\Q_p}}\Bigr),
\]
the Fitting ideal of the $G_{\Q_p}$-invariants of the torsion of $H^1(I_p,M)$. It is a nonzero ideal $\ell^a\Zl$, with $a=\len\bigl((H^1(I_p,M)_{\mathrm{tor}})^{G_{\Q_p}}\bigr)$, and we write
$v_\ell(\Tam_p(M))=a$.
\end{definition}

That the ideal is nonzero: $H^1(I_p,M)$ is a finitely generated $\Zl$-module (\Cref{ch:cohomology}: $I_p$ acts on the pro-$\ell$ module $M$ through a quotient that is an
extension of $\Zl(1)$ by a finite group, the wild inertia and the prime-to-$\ell$ part of tame inertia contributing nothing), so its torsion is finite, and so are its
invariants; the Fitting ideal of a finite $\Zl$-module $X$ is $\ell^{\len X}\Zl$ (\Cref{ch:modules}).

\begin{paperlink}
\sloppy Section 1.2 of the revised paper defines ``$\Tam_p(M)=\Fitt_{\Z_\ell}\bigl((H^1(I_p,M)_{\mathrm{tor}})^{G_{\Q_p}}\bigr)$, an ideal of $\Z_\ell$ \cite[p.~708]{DFG04},
\cite[I.4.2.2]{FPR94}''. The original version wrote $\Fitt_{\Z_\ell}\bigl(H^1(I_p,M)_{\mathrm{tor}}\bigr)^{G_{\Q_p}}$, with the invariants outside the Fitting ideal, which has no
meaning (an ideal of $\Zl$ has no Galois action); the computation in Lemma lem:tamfactor has always used the invariants inside.
\end{paperlink}

\begin{proposition}[the index description]\label{prop:blochkato:tam-index}
For $p\ne\ell$ and $M$ as in \Cref{def:blochkato:tamagawa},
\[
\ell^{v_\ell(\Tam_p(M))}=\#\bigl(H^1(I_p,M)_{\mathrm{tor}}\bigr)^{G_{\Q_p}}=\bigl[H^1_f(\Q_p,M):H^1_{\ur}(\Q_p,M)\bigr].
\]
\end{proposition}

\begin{proof}
\Cref{prop:blochkato:local-unramified}(iii), with $T=M$.
\end{proof}

So the Tamagawa factor measures how much larger the Bloch-Kato condition on the lattice is than the naive unramified condition. For $W$ the inclusion
goes the other way (\Cref{prop:blochkato:local-unramified}(iv)); both reflect the same finite group.

\begin{proposition}[basic properties]\label{prop:blochkato:tam-basic}
Let $p\ne\ell$.
\begin{enumerate}[label=(\roman*)]
\item If $I_p$ acts trivially on $M$, then $\Tam_p(M)=\Zl$.
\item If $I_p$ acts on $M$ through a finite quotient of order prime to $\ell$, then $\Tam_p(M)=\Zl$.
\item $\Tam_p(M)$ depends only on the isomorphism class of the $G_{\Q_p}$-module $M$; in particular $\Tam_p(\ell^bM)=\Tam_p(M)$ for a lattice $M$ in a representation.
\item $\Tam_p(\ad T)=\Tam_p(\ad(T\otimes\chi))$ for every character $\chi$ of $G_{\Q_p}$ with values in $\Zl^\times$.
\end{enumerate}
In particular $\Tam_p(\ad T_\ell E)=\Zl$ at every prime $p\nmid N\ell$, and at every prime with $\rho(I_p)=\{\pm1\}$.
\end{proposition}

\begin{proof}
(i) $H^1(I_p,M)=\Hom_{\mathrm{cont}}(I_p,M)$, and a continuous homomorphism into the pro-$\ell$ group $M$ factors through the maximal pro-$\ell$ quotient of $I_p$, which is $\Zl(1)$
for $p\ne\ell$ (\Cref{ch:local}); so $H^1(I_p,M)\simeq\Hom(\Zl(1),M)=M(-1)$ is torsion-free. (ii) is the case $\#\Delta$ prime to $\ell$ of \Cref{lem:blochkato:finite-inertia-image}
below: the torsion is $H^1(\Delta,M)$, killed by $\#\Delta$, which is a unit. (iii) The definition uses only the module. (iv) The two modules are equal
(\Cref{prop:blochkato:adjoint-self-dual}(ii)). For the last sentence: at a good prime $\ad T$ is unramified, and if $\rho(I_p)=\{\pm1\}$ then $I_p$ acts trivially on $\ad T$.
\end{proof}

\begin{theorem}[Fontaine and Perrin-Riou at $p\ne\ell$]\label{thm:blochkato:fpr-local}
Let $p\ne\ell$ and let $X$ be a free $\Zl$-module of finite rank with a continuous action of $G_{\Q_p}$ such that $H^0(\Q_p,X\otimes\Ql)=0$. Then $H^1_f(\Q_p,X)=H^1(\Q_p,X)_{\mathrm{tor}}$ is
finite, there is an exact sequence
\[
0\to X^{I_p}/(\Frob_p-1)X^{I_p}\to H^1(\Q_p,X)_{\mathrm{tor}}\to\bigl(H^1(I_p,X)_{\mathrm{tor}}\bigr)^{G_{\Q_p}}\to0,
\]
and
\begin{gather*}
\Fitt_{\Zl}H^1_f(\Q_p,X)=\det\bigl(1-\Frob_p^{-1}\bigm|X^{I_p}\bigr)\cdot\Tam_p(X),\\
\len H^1_f(\Q_p,X)=v_\ell\det\bigl(1-\Frob_p^{-1}\bigm|X^{I_p}\bigr)+v_\ell(\Tam_p(X)).
\end{gather*}
Moreover $H^1(\Q_p,X)_{\mathrm{tor}}\simeq H^0(\Q_p,X\otimes\Qlmod)$. For the paper's $X=T^D$ and every $p\ne\ell$:
\begin{gather*}
\Fitt_{\Zl}H^1_f(\Q_p,T^D)=L_p(A_f,1)^{-1}\cdot\Tam_p(T^D),\\
\len H^1_f(\Q_p,T^D)=\len H^0(\Q_p,W^D)=v_\ell\bigl(L_p(A_f,1)^{-1}\bigr)+v_\ell\bigl(\Tam_p(T^D)\bigr).
\end{gather*}
\end{theorem}

\begin{proof}
Write $G=G_{\Q_p}$, $I=I_p$, $F=\Frob_p$, $U=X\otimes\Ql$.

\emph{Step 1 (the unramified part is finite).} $H^0(\Q_p,U)=(U^I)^{F=1}=0$, so $1$ is not an eigenvalue of $F$, nor of $F^{-1}$, on $U^I$, and $\det(1-F^{-1}\mid X^I)\ne0$. By
\Cref{prop:blochkato:frobenius-dictionary}(c) the group $X^I/(F-1)X^I$ is finite of order $\ell^{v_\ell\det(1-F^{-1}\mid X^I)}$; its Fitting ideal is generated by
$\det(1-F^{-1}\mid X^I)$.

\emph{Step 2 (inflation-restriction).} By \Cref{prop:blochkato:local-unramified}(ii) the sequence $0\to X^I/(F-1)X^I\to H^1(\Q_p,X)\to H^1(I,X)^{G/I}\to0$ is exact.

\emph{Step 3 (torsion).} Taking torsion subgroups keeps it exact. Exactness at the left and in the middle is clear (the left term is finite, hence torsion). For
surjectivity, let $y\in(H^1(I,X)^{G/I})_{\mathrm{tor}}$, say $\ell^ky=0$, and lift it to $x\in H^1(\Q_p,X)$. Then $\ell^kx$ lies in the finite left term, so $\ell^{k+m}x=0$ for some $m$,
and $x$ is torsion. Finally $(H^1(I,X)^{G/I})_{\mathrm{tor}}=(H^1(I,X)_{\mathrm{tor}})^{G/I}$.

\emph{Step 4 (Fitting ideals).} Lengths are additive in short exact sequences of finite $\Zl$-modules, and $\Fitt(Y)=\ell^{\len Y}\Zl$ for finite $Y$ (\Cref{ch:modules}). So
$\Fitt H^1(\Q_p,X)_{\mathrm{tor}}=\det(1-F^{-1}\mid X^I)\cdot\Tam_p(X)$.

\emph{Step 5 (the Bloch-Kato condition).} By \Cref{prop:blochkato:local-unramified}(v), $H^1_f(\Q_p,X)=H^1(\Q_p,X)_{\mathrm{tor}}$. The isomorphism with $H^0(\Q_p,X\otimes\Qlmod)$ is the
connecting map of $0\to X\to U\to X\otimes\Qlmod\to0$, injective because $H^0(\Q_p,U)=0$, with image $\ker\iota$ (\Cref{prop:blochkato:propagation}(a)).

For $X=T^D$: $H^0(\Q_p,V^D)=0$ and $\det(1-F^{-1}\mid(T^D)^I)=L_p(A_f,1)^{-1}$ by \Cref{lem:blochkato:no-local-invariants} and \Cref{lem:blochkato:euler-integral}.
\end{proof}

\begin{proposition}[$\Tam_p(T^D)=\Tam_p(A)$]\label{prop:blochkato:tam-dual}
\sloppy Let $\ell$ be odd, $p\ne\ell$, $A=\ad T_\ell E$ and $T^D=\Hom(A,\Zl(1))\simeq A(1)$. Then $\Tam_p(T^D)=\Tam_p(A)$. More precisely, $H^1(I_p,A(1))=H^1(I_p,A)(1)$ as $\Zl$-modules
with $G_{\Q_p}/I_p$-action, so the two torsion subgroups are the same group with Frobenius actions differing by the factor $p$; and:
\begin{enumerate}[label=(\alph*)]
\item if the reduction at $p$ is good or potentially good, both torsion subgroups are $0$;
\item if the reduction at $p$ is potentially multiplicative, with the notation of \Cref{lem:blochkato:tamfactor} (and \Cref{prop:blochkato:tam-potmult}), both invariant
modules are isomorphic to $\Z/\ell^{a}\oplus\Z/\ell^{m}$ with $a=v_\ell(\nu)$ and $m=\min(a,v_\ell(p-1),v_\ell(c'))=\min(a,v_\ell(p-1),v_\ell(c(u)))$.
\end{enumerate}
\end{proposition}

\begin{proof}
$I_p$ acts trivially on $\Zl(1)$, so $A(1)=A\otimes\Zl(1)$ is the same $I_p$-module as $A$ tensored with a free rank-one module with trivial $I_p$-action; cochains and
cohomology commute with $\otimes\Zl(1)$, and $\Frob_p$ acts on $\Zl(1)$ by $p$. (a) At good and potentially good primes $H^1(I_p,A)_{\mathrm{tor}}=0$ by \Cref{prop:blochkato:tam-basic}
and \Cref{thm:blochkato:tam-additive}, so both factors are trivial. (b) By \Cref{lem:blochkato:tamfactor} (applied, at an additive potentially multiplicative prime, to the
Tate shape of \Cref{prop:blochkato:tam-potmult}), $H^1(I_p,A)_{\mathrm{tor}}=(\Z/\ell^a)\bar e\oplus(\Z/\ell^a)\bar h$ with $\Frob_p$ acting by $p^{-1}\Ad\rho(\Frob_p)$, so on the twist it acts by
$\Ad\rho(\Frob_p)$: $\bar e\mapsto p\bar e$, $\bar h\mapsto\bar h-2\varepsilon c'\bar e$. Thus $x\bar e+y\bar h$ is fixed iff $(p-1)x-2\varepsilon c'y\equiv0\pmod{\ell^a}$, the kernel of the linear form
$\phi(x,y)=(p-1)x-2\varepsilon c'y$ on $(\Z/\ell^a)^2$. The image of $\phi$ is the ideal generated by $p-1$ and $c'$ in $\Z/\ell^a$ ($2\varepsilon$ is a unit), of order $\ell^{a-m}$, so the kernel
has order $\ell^{2a-(a-m)}=\ell^{a+m}$. Its structure: if $m=v_\ell(p-1)\le v_\ell(c')$ (and $m<a$), write $\phi=\ell^m\phi'$ with $\phi'(x,y)=\gamma x-\delta y$, $\gamma$ a unit; the kernel is
$\{(x,y):\gamma x\equiv\delta y\pmod{\ell^{a-m}}\}$, and $y$ is free in $\Z/\ell^a$ while $x$ is determined modulo $\ell^{a-m}$, giving $\Z/\ell^a\oplus\Z/\ell^m$; the case $m=v_\ell(c')<v_\ell(p-1)$ is
symmetric, and if $m=a$ the form is $0$. The invariants for $A$ itself are $\{x\bar e+y\bar h:(p-1)y\equiv0,\ c'y\equiv0\}=\Z/\ell^a\oplus\Z/\ell^m$ (\Cref{lem:blochkato:tamfactor}).
\end{proof}

The general statement $\Tam_p(T^D)=\Tam_p(A)$ is \cite[(57)]{DFG04}, which the paper cites in Lemma lem:fpr; for the adjoint of an elliptic curve the proposition proves
it in every case that occurs. As a check, an exhaustive computation over $\ell\in\{3,5\}$, $v\in\{0,1,2\}$, all primes $p<120$ with $p\ne\ell$, $\varepsilon=\pm1$ and
$c'\in\{0,\dots,\ell^3-1\}$ ($26{,}448$ parameter choices) found both invariant modules of order $\ell^{v+\min(v,v_\ell(p-1),v_\ell(c'))}$ in every case.

\begin{proposition}[the Tamagawa factor of $T_\ell E$]\label{prop:blochkato:tam-of-E}
For $p\ne\ell$, $v_\ell(\Tam_p(T_\ell E))=v_\ell(c_p)$, the valuation of the Tamagawa number $c_p=[E(\Q_p):E_0(\Q_p)]$.
\end{proposition}

\begin{proof}
Let $T=T_\ell E$, $V$, $W=E[\ell^\infty]$, $I=I_p$. As in the proof of \Cref{prop:blochkato:local-unramified}(iii), the connecting map identifies $W^I/(W^I)_{\mathrm{div}}$ with
$H^1(I,T)_{\mathrm{tor}}=\ker(H^1(I,T)\to H^1(I,V))$, equivariantly for $G_{\Q_p}/I$ (the image of $V^I$ in $W^I$ is $(W^I)_{\mathrm{div}}$, as shown there). Now $W^I=E(\Q_p^{\ur})[\ell^\infty]$, and
by \Cref{ch:reduction} the group $E(\Q_p^{\ur})$ is an extension of $\Phig_p=\Phi_p(\Fbar_p)$ by $E_0(\Q_p^{\ur})$, which is an extension of the group of nonsingular points of the
reduction over $\Fbar_p$ (an elliptic curve, $\mathbb G_m$ or $\mathbb G_a$) by a pro-$p$ group; so $E_0(\Q_p^{\ur})$ is $\ell$-divisible and its $\ell$-power torsion is divisible
($(\Qlmod)^2$, $\mu_{\ell^\infty}$ or $0$). Hence $0\to E_0(\Q_p^{\ur})[\ell^\infty]\to W^I\to\Phig_p[\ell^\infty]\to0$ is exact (\Cref{ch:galrep}, inertia invariants and component
groups), the first term is $(W^I)_{\mathrm{div}}$ (a divisible subgroup with finite quotient), and $H^1(I,T)_{\mathrm{tor}}\simeq\Phig_p[\ell^\infty]$ as Frobenius modules. The
Frobenius invariants of $\Phig_p$ are $\Phi_p(\F_p)$, of order $c_p$.
\end{proof}

At a multiplicative prime this is also the last step of the proof of \Cref{lem:blochkato:tamfactor}: $H^1(I,T)_{\mathrm{tor}}=(\Z/\ell^{v_\ell(n)})\bar e_1$ with Frobenius acting by $\varepsilon$, so
$v_\ell(\Tam_p(T))=v_\ell(n)$ if the reduction is split and $0$ if it is nonsplit. On all curves of Cremona's tables of conductor below $3{,}000$ ($17{,}834$ curves), at every
multiplicative prime $q$ and $\ell\in\{3,5,7\}$, $\ell\ne q$ ($101{,}336$ pairs), $v_\ell(c_q)$ is $v_\ell(n_q)$ if $a_q=1$ and $0$ if $a_q=-1$, as it should be. At a prime $p\ne3$ of type $\kod{IV}$ or
$\kod{IV}^*$ with $\ell=3$, $H^1(I,T)_{\mathrm{tor}}=H^1(C_3,T)\simeq\F_3$ (\Cref{prop:blochkato:ad-free-c3}), and the proposition says that Frobenius acts on it trivially exactly when
$3\mid c_p$ (for example \curve{50a1} at $p=5$, type $\kod{IV}$, $c_5=3$).

\section[Lemma lem:tamfactor]{The Tamagawa factor of \texorpdfstring{$\ad T$}{ad T} at a multiplicative prime: Lemma lem:tamfactor}\label{sec:blochkato:tamfactor}

This section proves the paper's Lemma lem:tamfactor in full, with every step the paper compresses expanded, and extends it to the additive potentially
multiplicative primes, $q=2$ included. The lemma is the source of the terms $t_\ell(E,q)$ in (eq:thm) and of the refinement tested in Section sec:tam. The
inputs are: the Tate shape of $\rho|_{G_{\Q_q}}$ (\Cref{ch:galrep}), the tame relation (\Cref{ch:local}), the cohomology of procyclic groups (\Cref{ch:cohomology}),
and the identification of the Frobenius value of the Kummer cocycle of $q_E$ with the class of $\bar u$ (\Cref{ch:galcoh}).

\begin{convention}[Tate data]\label{conv:blochkato:tate-data}
Let $q\ne\ell$ be a prime of multiplicative reduction, $n=n_q$, $G=G_{\Q_q}$, $I=I_q$, and $\Frob$ an arithmetic Frobenius in $G$. We use:
\begin{enumerate}[label=(\alph*)]
\item (\Cref{ch:galrep}, the Tate shape) a $\Zl$-basis $e_1,e_2$ of $T$ in which $\rho(g)=\sm{\chi_\ell(g)\varepsilon(g)}{\varepsilon(g)\kappa(g)}{0}{\varepsilon(g)}$ for $g\in G$, where $\kappa$ is a Kummer
cocycle of $q_E$ with values in $\Zl(1)\simeq\Zl$ (the corner is $\varepsilon\kappa$, not $\kappa$: only $\varepsilon\kappa$ satisfies the cocycle relation for these matrices when $\varepsilon\ne1$);
\item (\Cref{ch:local}) $\kappa|_I=n\,t_\ell$ with $t_\ell$ the tame character, so a topological generator $t$ of the pro-$\ell$ tame quotient of $I$ acts by $te_1=e_1$,
$te_2=e_2+n'e_1$ with $v_\ell(n')=v_\ell(n)$, and the kernel of $I\to\langle t\rangle\simeq\Zl$ has no nontrivial pro-$\ell$ quotient;
\item (\Cref{ch:local}) the tame relation $\Frob\,t\,\Frob^{-1}=t^q$, equivalently $\Frob^{-1}t\,\Frob=t^{q^{-1}}$ with $q^{-1}\in\Zl^\times$;
\item (\Cref{ch:galcoh}) $\kappa=n\kappa_q+\kappa_u$ for $q_E=q^nu$, with $\kappa_u$ unramified, and the canonical splitting $H^1(\Q_q,\Zl(1))=\Zl\kappa_q\oplus H^1_{\ur}(\Q_q,\Zl(1))$ given by the
uniformiser $q$;
\item (\Cref{ch:galcoh}) $H^1_{\ur}(\Q_q,\Zl(1))=H^1(G/I,\Zl(1))=\Zl(1)/(q-1)\simeq\F_q^\times\otimes\Zl$, under which the class of $\kappa_u$ is the class of $\bar u$.
\end{enumerate}
We write $\rho(\Frob)=\sm{\varepsilon q}{c'}{0}{\varepsilon}$ with $\varepsilon=\varepsilon(\Frob)=a_q(E)$ and $c'=\varepsilon\kappa(\Frob)$.
\end{convention}

\begin{lemma}[the Tamagawa factor at a multiplicative prime]\label{lem:blochkato:tamfactor}
Let $q$ be a prime of multiplicative reduction, $n=n_q=v_q(\Delta_E)$, $\ell$ an odd prime different from $q$, $M=\ad T$, and $q_E$, $u$, $c(u)$ as in
\Cref{conv:blochkato:setup}(f). Then $H^1(I_q,M)_{\mathrm{tor}}\simeq(\Z/\ell^{v_\ell(n)})^2$ and
\[
v_\ell\bigl(\Tam_q(\ad T)\bigr)=v_\ell(n)+\min\bigl(v_\ell(n),\,v_\ell(q-1),\,v_\ell(c(u))\bigr).
\]
In particular $v_\ell(\Tam_q(\ad T))\ge t_\ell(E,q)=v_\ell(n)$, with equality unless $q\equiv1\pmod\ell$, $\ell\mid n$ and $\bar u$ is an $\ell$-th power in $\F_q^\times$.
\end{lemma}

\begin{proof}
\emph{Step 1 (the shape on inertia).} By \Cref{conv:blochkato:tate-data}(a), (b), $I$ acts on $T$ through $\langle t\rangle$, with $t=\sm1{n'}01$ in the basis $e_1,e_2$ ($\varepsilon$ is trivial on
$I$).

\emph{Step 2 (reduction to a procyclic group).} Let $K$ be the kernel of $I\to\langle t\rangle$. It acts trivially on $M$, and $H^1(K,M)=\Hom_{\mathrm{cont}}(K,M)=0$ because $K$ has no
nontrivial pro-$\ell$ quotient and $M$ is pro-$\ell$. By inflation-restriction (\Cref{ch:cohomology}), $H^1(I,M')=H^1(\langle t\rangle,M')$ for every finitely generated
$\Zl$-module $M'$ with this action, and for the procyclic group $\langle t\rangle\simeq\Zl$, $c\mapsto c(t)$ identifies $H^1(\langle t\rangle,M')$ with $M'/(t-1)M'$ (\Cref{ch:cohomology}).

\emph{Step 3 (the computation).} In the basis $e=e_1\otimes e_2^\vee=\sm0100$, $h=e_1\otimes e_1^\vee-e_2\otimes e_2^\vee=\sm100{-1}$, $f=e_2\otimes e_1^\vee=\sm0010$ of $M$, conjugation by $t$
gives
\[
tet^{-1}=e,\qquad tht^{-1}=h-2n'e,\qquad tft^{-1}=f+n'h-n'^2e,
\]
so $(t-1)e=0$, $(t-1)h=-2n'e$, $(t-1)f=n'h-n'^2e$. As $2$ is a unit, $(t-1)M$ contains $n'e$ and then $n'h$, and it is contained in $n'(\Zl e\oplus\Zl h)$; so
$(t-1)M=n'(\Zl e\oplus\Zl h)$ and
\[
H^1(I,M)=(\Zl/n')\bar e\oplus(\Zl/n')\bar h\oplus\Zl\bar f,\qquad H^1(I,M)_{\mathrm{tor}}\simeq(\Z/\ell^{v_\ell(n)})^2 .
\]

\emph{Step 4 (the Frobenius action).} For a cocycle $c$ on $\langle t\rangle$ and $a\in\Zl$, $c(t^a)\equiv a\,c(t)$ modulo $(t-1)M$ (for $a\in\N$, $c(t^a)=(1+t+\dots+t^{a-1})c(t)$, and
the general case follows by continuity). Hence, by \Cref{conv:blochkato:tate-data}(c),
\[
(\Frob\cdot c)(t)=\rho(\Frob)\,c(\Frob^{-1}t\,\Frob)\,\rho(\Frob)^{-1}=\Ad\rho(\Frob)\,c(t^{q^{-1}})\equiv q^{-1}\Ad\rho(\Frob)\,c(t)
\]
modulo $(t-1)M$:
the \emph{arithmetic} Frobenius acts on $M/(t-1)M$ as $q^{-1}\Ad\rho(\Frob)$. With $\rho(\Frob)=\sm{\varepsilon q}{c'}{0}{\varepsilon}$ one computes $\rho(\Frob)e\rho(\Frob)^{-1}=qe$ and
$\rho(\Frob)h\rho(\Frob)^{-1}=h-2\varepsilon c'e$ (the inverse is $\sm{\varepsilon q^{-1}}{-c'q^{-1}}{0}{\varepsilon}$, as $\varepsilon^2=1$). So on the torsion
\[
\bar e\mapsto\bar e,\qquad\bar h\mapsto q^{-1}(\bar h-2\varepsilon c'\bar e),
\]
and $x\bar e+y\bar h$ is fixed iff $q^{-1}y\equiv y$ and $2\varepsilon c'q^{-1}y\equiv0$ modulo $\ell^{v_\ell(n)}$, that is, iff $(q-1)y\equiv0$ and $c'y\equiv0\pmod{\ell^{v_\ell(n)}}$. The solutions
are $\ell^{v_\ell(n)}$ values of $x$ and the $y\in(\Z/\ell^{v_\ell(n)})[\ell^{\min(v_\ell(q-1),v_\ell(c'))}]$, a cyclic group of order $\ell^{\min(v_\ell(n),v_\ell(q-1),v_\ell(c'))}$ (\Cref{ch:groups}):
\[
\bigl(H^1(I,M)_{\mathrm{tor}}\bigr)^{G}\simeq\Z/\ell^{v_\ell(n)}\oplus\Z/\ell^{\min(v_\ell(n),\,v_\ell(q-1),\,v_\ell(c'))}.
\]

\emph{Step 5 (the class $c'$).} The number $c'=\varepsilon\kappa(\Frob)$ is a cocycle value, not a class. Replacing $\Frob$ by $\Frob\,t^a$ replaces $\rho(\Frob)$ by
$\rho(\Frob)\sm1{an'}01$, whose corner is $c'+\varepsilon q\,a n'$; replacing $\kappa$ by a cohomologous cocycle $\kappa+(g-1)z$ changes $\kappa(\Frob)$ by $(q-1)z$. So $c'$ is well
defined only modulo the ideal $(n,q-1)\Zl=\ell^{\min(v_\ell(n),v_\ell(q-1))}\Zl$, and $\min(v_\ell(n),v_\ell(q-1),v_\ell(c'))$ is unchanged by either change (ultrametric
inequality). By \Cref{ch:galcoh} (the identification of $c'$): for the lift $\Frob_0$ fixing a compatible system of $\ell$-power roots of $q$, $\kappa_q(\Frob_0)=0$, so
$c'=\varepsilon\kappa_u(\Frob_0)$, and the class of $\kappa_u(\Frob_0)$ in $\Zl(1)/(q-1)\simeq\F_q^\times\otimes\Zl\simeq\Z/\ell^{v_\ell(q-1)}$ is the class of $\bar u$. An element of the cyclic group
$\F_q^\times/(\F_q^\times)^{\ell^{v_\ell(q-1)}}$ lies in the subgroup of $\ell^k$-th powers ($k\le v_\ell(q-1)$) iff, under any isomorphism with $\Z/\ell^{v_\ell(q-1)}$, it is divisible by
$\ell^k$. Hence, for $\Frob_0$, $\min(v_\ell(q-1),v_\ell(c'))=v_\ell(c(u))$ in the sense of \Cref{conv:blochkato:setup}(f), and therefore
$\min(v_\ell(n),v_\ell(q-1),v_\ell(c'))=\min(v_\ell(n),v_\ell(q-1),v_\ell(c(u)))$ for every choice of lift and cocycle.

\emph{Step 6 (conclusion).} By \Cref{def:blochkato:tamagawa} and Steps 4, 5, $v_\ell(\Tam_q(\ad T))$ is the length $v_\ell(n)+\min(v_\ell(n),v_\ell(q-1),v_\ell(c(u)))$. The excess
$\min(\dots)$ is positive iff $\ell\mid n$, $\ell\mid q-1$ and $\bar u\in(\F_q^\times)^\ell$. As a check, the same steps for $M=T$ give $(t-1)T=n'\Zl e_1$, $H^1(I,T)_{\mathrm{tor}}=(\Z/\ell^{v_\ell(n)})\bar e_1$, and
$\Frob$ acts by $q^{-1}\cdot\varepsilon q=\varepsilon$; the invariants have length $v_\ell(n)$ if $\varepsilon=1$ and $0$ if $\varepsilon=-1$ ($\ell$ odd), the $\ell$-part of $c_q$ (\Cref{prop:blochkato:tam-of-E}).
\end{proof}

\begin{paperlink}
The proof of Lemma lem:tamfactor in the paper is Steps 1 to 6 in compressed form. Three points are corrected or made precise above, and the revised
version of the paper has adopted all three. The corner of the matrices is $\varepsilon\kappa$ (Step 1; the original wrote $\kappa$). The Frobenius in the matrices is
arithmetic (Step 4; the paper now says so). The original identified $c'$ with ``the image of $\kappa$ in $H^1(G/I,\Z_\ell(1)^I)$ \dots, well defined modulo $q-1$''; the image
exists only through the splitting attached to the uniformiser $q$, $c'$ equals $\varepsilon$ times that image only for the lift $\Frob_0$, and the indeterminacy is
$(n,q-1)$ (Step 5). The revised proof reads: ``it is $\varepsilon$ times the class of $u$ for the Frobenius lift fixing a compatible system of $\ell$-power roots of $q$, and in
general it is well defined only modulo $\ell^{\min(v_\ell(n),v_\ell(q-1))}$''. None of this changes the formula. The identification of $\Tam_q$ with the Fitting ideal of the
invariants of the torsion is \Cref{def:blochkato:tamagawa}.
\end{paperlink}

\begin{corollary}[consequences]\label{cor:blochkato:tamfactor-consequences}
In the situation of \Cref{lem:blochkato:tamfactor}:
\begin{enumerate}[label=(\roman*)]
\item $v_\ell(\Tam_q(\ad T))-t_\ell(E,q)=\min(v_\ell(n),v_\ell(q-1),v_\ell(c(u)))\in\{0,1,\dots,v_\ell(n)\}$;
\item the excess is positive iff $\ell\mid n$, $q\equiv1\pmod\ell$ and $\bar u\in(\F_q^\times)^\ell$; it is $0$ for $q=2$;
\item (Remark rem:ko(1)) at such a prime $v_\ell(\Tam_q)\ge t_\ell(E,q)+1$, and the excess exceeds $1$ only if $\ell^2\mid n$, $\ell^2\mid q-1$ and $\bar u\in(\F_q^\times)^{\ell^2}$;
\item the excess does not depend on whether the reduction is split.
\end{enumerate}
\end{corollary}

\begin{proof}
(i) to (iii) are read off the formula; $v_\ell(2-1)=0$. (iv) The formula involves only $n$, $q$ and $\bar u$; $\varepsilon$ entered Step 4 only through the unit $2\varepsilon$.
\end{proof}

\begin{paperlink}
Remark rem:ko(1) says that at a potentially multiplicative prime $q\equiv1\pmod\ell$ with $\ell\mid n$ whose Tate parameter has an $\ell$-th power as unit part,
Lemma lem:tamfactor gives $v_\ell(\Tam_q)\ge t_\ell(E,q)+1$, and that (eq:identity) raises the bound accordingly. This is (iii), and \Cref{prop:blochkato:tam-potmult} covers the potentially
multiplicative primes. The revised abstract now says ``by at least one'' and includes the condition $\ell\mid n_q$; the original abstract said ``by one'' and omitted
$\ell\mid n_q$, without which the excess is $0$ whatever the unit part.
\end{paperlink}

\begin{proposition}[potentially multiplicative additive primes]\label{prop:blochkato:tam-potmult}
Let $q\ne\ell$ ($q=2$ allowed) be a prime of additive, potentially multiplicative reduction, $\nu=-v_q(j_E)$, $q_E$ the Tate parameter (determined by $j(q_E)=j_E$), and
$u=q_Eq^{-\nu}$. Then
\[
v_\ell\bigl(\Tam_q(\ad T)\bigr)=v_\ell(\nu)+\min\bigl(v_\ell(\nu),v_\ell(q-1),v_\ell(c(u))\bigr)\ \ge\ t_\ell(E,q).
\]
\end{proposition}

\begin{proof}
Over $\Q_q$, $E$ is the twist of the Tate curve $E_{q_E}$ by a ramified quadratic character $\chi$ of $G_{\Q_q}$ (\Cref{ch:reduction}, \Cref{ch:galrep}), so $T_\ell E\simeq T_\ell E_{q_E}\otimes\chi$
and $\ad T_\ell E\simeq\ad T_\ell E_{q_E}$ as $G_{\Q_q}$-modules (\Cref{prop:blochkato:adjoint-self-dual}(ii)). The Tate curve $E_{q_E}$ has split multiplicative reduction with
$v_q(\Delta)=v_q(q_E)=\nu$ and Tate parameter $q_E$, and the proof of \Cref{lem:blochkato:tamfactor} uses only the $G_{\Q_q}$-module $\ad T$ in its Tate shape (with $\varepsilon=1$); so the formula
holds with $n=\nu$. At $q=2$, $v_\ell(q-1)=0$ and the value is $t_\ell(E,2)=v_\ell(-v_2(j_E))$.
\end{proof}

For odd $q$ one may instead twist globally, by $\chi_{q^*}$, as the original proof of Theorem thm:main did with the curve $E_{\{p\}}$ of Lemma lem:twist; for $q=2$ that lemma
does not apply, and the local twist is needed. The revised paper states the formula at additive potentially multiplicative primes in Lemma lem:tamfactor itself,
with the local argument. At $q=2$ the index of a Kodaira symbol $\kod{I}_m^*$ need not be $-v_2(j)$, and the reduction may be potentially good (\Cref{ch:reduction}); the
proposition uses $\nu=-v_2(j_E)$.

\begin{example}[Tate data and Tamagawa factors]\label{ex:blochkato:tamfactor-examples}
The table lists, for some curves and primes $q$ of potentially multiplicative reduction, the Kodaira type, $\nu$, the reduction $\bar u$ of the unit part of the Tate
parameter, $\ell$, $v=v_\ell(q-1)$, $k=v_\ell(c(u))$, the resulting $v_\ell(\Tam_q(\ad T))$, and $t_\ell=v_\ell(\nu)$. The unit part was computed from $j_E$: since $q_E=j_E^{-1}+744j_E^{-2}+\cdots$
with integer coefficients and $v_q(j_E^{-1})=\nu\ge1$, $\bar u$ is the reduction of $j_E^{-1}q^{-\nu}$; for \curve{11a1} and \curve{791a1} this agrees with PARI's
\texttt{ellinit(E,O(q\^{}20)).tate}.
\begin{center}\small\setlength{\tabcolsep}{4pt}
\begin{tabular}{llrrrrrrr}
\toprule
curve, $q$ & type & $\nu$ & $\bar u$ & $\ell$ & $v$ & $k$ & $v_\ell(\Tam_q(\ad T))$ & $t_\ell$\\
\midrule
\curve{11a1}, $11$ & $\kod{I}_5$ & $5$ & $10$ & $5$ & $1$ & $1$ & $2$ & $1$\\
\curve{14a1}, $7$ & $\kod{I}_3$ & $3$ & $6$ & $3$ & $1$ & $1$ & $2$ & $1$\\
\curve{539c1}, $7$ & $\kod{I}_3^*$ & $3$ & $5$ & $3$ & $1$ & $0$ & $1$ & $1$\\
\curve{791a1}, $7$ & $\kod{I}_3$ & $3$ & $6$ & $3$ & $1$ & $1$ & $2$ & $1$\\
\curve{791a1}, $113$ & $\kod{I}_2$ & $2$ & $106$ & $3$ & $0$ & $0$ & $0$ & $0$\\
\curve{1184f1}, $37$ & $\kod{I}_3$ & $3$ & $36$ & $3$ & $2$ & $2$ & $2$ & $1$\\
\curve{1748e1}, $19$ & $\kod{I}_6$ & $6$ & $12$ & $3$ & $2$ & $1$ & $2$ & $1$\\
\bottomrule
\end{tabular}
\end{center}
For \curve{11a1} at $\ell=5$ and \curve{14a1} at $\ell=3$ the residual representation is reducible: the lemma applies (it is local and needs no irreducibility), but Theorem
thm:main does not; \curve{14a1} has $\degphi=1$. \curve{539c1} has the same $j$-invariant, hence the same Tate parameter at $7$, as its twist \curve{77c1}; $5$ is not a cube
modulo $7$. In \curve{1184f1} the cap $v_\ell(\nu)=1$ binds although $\bar u=-1$ is a ninth power modulo $37$; in \curve{1748e1}, $12$ is a cube but not a ninth power modulo
$19$. The curves \curve{791a1}, \curve{1184f1}, \curve{1748e1} have $3\nmid N$, $E[3]$ irreducible, mod $3$ image not \texttt{3Ns}, and no CM; they reappear in
\Cref{ex:blochkato:selmer-examples}.
\end{example}

\section[Additive primes and Open question 3]{Tamagawa factors at additive primes, and Open question 3}\label{sec:blochkato:additive}

The original version of the paper left the Tamagawa factors of $\ad T$ at the additive primes uncomputed (``$v_\ell(\Tam_p(\ad T))\ge0$ at an additive prime'') and
asked for them in Open question 3. This section computes them: they are trivial at every potentially good additive prime, for every odd $\ell$, and given by
\Cref{prop:blochkato:tam-potmult} at the potentially multiplicative ones. The revised paper contains the first statement as Lemma lem:tamadd and the resulting
explicit identity as Corollary cor:explicit.

\begin{lemma}[finite inertia image]\label{lem:blochkato:finite-inertia-image}
Let $q\ne\ell$ and let $M$ be a free $\Zl$-module of finite rank on which $I=I_q$ acts continuously through a finite quotient $\Delta=I/K$. Then inflation identifies
$H^1(\Delta,M)$ with $H^1(I,M)_{\mathrm{tor}}$, and $H^1(I,M)/H^1(\Delta,M)$ embeds in $\Hom(K,M)^\Delta\subseteq M(-1)$, which is torsion-free.
\end{lemma}

\begin{proof}
Inflation-restriction (\Cref{ch:cohomology}) gives $0\to H^1(\Delta,M)\to H^1(I,M)\to\Hom_{\mathrm{cont}}(K,M)^\Delta$. The group $K$ is the inertia group of a finite extension of $\Q_q^{\ur}$, and
a continuous homomorphism from $K$ to the pro-$\ell$ group $M$ factors through the maximal pro-$\ell$ abelian quotient of $K$, which is $\Zl(1)$ as $q\ne\ell$ (\Cref{ch:local}); so
$\Hom(K,M)\simeq M(-1)$ is torsion-free. $H^1(\Delta,M)$ is killed by $\#\Delta$, hence torsion, and it is therefore the whole torsion subgroup.
\end{proof}

\begin{proposition}[the adjoint lattice is free over $\Z_3\lbrack C_3\rbrack$]\label{prop:blochkato:ad-free-c3}
Let $\sigma\in\GL_2(\Z_3)$ have order $3$, $T=\Z_3^2$ with $\sigma$ acting, and $P=\langle\sigma\rangle$. Then $\ad T$, with $P$ acting by conjugation, is a free $\Z_3[P]$-module of rank one; hence
$H^1(P,\ad T)=0$. By contrast $H^1(P,T)\simeq\F_3$.
\end{proposition}

\begin{proof}
The eigenvalues of $\sigma$ are cube roots of unity; as $\Q_3$ contains no primitive cube root of unity, the characteristic polynomial, which has coefficients in $\Z_3$, is
$(X-1)^2$ or $X^2+X+1$, and $(X-1)^2$ would make the semisimple element $\sigma$ (it has finite order) equal to $1$. So $\Z_3[\sigma]\simeq\Z_3[X]/(X^2+X+1)\simeq\Z_3[\zeta_3]$, a
discrete valuation ring with uniformiser $1-\zeta_3$ (\Cref{ch:ant}), and $T$, torsion-free of $\Z_3$-rank $2$, is free of rank one over it. Choosing a generator $e_1$ and the
basis $e_1,\sigma e_1$ gives $\sigma=\sm0{-1}1{-1}$. In the basis $e,h,f$ of $\ad T$,
\[
e=\sm0100,\qquad\sigma e\sigma^{-1}=\sm00{-1}0=-f,\qquad\sigma^2e\sigma^{-2}=\sm{-1}1{-1}1=e-h-f,
\]
and the matrix of coordinates $\begin{psmallmatrix}1&0&1\\0&0&-1\\0&-1&-1\end{psmallmatrix}$ has determinant $-1$. So $e,\sigma e,\sigma^2e$ is a $\Z_3$-basis and $\ad T=\Z_3[P]\,e\simeq\Z_3[P]$. For a cyclic
group, $H^1(P,\Z_3[P])=\ker N/(\sigma-1)\Z_3[P]=0$, where $N=1+\sigma+\sigma^2$. On $T$, $N=0$ (as $\sigma^2+\sigma+1=0$), so $H^1(P,T)=T/(\sigma-1)T\simeq\Z_3[\zeta_3]/(1-\zeta_3)\simeq\F_3$.
\end{proof}

\begin{sloppypar}Conceptually: modulo $3$, $\sigma$ is a nontrivial unipotent on $T/3T$, and $\ad(T/3T)\simeq\Sym(T/3T)\otimes\det^{-1}$ is the symmetric square of a $2\times2$ Jordan block, a single
$3\times3$ Jordan block in characteristic $3$, that is, the regular module $\F_3[C_3]$; Nakayama's lemma lifts this to $\Z_3$. The coordinates above were also checked in PARI.\end{sloppypar}

\begin{theorem}[Tamagawa factors at additive primes]\label{thm:blochkato:tam-additive}
Let $\ell$ be odd, $q\ne\ell$ a prime of additive reduction of $E$, and $A=\ad T_\ell E$.
\begin{enumerate}[label=(\roman*)]
\item If the reduction at $q$ is potentially good, then $H^1(I_q,A)_{\mathrm{tor}}=0$, so $\Tam_q(A)=\Tam_q(A(1))=\Zl$.
\item If it is potentially multiplicative, $\Tam_q(A)$ is given by \Cref{prop:blochkato:tam-potmult}.
\end{enumerate}
\end{theorem}

\begin{proof}
(ii) is \Cref{prop:blochkato:tam-potmult}. (i) The image $\rho(I_q)$ is finite and its elements have characteristic polynomials in $\Z[X]$ (Serre and Tate; \Cref{ch:galrep}), so
their orders lie in $\{1,2,3,4,6\}$. Let $\Delta$ be the image of $I_q$ in $\Aut(A)$, that is, $\rho(I_q)/(\rho(I_q)\cap\{\pm1\})$. By \Cref{lem:blochkato:finite-inertia-image},
$H^1(I_q,A)_{\mathrm{tor}}=H^1(\Delta,A)$, which is killed by $\#\Delta$, a product of powers of $2$ and $3$. If $\ell\ge5$ it is $0$. Let $\ell=3$. Restriction followed by corestriction is
multiplication by the index, so the $3$-primary group $H^1(\Delta,A)$ injects into $H^1(P,A)$ for a Sylow $3$-subgroup $P$ of $\Delta$ (\Cref{ch:cohomology}). A Sylow $3$-subgroup $P'$ of
$\rho(I_q)$ maps isomorphically onto $P$ (the kernel $\{\pm1\}$ has order prime to $3$). $P'$ has exponent $3$; if it had order at least $9$ it would contain a subgroup $C_3\times C_3$
(a group of order $9$ is abelian), whose elements, being commuting and of finite order, are simultaneously diagonalisable over $\C$ inside $\SL_2$ ($\det\rho=\chi_\ell$ is trivial
on $I_q$), hence of the form $\operatorname{diag}(a,a^{-1})$ with $a^3=1$, a group of order $3$; a contradiction. So $P$ is trivial or generated by the image of an element $\sigma$ of order
$3$, and $H^1(P,A)=0$ by \Cref{prop:blochkato:ad-free-c3}. Finally $H^1(I_q,A(1))=H^1(I_q,A)(1)$ (\Cref{prop:blochkato:tam-dual}).
\end{proof}

So $\Tam_q(\ad T)$ is trivial at every potentially good additive prime, of every Kodaira type, at every $q\ne\ell$ including $q=2,3$, even for $\ell=3$ at the types $\kod{II}$,
$\kod{II}^*$, $\kod{IV}$, $\kod{IV}^*$, where $\Tam_q(T)$ can be divisible by $3$ (\Cref{prop:blochkato:tam-of-E}).

\begin{paperlink}
Lemma lem:tamadd of the revised paper is (i), with the same proof: inflation-restriction through the finite image of inertia, element orders $1,2,3,4,6$, and the
freeness of $\ad T$ over $\Z_3[P]$ shown by the same three coordinate vectors. The original proof of Theorem thm:main said only ``$v_\ell(\Tam_p(\ad T))\ge0$'' at the
additive primes.
\end{paperlink}

\begin{corollary}[the explicit identity]\label{cor:blochkato:explicit-identity}
Let $\ell$ be odd, $\ell\nmid N$. Then
\[
\sum_{q\mid N}v_\ell\bigl(\Tam_q(\ad T)\bigr)=\sum_{q:\,v_q(j_E)<0}\Bigl(t_\ell(E,q)+\min\bigl(t_\ell(E,q),v_\ell(q-1),v_\ell(c(u_q))\bigr)\Bigr),
\]
with $u_q$ the unit part of the Tate parameter at $q$. Consequently, under the hypothesis of Theorem thm:main ($E$ optimal, $\ell\notin S_f$), (eq:identity) reads
\begin{multline*}
v_\ell(\degphi)=\len H^1_f(\Q,W)+\sum_{q\mid N}e_\ell(E,q)\\
+\sum_{q:\,v_q(j_E)<0}\Bigl(t_\ell(E,q)+\min\bigl(t_\ell(E,q),v_\ell(q-1),v_\ell(c(u_q))\bigr)\Bigr).
\end{multline*}
\end{corollary}

\begin{proof}
The first identity: \Cref{lem:blochkato:tamfactor} at the multiplicative primes, \Cref{prop:blochkato:tam-potmult} at the additive potentially multiplicative ones, and
\Cref{thm:blochkato:tam-additive}(i) at the potentially good additive primes; $t_\ell(E,q)=v_\ell(\nu_q)$. The second: Theorem thm:main, whose proof is assembled in
\Cref{ch:congruence} from \Cref{thm:blochkato:selmer-identity} below and the congruence-ideal theorems, with the equality of (eq:identity) for the optimal curve proved in
\Cref{ch:moddeg} (Proposition prop:pullback).
\end{proof}

This is the paper's Corollary cor:explicit. The excess of $v_\ell(\degphi)$ over the explicit local bound is the length of the adjoint Bloch-Kato Selmer group. What
remains of Open question 3 is that Selmer group.

\begin{remark}[the excess on the tables]\label{rem:blochkato:open-question-3}
We computed the excess $v_\ell(\degphi)-\bigl[\sum_qe_\ell(E,q)+\sum_{v_q(j)<0}\bigl(t_\ell+\min(t_\ell,v_\ell(q-1),v_\ell(c(u_q)))\bigr)\bigr]$ on the optimal curves of Cremona's tables of
conductor below $10{,}000$ without complex multiplication, with $\ell\nmid N$, $E[\ell]$ irreducible and, for $\ell=3$, mod $3$ image not \texttt{3Ns} (the optimal curve is the one
listed in Cremona's \texttt{degphi} file; irreducibility and the image are read from \texttt{galrep}; $e_\ell$ from PARI's \texttt{lfunsympow}, as in the paper's \texttt{rule}):
\begin{center}\small
\begin{tabular}{lrrrrrr}
\toprule
$\ell$ & curves & excess $0$ & $1$ & $2$ & $3$ & $\ge4$\\
\midrule
$3$ & $13{,}203$ & $8{,}688$ & $3{,}288$ & $942$ & $248$ & $37$\\
$5$ & $23{,}361$ & $19{,}019$ & $3{,}993$ & $328$ & $19$ & $2$\\
$7$ & $28{,}134$ & $25{,}133$ & $2{,}888$ & $111$ & $2$ & $0$\\
\bottomrule
\end{tabular}
\end{center}
The excess is never negative. By \Cref{cor:blochkato:explicit-identity} it is $\len H^1_f(\Q,W)$; so the adjoint Selmer group vanishes for $65.8\%$ of these curves at $\ell=3$,
$81.4\%$ at $\ell=5$ and $89.3\%$ at $\ell=7$. The Tamagawa factor at $q=\ell$ is not defined by \Cref{def:blochkato:tamagawa} (it needs $p$-adic Hodge theory) and belongs to Open
question 1 (\Cref{ch:congruence}).
\end{remark}

\section[Lemma lem:dfg21]{The comparison of Selmer groups: Lemma lem:dfg21 and \texorpdfstring{$H^0(\Q,W^D)$}{H0(Q,WD)}}\label{sec:blochkato:comparison}

This section proves the paper's Lemma lem:dfg21 (Lemma 2.1 of Diamond, Flach and Guo) by specialising the abstract comparison of Selmer structures of
\Cref{ch:galcoh} to the Bloch-Kato conditions, identifies its last term with $H^0(\Q,W^D)$ (Lemma lem:h0), and derives the length identity that the proof of
Theorem thm:main and Remark rem:ns use. It also states Theorem thm:dfg27, which supplies the hypothesis of the lemma.

\begin{definition}[the exceptional set]\label{def:blochkato:exceptional-set}
For $f=f_E$ (weight $2$, trivial character, rational coefficients) the \emph{exceptional set} of Diamond, Flach and Guo, as the paper reads \cite[(31)]{DFG04}, is
\[
S_f=\bigl\{\ell:\ \ell\mid2N,\ \text{or}\ \bar\rho_{E,\ell}|_{\Gal(\Qbar/\Q(\sqrt{\ell^*}))}\ \text{is not absolutely irreducible}\bigr\}.
\]
The hypothesis of Theorem thm:main is $\ell\notin S_f$.
\end{definition}

For $\ell\nmid2N$ the following are equivalent (the paper's Proposition prop:Sf, proved in \Cref{ch:galrep}): $\ell\notin S_f$; $E[\ell]$ is irreducible and $h_\ell(E)=0$; $E[\ell]$
is irreducible and, if $\ell=3$, the mod $3$ image is not $N_s(\F_3)$. In particular for $\ell\ge5$ the condition on $\Q(\sqrt{\ell^*})$ is automatic once $E[\ell]$ is irreducible.

\begin{theorem}[Diamond, Flach and Guo, Theorem 2.7]\label{thm:blochkato:dfg27}
If $\ell\notin S_f$, then $H^0(\Q,V^D)=H^1_f(\Q,V^D)=0$ and $H^0(\Q,A_f)=H^1_f(\Q,A_f)=0$.
\end{theorem}

\begin{proof}[Sketch of proof]
This is \cite[Theorem 2.7]{DFG04}, the paper's Theorem thm:dfg27; we do not prove it. The $H^0$ statements hold for every $E$ with $\bar\rho$ absolutely irreducible
(\Cref{prop:blochkato:global-invariants}). The content is the vanishing of the adjoint Selmer group with $\Ql$ coefficients, which comes from the Taylor-Wiles method:
the identification of a universal deformation ring with a Hecke algebra forces the tangent space $H^1_f(\Q,A_f)$ to vanish. One way to see the implication, given the
other theorem of Diamond, Flach and Guo that the paper uses (Theorem thm:dfg37, \Cref{ch:congruence}): it says that $H^1_\Sigma(\Q,W)$ is finite, so $H^1_f(\Q,W)\subseteq H^1_\Sigma(\Q,W)$
is finite, so $H^1_f(\Q,A_f)=0$ (\Cref{prop:blochkato:selmer-relations}(iii)), so $H^1_f(\Q,V^D)=0$ (\Cref{prop:blochkato:euler-characteristic}).
\end{proof}

\begin{proposition}[Lemma lem:h0 in Selmer form]\label{prop:blochkato:h0-selmer}
Let $\ell$ be odd and $\bar\rho$ irreducible (equivalently, absolutely irreducible: complex conjugation has distinct eigenvalues $\pm1$).
\begin{enumerate}[label=(\roman*)]
\item If $H^0(\Q,V^D)=H^1_f(\Q,V^D)=0$, then $H^1_f(\Q,T^D)=H^1(\Q,T^D)_{\mathrm{tor}}$, and the connecting map is an isomorphism $H^0(\Q,W^D)\xrightarrow{\sim}H^1_f(\Q,T^D)$.
\item (Lemma lem:h0, \Cref{ch:galrep}) $H^0(\Q,W^D)=H^0(\Q,\ad T(1)\otimes\Qlmod)$ is finite and cyclic, of length $h_\ell(E)$, the largest $k$ with $\rho\bmod\ell^k\simeq(\rho\otimes\chi_\ell)\bmod\ell^k$.
Moreover $h_\ell(E)=0$ for $\ell\ge5$, and $h_3(E)\le1$ with equality iff $\bar\rho_{E,3}\simeq\bar\rho_{E,3}\otimes\chi_{-3}$, iff the image of $\bar\rho_{E,3}$ is conjugate to $N_s(\F_3)$. This
holds for every $E$, with or without complex multiplication, and whether or not $\ell\mid N$.
\end{enumerate}
\end{proposition}

\begin{proof}
(i) By \Cref{prop:blochkato:selmer-relations}(i) for the lattice $T^D$, $H^1_f(\Q,T^D)=\iota^{-1}(H^1_f(\Q,V^D))=\ker\iota=H^1(G_{\Q,S},T^D)_{\mathrm{tor}}$, the image of $H^0(\Q,W^D)$ under the
connecting map; the connecting map is injective because $H^0(\Q,V^D)=0$. (ii) is proved in \Cref{ch:galrep}. In brief: the $\ell$-torsion of $H^0(\Q,W^D)$ is the space of
$G_\Q$-maps $\bar\rho\to\bar\rho\otimes\bar\chi_\ell$, of dimension at most one by Schur's lemma; an element of exact order $\ell^k$ is an isomorphism $\rho\bmod\ell^k\simeq(\rho\otimes\chi_\ell)\bmod\ell^k$ by
Nakayama's lemma; taking determinants gives $\chi_\ell^2\equiv1\pmod{\ell^k}$ on $G_\Q$, and as $\chi_\ell(G_\Q)=\Zl^\times\ni2$ this forces $k=0$ for $\ell\ge5$ and $k\le1$ for $\ell=3$; and
$\bar\chi_3=\chi_{-3}$.
\end{proof}

\begin{paperlink}
Lemma lem:h0 in the revised paper states (ii) and proves it by the determinant argument. The original version asserted finiteness ``for $E$ without complex
multiplication'', by an argument with $a_p=0$ at inert primes that confused the cyclotomic $\chi_\ell$ with the quadratic $\chi_{\ell^*}$; it also stated the criterion for
$h_\ell(E)>0$ with $\chi_{\ell^*}$ instead of $\bar\chi_\ell$. The two agree only for $\ell=3$. Conjecture conj:euler defines $h_\ell(E)$ as the length of
$H^0(\Q,\ad T(1)\otimes\Q_\ell/\Z_\ell)$, which is $\len H^1_f(\Q,T^D)$ by (i) whenever $H^1_f(\Q,V^D)=0$.
\end{paperlink}

\begin{lemma}[Diamond, Flach and Guo, Lemma 2.1]\label{lem:blochkato:dfg21}
Let $\ell$ be odd. Suppose $H^0(\Q,V^D)=H^1_f(\Q,V^D)=0$, and let $\Sigma$ be a nonempty finite set of primes, $\ell\notin\Sigma$, containing every prime $p\ne\ell$ at which $W$ is
ramified. Then there is an exact sequence
\[
0\to H^1_f(\Q,W)\to H^1_\Sigma(\Q,W)\to\bigoplus_{p\in\Sigma\cup\{\infty\}}H^1_f(\Q_p,T^D)^*\to H^0(\Q,W^D)^*\to0,
\]
in which the term at $\infty$ is $0$ and, for $p\in\Sigma$, $H^1_f(\Q_p,T^D)=H^1(\Q_p,T^D)_{\mathrm{tor}}\simeq H^0(\Q_p,W^D)$ is finite.
\end{lemma}

\begin{proof}
\Cref{ch:galcoh} proves the following (its abstract form of the lemma): let $T$ be unramified outside a finite set $S$ of places containing $\ell$, $\infty$ and $\Sigma$, and let
$(L_v)_v$ be local conditions on $V$ with $L_v=H^1_{\ur}(\Q_v,V)$ for $v\notin S$ and $L_v^\perp=L^D_v$ for every $v$, where $(L^D_v)$ are conditions on $V^D$; propagate them to
$W$ and to $T^D$ (preimages and images), and let $\mathcal L^\Sigma$ be the structure on $W$ obtained by dropping the condition at the primes of $\Sigma$. Assume (a) the Selmer
group of $V^D$ for $(L^D_v)$ vanishes, and (b) $H^0(\Q_p,V^D)=0$ for some $p\in\Sigma$. Then
\[
0\to H^1_{\mathcal L}(\Q,W)\to H^1_{\mathcal L^\Sigma}(\Q,W)\to\bigoplus_{p\in\Sigma\cup\{\infty\}}(L^D_{T^D,p})^*\to H^0(\Q,W^D)^*\to0
\]
is exact, with $L^D_{T^D,p}$ the propagated condition on $T^D$ at $p$. We apply this to the triple $(A_f,A,W)$ with $L_v=H^1_f(\Q_v,A_f)$ and $L^D_v=H^1_f(\Q_v,V^D)$, and check
each point.
\begin{itemize}
\item $L_v=H^1_{\ur}$ for $v\notin S$: \Cref{def:blochkato:local-conditions}(i).
\item $L_v^\perp=L^D_v$ at every place: \Cref{prop:blochkato:local-duality-f} ($v=p\ne\ell$: unramified orthogonality; $v=\ell$: Bloch and Kato, $A_f$ being crystalline for $\ell\nmid N$,
\Cref{prop:blochkato:crystalline-closure}, and in general de Rham; $v=\infty$: all groups vanish).
\item The propagated structures are the Bloch-Kato conditions on $W$ and $T^D$ (\Cref{def:blochkato:propagation}), so $H^1_{\mathcal L}(\Q,W)=H^1_f(\Q,W)$ and $L^D_{T^D,p}=H^1_f(\Q_p,T^D)$.
\item $\mathcal L^\Sigma$ is the structure of $H^1_\Sigma$: no condition on $\Sigma$, Bloch-Kato at $\ell$, and at $p\notin\Sigma\cup\{\ell\}$ the Bloch-Kato condition, which is the unramified one because
$W$ is unramified there (\Cref{lem:blochkato:selmer-S}); at $\infty$ there is nothing to impose.
\item (a) is the hypothesis $H^1_f(\Q,V^D)=0$.
\item (b) holds at every $p\in\Sigma$ (\Cref{lem:blochkato:no-local-invariants}(i)), and $\Sigma\ne\emptyset$.
\end{itemize}
The term at $\infty$ is $H^1(\R,T^D)^*=0$, a group killed by $2$ and by a power of $\ell$. The description of the local terms is \Cref{prop:blochkato:local-unramified}(v) and
\Cref{thm:blochkato:fpr-local}.
\end{proof}

What each hypothesis does, in the abstract proof of \Cref{ch:galcoh}: $H^1_f(\Q,V^D)=0$ makes the dual Selmer group $H^1_f(\Q,T^D)$ torsion, hence equal to
$H^0(\Q,W^D)$ (\Cref{prop:blochkato:h0-selmer}(i)); $\Sigma\ne\emptyset$ with $H^0(\Q_p,V^D)=0$ at some $p\in\Sigma$ makes the strict dual Selmer group (classes of $H^1_f(\Q,T^D)$ that vanish
at $\Sigma$) zero, because the local connecting map $H^0(\Q_p,W^D)\to H^1(\Q_p,T^D)$ is then injective, and this gives the surjection onto $H^0(\Q,W^D)^*$; and $\ell$ odd kills
the archimedean terms. The paper's hypothesis $H^0(\Q,V^D)=0$ is implied by (b), as the paragraph after Lemma lem:dfg21 observes: there it is automatic, and only
$H^1_f(\Q,V^D)=0$ needs Theorem thm:dfg27.

\begin{corollary}[the length identity]\label{cor:blochkato:length-identity}
Under the hypotheses of \Cref{lem:blochkato:dfg21} for the paper's objects, with $\Sigma$ the set of primes dividing $N$ and $\bar\rho$ irreducible, all terms are finite and
\[
\len H^1_\Sigma(\Q,W)=\len H^1_f(\Q,W)+\sum_{p\in\Sigma}\Bigl(v_\ell\bigl(L_p(A_f,1)^{-1}\bigr)+v_\ell\bigl(\Tam_p(\ad T)\bigr)\Bigr)-h_\ell(E).
\]
\end{corollary}

\begin{proof}
By \Cref{prop:blochkato:euler-characteristic}, $H^1_f(\Q,V^D)=0$ gives $H^1_f(\Q,A_f)=0$, so $H^1_f(\Q,W)$ is finite (\Cref{prop:blochkato:selmer-relations}(iii)); the local terms are
finite (\Cref{thm:blochkato:fpr-local}), and $H^0(\Q,W^D)$ is finite (\Cref{prop:blochkato:h0-selmer}(ii)). So $H^1_\Sigma(\Q,W)$ is finite, and in an exact sequence of finite
$\Zl$-modules the alternating sum of lengths is $0$ (\Cref{ch:cohomology}); Pontryagin duality preserves lengths. The local lengths are
$v_\ell(L_p(A_f,1)^{-1})+v_\ell(\Tam_p(T^D))$ (\Cref{thm:blochkato:fpr-local}), with $\Tam_p(T^D)=\Tam_p(\ad T)$ (\Cref{prop:blochkato:tam-dual}), and $\len H^0(\Q,W^D)=h_\ell(E)$.
\end{proof}

\begin{paperlink}
The proof of Theorem thm:main contains the display
``$\len H^1_\Sigma(\Q,W)=\len H^1_f(\Q,W)+\sum_{p\in\Sigma}\len H^1_f(\Q_p,T^D)$'' (with $H^0(\Q,W^D)=0$), and then says that by Lemma lem:fpr each local term is
$v_\ell(L_p(A_f,1)^{-1})+v_\ell(\Tam_p(A))$; Remark rem:ns contains the display with $-h_\ell(E)$. Both are
\Cref{cor:blochkato:length-identity}.
\end{paperlink}

\section[The Selmer half of Theorem thm:main]{The Selmer half of the proof of Theorem thm:main}\label{sec:blochkato:theorem}

This section states and proves the part of Theorem thm:main that is about Selmer groups: the identity between the lengths of $H^1_\Sigma(\Q,W)$ and $H^1_f(\Q,W)$ with
every local term made explicit. \Cref{ch:congruence} combines it with the congruence-ideal theorems of Diamond, Flach and Guo and with the identification of the
congruence ideal with the modular degree (\Cref{ch:moddeg}) to prove Theorem thm:main.

\begin{theorem}[the Selmer identity]\label{thm:blochkato:selmer-identity}
Let $E/\Q$ be an elliptic curve of conductor $N$, $\ell\notin S_f$, and $\Sigma$ the set of primes dividing $N$. Then $H^1_f(\Q,W)$ and $H^1_\Sigma(\Q,W)$ are finite and
\begin{multline*}
\len H^1_\Sigma(\Q,W)=\len H^1_f(\Q,W)+\sum_{q\mid N}v_\ell\bigl(L_q(A_f,1)^{-1}\bigr)\\
+\sum_{q:\,v_q(j_E)<0}\Bigl(t_\ell(E,q)+\min\bigl(t_\ell(E,q),v_\ell(q-1),v_\ell(c(u_q))\bigr)\Bigr).
\end{multline*}
Here $v_\ell(L_q(A_f,1)^{-1})=v_\ell(q^2-1)$ at the potentially multiplicative primes and $e_\ell(E,q)$ at the potentially good additive primes. Optimality of $E$ is not used.
\end{theorem}

\begin{proof}
$\ell\notin S_f$ means $\ell\nmid2N$ and $\bar\rho$ absolutely irreducible over $\Q(\sqrt{\ell^*})$, hence over $\Q$. $\Sigma\ne\emptyset$ because no elliptic curve over $\Q$ has good reduction
everywhere (\Cref{ch:reduction}), and $W$ is unramified outside $\Sigma\cup\{\ell\}$. By \Cref{thm:blochkato:dfg27} the hypotheses of \Cref{lem:blochkato:dfg21} hold. By
\Cref{prop:blochkato:h0-selmer}(ii) and the equivalences after \Cref{def:blochkato:exceptional-set}, $h_\ell(E)=0$. \Cref{cor:blochkato:length-identity} gives the identity with
$\sum_{p\in\Sigma}v_\ell(\Tam_p(\ad T))$, which \Cref{cor:blochkato:explicit-identity} (its first part, which is proved in this chapter) makes explicit. The values of the Euler terms
are \Cref{lem:blochkato:euler-integral}.
\end{proof}

\begin{remark}[the assembly, in \Cref{ch:congruence}]\label{rem:blochkato:assembly}
For the optimal curve and $\ell\notin S_f$: $v_\ell(\degphi)=v_\lambda(\eta^\emptyset_{f,\lambda})$ (Proposition prop:pullback and Lemma lem:hecke, \Cref{ch:moddeg}; an equality with no
multiplicity-one hypothesis); $v_\lambda(\eta^\Sigma_{f,\lambda})=v_\lambda(\eta^\emptyset_{f,\lambda})+\sum_{p\in\Sigma}v_\ell(L^{\nv}_p(A_f,1)^{-1})$ (Theorem thm:dfg14); and $v_\lambda(\eta^\Sigma_{f,\lambda})=\len H^1_\Sigma(\Q,W)$
(Theorem thm:dfg37). With \Cref{thm:blochkato:selmer-identity}, the naive factors cancel the true ones at the multiplicative primes ($\delta_p=1$), and at the additive primes
($\delta_p=0$) the naive factor is $1$, so the surviving Euler terms are $\sum_qe_\ell(E,q)$ (Definition def:terms). The result is (eq:identity), with equality, and (eq:thm)
follows from $\len H^1_f(\Q,W)\ge0$ and $\min(\dots)\ge0$, including the primes of type $\kod{I}_n^*$ at $q=2$ (\Cref{prop:blochkato:tam-potmult}). These steps are proved in
\Cref{ch:congruence}; they are recalled here only to show where \Cref{thm:blochkato:selmer-identity} enters.
\end{remark}

\begin{paperlink}
The proof of Theorem thm:main, from the sentence that invokes Theorem thm:dfg27 for the hypotheses of Lemma lem:dfg21 to the display of $v_\ell(\degphi)$, is \Cref{thm:blochkato:selmer-identity}. The final step, at the potentially multiplicative primes of additive
reduction, uses \Cref{prop:blochkato:tam-potmult}: the revised paper applies Lemma lem:tamfactor to the local twist, as here; the original applied it to the global twist
$E_{\{p\}}$ of Lemma lem:twist, which does not exist at $p=2$.
\end{paperlink}

\begin{example}[four curves at $\ell=3$]\label{ex:blochkato:selmer-examples}
For the following optimal curves, $3\nmid N$, $E[3]$ is irreducible, the mod $3$ image is not \texttt{3Ns}, and $E$ has no CM, so $3\notin S_f$. The explicit local bound is
$\sum_qe_3(E,q)+\sum_{v_q(j)<0}\bigl(t_3+\min(t_3,v_3(q-1),v_3(c(u_q)))\bigr)$; by \Cref{cor:blochkato:explicit-identity} the Selmer length is $v_3(\degphi)$ minus it.
\begin{center}\small
\begin{tabular}{lllrrr}
\toprule
curve & local data & & bound & $\degphi$ & $\len H^1_f(\Q,W)$\\
\midrule
\curve{121a1} & $\kod{II}$ at $11$, $e_3=v_3(12)=1$ & & $1$ & $6$ & $0$\\
\curve{791a1} & $\kod{I}_3$ at $7$ ($2$), $\kod{I}_2$ at $113$ ($0$) & & $2$ & $144$ & $0$\\
\curve{1184f1} & $\kod{III}^*$ at $2$ ($0$), $\kod{I}_3$ at $37$ ($2$) & & $2$ & $1728$ & $1$\\
\curve{1748e1} & $\kod{IV}$ at $2$ ($e_3=1$), $\kod{I}_6$ at $19$ ($2$), $\kod{I}_1$ at $23$ ($0$) & & $3$ & $3240$ & $1$\\
\bottomrule
\end{tabular}
\end{center}
For \curve{791a1} the bound (eq:thm) is $1$, and the Tamagawa excess at $7$ is needed for equality. The numbers in parentheses are the Tamagawa contributions
(\Cref{ex:blochkato:tamfactor-examples}).
\end{example}

\section[Beyond the hypothesis]{Beyond the hypothesis: Remark rem:ns, \texorpdfstring{$h_\ell(E)$}{h(E)}, Open question 2}\label{sec:blochkato:beyond}

When $\ell=3$ and the mod $3$ image is $N_s(\F_3)$, the hypothesis $\ell\notin S_f$ fails, $H^0(\Q,W^D)$ is not zero, and the paper's Remark rem:ns records what the
exact sequence of Lemma lem:dfg21 then says. This section makes the conditional statement precise, identifies the correction term $h_\ell(E)$ of Conjecture
conj:euler, and states what is missing for Open question 2.

\begin{proposition}[Remark rem:ns]\label{prop:blochkato:rem-ns}
Let $\ell=3$, $3\nmid N$, and suppose the image of $\bar\rho_{E,3}$ is conjugate to $N_s(\F_3)$. Then $H^0(\Q,V^D)=0$, $h_3(E)=1$, and, \emph{if} $H^1_f(\Q,V^D)=0$ (equivalently
$H^1_f(\Q,\ad V_3E)=0$), the groups $H^1_f(\Q,W)$ and $H^1_\Sigma(\Q,W)$ are finite and
\[
\len H^1_\Sigma(\Q,W)=\len H^1_f(\Q,W)+\sum_{p\in\Sigma}\Bigl(v_3\bigl(L_p(A_f,1)^{-1}\bigr)+v_3\bigl(\Tam_p(\ad T)\bigr)\Bigr)-1 .
\]
\end{proposition}

\begin{proof}
$\bar\rho$ is irreducible (the image $N_s(\F_3)$ fixes no line), $H^0(\Q,V^D)=0$ (\Cref{prop:blochkato:global-invariants}(ii)), $h_3(E)=1$ (\Cref{prop:blochkato:h0-selmer}(ii)), and
$\Sigma\ne\emptyset$. The equivalence is \Cref{prop:blochkato:euler-characteristic}. Under the condition, \Cref{lem:blochkato:dfg21} and \Cref{cor:blochkato:length-identity} apply.
\end{proof}

Theorem thm:dfg27 does not supply the vanishing of $H^1_f(\Q,V^D)$ here, since $3\in S_f$; the proposition states it as a hypothesis, as Remark rem:ns of the revised
paper does (``we assume it here''). It is predicted by the Bloch-Kato conjecture (\Cref{conj:blochkato:tnc}), since $L(\ad f,1)\ne0$; we know of no proof for these
curves and have not searched the literature after Diamond, Flach and Guo for one. Section 1.3 of the revised paper also states the condition.

\begin{remark}[the correction term of Conjecture conj:euler]\label{rem:blochkato:conj-euler-h}
In Conjecture conj:euler the correction is $h_\ell(E)=\len H^0(\Q,\ad T(1)\otimes\Qlmod)$: by \Cref{prop:blochkato:h0-selmer} it is $0$ for $\ell\ge5$ and is $1$ exactly for the curves
with mod $3$ image $N_s(\F_3)$, for every conductor. It is the term $H^0(\Q,W^D)^*$ of \Cref{lem:blochkato:dfg21}, and it enters \Cref{prop:blochkato:rem-ns} with the sign $-1$,
the sign with which it appears in (eq:rule). The paper's function \texttt{ns3(E)} (two quadratic factors of the $3$-division polynomial) computes it (\Cref{ch:galrep}).
Example: \curve{726a1} $=[1,1,0,-35,-51]$, with $3\mid N$, so outside every theorem; its $3$-division polynomial factors as $(x^2-x-74)(3x^2+8x+20)$, so $h_3=1$; it has
type $\kod{III}$ at $11$ with $e_3=v_3(12)=1$, $\degphi=144$, and the paper's \texttt{rule(E,3)} and \texttt{rule(E,3,1)} both return $[2,2]$: in the second the type
$\kod{III}$ term and the correction cancel.
\end{remark}

\begin{remark}[Open question 2]\label{rem:blochkato:open-question-2}
Open question 2 asks to prove the corrected identity of Remark rem:ns for the curves with image $N_s(\F_3)$, that is, to establish $H^1_f(\Q,\ad V_3E)=0$ and
$v_\lambda(\eta^\Sigma_{f,\lambda})=\len H^1_\Sigma(\Q,W)$ in this dihedral case. On the Selmer side \Cref{prop:blochkato:rem-ns} shows what the identity would give and that the
only missing input is $H^1_f(\Q,A_f)=0$; on the congruence side the missing input is a Taylor-Wiles theorem for a residual representation induced from $\Q(\sqrt{-3})$
(\Cref{ch:congruence}). The data of Section sec:anomaly (failures of (eq:thm) by exactly one, never more) are consistent with the identity with $-h_3(E)$.
\end{remark}

\section[Section sec:tam]{Section sec:tam: testing the Tamagawa factors on the tables}\label{sec:blochkato:sectam}

The paper's Section sec:tam tests the refinement of (eq:thm) by the Tamagawa excess of \Cref{cor:blochkato:tamfactor-consequences}. This section reproduces the test for
conductors below $10{,}000$ and says what it does and does not test.

\begin{example}[the test]\label{ex:blochkato:sec-tam}
The population: optimal curves without CM satisfying the hypothesis of Theorem thm:main ($\ell\nmid N$, $E[\ell]$ irreducible, image not \texttt{3Ns} at $\ell=3$) that have a
potentially multiplicative prime $q\equiv1\pmod\ell$ with $\ell\mid\nu_q$. For each, compare $v_\ell(\degphi)$ with the bound $\sum_qe_\ell+\sum_qt_\ell$ of (eq:thm) and with the
refined bound that adds the Tamagawa excesses. The paper's population is the working set (conductors up to $500{,}000$) with primes of type $\kod{I}_n$ only; ours is
conductors below $10{,}000$ with primes of type $\kod{I}_n$ or $\kod{I}_n^*$ (\Cref{prop:blochkato:tam-potmult}):
\begin{center}\small\setlength{\tabcolsep}{4pt}
\begin{tabular}{llrrrr}
\toprule
source & $\ell$ & curves & positive excess & equality, bound (eq:thm) & equality, refined bound\\
\midrule
paper & $3$ & $46{,}222$ & $12{,}095$ & $20{,}822$ ($45.0\%$) & $28{,}007$ ($60.6\%$)\\
paper & $5$ & $10{,}940$ & $1{,}646$ & $7{,}333$ ($67.0\%$) & $8{,}610$ ($78.7\%$)\\
here & $3$ & $1{,}050$ & $186$ & $596$ ($56.8\%$) & $714$ ($68.0\%$)\\
here & $5$ & $252$ & $24$ & $203$ ($80.6\%$) & $220$ ($87.3\%$)\\
here & $7$ & $18$ & $1$ & $17$ & $18$\\
\bottomrule
\end{tabular}
\end{center}
No curve violates the refined inequality, in these populations or in the full populations of \Cref{rem:blochkato:open-question-3} ($13{,}203$, $23{,}361$ and $28{,}134$ curves). The
rates are not comparable with the paper's (smaller conductors; $\kod{I}_n^*$ primes included); the revised paper also reports our figures at $\ell=3,5$.
\end{example}

\begin{remark}[what Section sec:tam tests]\label{rem:blochkato:what-sec-tam-tests}
Given \Cref{thm:blochkato:selmer-identity} and its assembly with the equality $v_\ell(\degphi)=v_\lambda(\eta^\emptyset_{f,\lambda})$ (\Cref{ch:moddeg}, \Cref{ch:congruence}), a violation of the
refined inequality is impossible, and the proportion of equality is the proportion of curves with $H^1_f(\Q,W)=0$. So Section sec:tam is a consistency check of the
Tamagawa term of (eq:identity) and of the computation of Tate parameters, not an independent test of a conjecture. The revised paper's phrasing (``the extra
divisibility of the modular degree predicted by the Tamagawa factor of the adjoint is always present, the Selmer length accounting for the rest'') says this; the
original said that the extra divisibility occurs ``exactly where the Tamagawa factor of the adjoint exceeds $\ell^{t_\ell}$'', which is too strong, since extra
divisibility also comes from the Selmer group.
\end{remark}

\section[The Tamagawa number conjecture]{The Tamagawa number conjecture and the adjoint motive}\label{sec:blochkato:motive}

This section places Theorem thm:main in the framework of Bloch and Kato. Nothing here is used in a proof. The adjoint motive of $f$, its $L$-function, the shape of the
Tamagawa number conjecture, the theorem of Diamond, Flach and Guo that proves it for the adjoint at $\ell\notin S_f$, and Flach's earlier finiteness theorem are stated with
references, and a labelled heuristic explains why (eq:identity) is the conjecture read through the modular degree.

\begin{definition}[the adjoint motive]\label{def:blochkato:adjoint-motive}
Let $f$ be a weight-$2$ newform with rational coefficients. The \emph{adjoint motive} $\ad M_f$ is the motive with $\ell$-adic realisations $A_f=\ad V_\ell(f)\simeq\Sym V_\ell(f)\otimes\det^{-1}$
(\Cref{ch:reps}), and Betti and de Rham realisations built from those of $M_f$ in the same way; it is pure of weight $0$ with Hodge numbers $h^{1,-1}=h^{0,0}=h^{-1,1}=1$, and
\[
L(\ad f,s)=\prod_pL_p(A_f,s)=L(\Sym f,s+1),
\]
with the true Euler factors (at $p=\ell$ computed through $\Dcris$, or with $V_{\ell'}$ for a prime $\ell'\ne p$) and the normalisation $L_p(\Sym f,s)=L_p(A_f,s-1)$ of
Appendix B, B-C7. The point $s=1$ is critical in the sense of Deligne \cite{Deligne79}, and $L(\ad f,1)\ne0$.
\end{definition}

The last assertion follows from \Cref{ch:lfunctions}: the formula $\langle f,f\rangle=N\,L^{\nv}(\Sym f,2)/(8\pi^3)$ of Rankin and Shimura, quoted in the paper's Remark rem:degree, shows
$L^{\nv}(\ad f,1)>0$, and $L(\ad f,1)$ differs from it by the nonzero factors $L_q(A_f,1)/L^{\nv}_q(A_f,1)$ (\Cref{lem:blochkato:no-local-invariants}). For the critical value and the
Hodge numbers see \cite{Deligne79} and \cite[Section 1]{DFG04}; for the symmetric square, Shimura \cite{Shimura76} and Gelbart and Jacquet \cite{GelbartJacquet}.

\begin{conjecture}[Bloch and Kato, in outline]\label{conj:blochkato:tnc}
Let $M$ be a motive over $\Q$ whose $L$-function is critical and nonzero at $s=0$, $\ell$ a prime, $T$ a stable lattice in the $\ell$-adic realisation of $M$, $W_M=T\otimes\Qlmod$ and
$W_M^D=T^D\otimes\Qlmod$. Then, for a suitably normalised period $\Omega(M)$,
\[
\frac{L(M,0)}{\Omega(M)}\ \sim_\ell\ \frac{\prod_{p\ne\ell}\#\bigl(\Zl/\Tam_p(T)\bigr)\cdot\#H^1_f(\Q,W_M^D)}{\#H^0(\Q,W_M)\cdot\#H^0(\Q,W_M^D)},
\]
where $\sim_\ell$ means equality up to an $\ell$-adic unit, and the contributions of $\ell$ and $\infty$ are absorbed in $\Omega(M)$ (the Tamagawa measure of Bloch and Kato).
\end{conjecture}

This is a schematic form of the conjecture of Bloch and Kato \cite[Section 5]{BlochKato}, in the formulation of Fontaine and Perrin-Riou \cite{FPR94}; the precise
normalisations of the period and of the terms at $\ell$ differ between these sources and \cite[Section 2]{DFG04}, and we do not reconcile them. The book never uses the
conjecture in a proof. For $M=\ad M_f(1)$, whose realisation is $A_f(1)=V^D$ with lattice $T^D$, one has $L(M,0)=L(\ad f,1)$, $W_M=W^D$ and $W_M^D=W$; so the right side is
\[
\frac{\prod_{p\ne\ell}\#(\Zl/\Tam_p(T^D))\cdot\#H^1_f(\Q,W)}{\#H^0(\Q,W^D)\cdot\#H^0(\Q,W)}:
\]
the adjoint Selmer group of Theorem thm:main, the Tamagawa factors of Lemma lem:fpr, and $h_\ell(E)$ in the denominator, with the same sign as in Remark rem:ns.

\begin{example}[Birch and Swinnerton-Dyer]\label{ex:blochkato:bsd}
For $M=h^1(E)(1)$, with realisation $V_\ell E$ and $L(M,0)=L(E,1)$, \Cref{conj:blochkato:tnc} is the $\ell$-part of the formula of Birch and Swinnerton-Dyer in rank $0$:
\begin{center}\small
\begin{tabular}{ll}
\toprule
term of \Cref{conj:blochkato:tnc} & term of the BSD formula\\
\midrule
$\Omega(M)$ & the real period $\Omega_E$\\
$\#(\Zl/\Tam_p(T_\ell E))$ & $\ell$-part of $c_p$ (\Cref{prop:blochkato:tam-of-E})\\
$\#H^1_f(\Q,W_M^D)$, $W_M^D\simeq E[\ell^\infty]$ & $\#\Sha(E/\Q)[\ell^\infty]$ (rank $0$; \Cref{ex:blochkato:classical-selmer})\\
$\#H^0(\Q,W_M)\cdot\#H^0(\Q,W_M^D)$ & $\#E(\Q)[\ell^\infty]^2$\\
\bottomrule
\end{tabular}
\end{center}
Here $T^D\simeq T$ by the Weil pairing (\Cref{prop:blochkato:twists}(iii)). This is orientation only; see \cite[Section 6]{BlochKato} and \Cref{ch:ellcurves}.
\end{example}

\begin{theorem}[Diamond, Flach and Guo, Theorem 0.1]\label{thm:blochkato:dfg-01}
Let $f$ be a newform (of weight $k\ge2$) and $\lambda$ a prime of its coefficient field not in the exceptional set $S_f$. Then the $\lambda$-part of the Tamagawa number conjecture
of Bloch and Kato holds for the adjoint motive of $f$ at the critical point $s=1$ \cite[Theorem 0.1]{DFG04}.
\end{theorem}

\begin{proof}[Reference]
We do not prove this; it is the theorem the paper cites when it says that its hypothesis ``is the one under which Diamond, Flach and Guo prove the Tamagawa
number conjecture for the adjoint motive of $f$''. The exact normalisation of the period is in \cite[Section 1]{DFG04}. The proof combines the identification of
the congruence ideal with the special value (Hida, and the Rankin-Shimura method), Theorem thm:dfg37, Theorem thm:dfg14, and the local computations of Lemma
lem:fpr; \Cref{ch:congruence} gives the parts the paper uses.
\end{proof}

\begin{theorem}[Diamond, Flach and Guo, Theorem 0.2]\label{thm:blochkato:dfg-02}
If $\ell\notin S_f$ and $\Sigma$ is a finite set of primes, $\ell\notin\Sigma$, outside which $M_{f,\lambda}$ is minimally ramified (for $f=f_E$: any $\Sigma$ containing the primes dividing $N$),
then $H^1_\Sigma(\Q,W)$ is finite of length $v_\lambda(\eta^\Sigma_{f,\lambda})$.
\end{theorem}

\begin{proof}[Reference]
This is the paper's Theorem thm:dfg37, which the paper says is ``the form given as Theorem 0.2'' of \cite{DFG04}. We do not prove it; the statement, the meaning of
``minimally ramified'' and the outline of the Taylor-Wiles argument are in \Cref{ch:congruence}. It is recalled here only so that \Cref{thm:blochkato:selmer-identity} can be
read against it.
\end{proof}

\begin{theorem}[Flach, in outline]\label{thm:blochkato:flach}
Let $E/\Q$ be modular. For every prime $\ell$ outside an explicit finite set of exceptional primes (which contains the primes dividing $6N$ and those for which $\bar\rho_{E,\ell}$ is not
surjective, among others), the Bloch-Kato Selmer group of $\Sym T_\ell E\otimes\Qlmod$ is finite, and it is annihilated by an explicit integer built from the modular
degree \cite{Flach92}.
\end{theorem}

\begin{proof}[Reference]
We do not prove this, and we do not reproduce the exact list of exceptional primes or the annihilator; see \cite{Flach92}. Note that $\Sym V_\ell E\simeq A_f(1)=V^D$ (since
$\ad V\simeq\Sym V\otimes\det^{-1}$ and $\det=\chi_\ell$), so Flach's group is $H^1_f(\Q,W^D)$ in the paper's notation; when $\bar\rho$ is absolutely irreducible its finiteness is
equivalent to $H^1_f(\Q,V^D)=0$ (\Cref{prop:blochkato:selmer-relations}(iii)), hence to $H^1_f(\Q,A_f)=0$ (\Cref{prop:blochkato:euler-characteristic}), hence to the finiteness of
$H^1_f(\Q,W)$.
\end{proof}

\begin{paperlink}
Section 1 of the paper: ``Flach's finiteness theorem for the symmetric square \cite{Flach92} is the first place where the modular degree meets the Selmer group
of $\Sym f$''.
\end{paperlink}

\begin{remark}[the modular degree as an adjoint $L$-value: a heuristic]\label{rem:blochkato:degree-as-l-value}
The following is a heuristic, not a proof. By the formula of Zagier and Watkins, $\degphi\cdot\operatorname{vol}(\Lambda_E)=4\pi^2c_E^2\langle f,f\rangle$ (\Cref{ch:moddeg}); by the
Rankin-Shimura method $\langle f,f\rangle$ is $L^{\nv}(\Sym f,2)=L^{\nv}(\ad f,1)$ times explicit factors (\Cref{ch:lfunctions}). The paper's Remark rem:degree turns this into
$v_\ell(\degphi)=v_\ell(B_E)+\sum_qe_\ell(E,q)$ for $\ell\nmid N$, with $B_E=c_E^2N\,L(\ad f,1)/(2\pi\operatorname{vol}(\Lambda_E))$, so the ratio of the naive to the true local factors produces the
terms $e_\ell(E,q)$. If one reads $2\pi\operatorname{vol}(\Lambda_E)/(c_E^2N)$ as the period $\Omega(\ad M_f(1))$, then \Cref{conj:blochkato:tnc}, proved by \Cref{thm:blochkato:dfg-01}, predicts
$v_\ell(B_E)=\len H^1_f(\Q,W)+\sum_qv_\ell(\Tam_q(\ad T))$, which is (eq:identity). Theorem thm:main proves this identity through congruence ideals instead of $L$-values, so the
identification of the periods is never needed; whether this reading of the period agrees with the normalisation of \cite[Section 1]{DFG04} we have not checked.
\end{remark}

\section{Exercises}\label{sec:blochkato:exercises}

\begin{exercise}[the two Frobenii]\label{ex:blochkato:ex-frobenius}
Let $p\ne\ell$. Compute $\det(1-\Frob_p^{-1}p^{-s}\mid\Ql(1))^{-1}$ and $\det(1-\Frob_pp^{-s}\mid\Ql(1))^{-1}$. At a prime of multiplicative reduction, show that the geometric
Frobenius gives $\det(1-\Frob_p^{-1}\mid(V^D)^{I_p})=1-p^{-2}$, which is the paper's $L_p(A_f,1)^{-1}=L_p(A_f(1),0)^{-1}$, while the arithmetic one gives $1-p^2$, a number with the
same $\ell$-adic valuation.
\emph{Answer:} $(1-p^{-1-s})^{-1}$ and $(1-p^{1-s})^{-1}$; on $(V^D)^{I_p}=\Ql e\otimes\Ql(1)$ the arithmetic Frobenius acts by $p\cdot p=p^2$; and $1-p^2=-p^2(1-p^{-2})$.
\end{exercise}

\begin{exercise}[the trace form at $\ell=2$]\label{ex:blochkato:ex-gram}
Compute the Gram matrix of the trace form of $\ad T$ in the basis $e,h,f$. Show that for $\ell=2$ the map $A\to A^\vee$, $x\mapsto\tr(x\,\cdot)$, is injective with cokernel of order $2$,
and find an element of $A^\vee$ not in its image.
\emph{Hint:} the dual basis vector $h^\vee$ is $\tfrac12\tr(h\,\cdot)$.
\end{exercise}

\begin{exercise}[local ranks]\label{ex:blochkato:ex-local-ranks}
\sloppy For each row of the table of \Cref{prop:blochkato:local-adjoint}, find $\dim H^0(\Q_p,V^D)$, $\dim H^0(\Q_p,A_f)$ and, granting the local Euler characteristic formula, the
$\Zl$-rank of $H^1(\Q_p,T^D)$. Which bad primes have $H^1(\Q_p,T^D)$ infinite? Check in each case the paper's statement that $H^1_f(\Q_p,T^D)$ is the torsion of $H^1(\Q_p,T^D)$.
\emph{Answer:} $H^0(\Q_p,V^D)=0$ always; the rank is $1$ in the row $\{\pm1\}$ (the eigenvalue $1$ has multiplicity one, as at a good prime) and in the cyclic row with $c=+1$, and
$0$ otherwise. The statement holds in all rows by \Cref{prop:blochkato:local-unramified}(v).
\end{exercise}

\begin{exercise}[multiplicative primes]\label{ex:blochkato:ex-mult}
Let $q$ be a prime of multiplicative reduction, $\ell\ne q$ odd, and $W=\ad T\otimes\Qlmod$. Show that $H^1_f(\Q_q,W)=0$; that $W^{I_q}=(\Qlmod)\,e$ if $\ell\nmid n_q$, so that
$H^1_{\ur}(\Q_q,W)=0$; and that if $\ell\mid n_q$, $W^{I_q}$ contains $\ell^{-v}\Zl h/\Zl h$ modulo $(\Qlmod)e$, with $v=v_\ell(n_q)$. For $W=E[\ell^\infty]$ with split reduction and $\ell\mid n_q$, show
$H^1_{\ur}(\Q_q,W)\simeq\Z/\ell^{v}$ while $H^1_f(\Q_q,W)=0$.
\emph{Hint:} for $x=ae+bh+cf$, $(t-1)x=-(2n'b+n'^2c)e+n'c\,h$; and multiplication by $q-1\ne0$ is surjective on $\Qlmod$. Compare \Cref{ex:blochkato:ur-versus-f}.
\end{exercise}

\begin{exercise}[the lemma for \curve{14a1}]\label{ex:blochkato:ex-14a1}
The curve \curve{14a1} $=[1,0,1,4,-6]$ has split multiplicative reduction at $7$ with $n_7=3$, and the unit part of its Tate parameter reduces to $6\equiv-1\pmod7$. Check the
paper's statement $H^1(I_q,\ad T)_{\mathrm{tor}}\simeq(\Z/\ell^{v_\ell(n)})^2$ for $q=7$, $\ell=3$; compute the Frobenius invariants and show $v_3(\Tam_7(\ad T))=2$. Explain why $\degphi=1$ is no
contradiction with (eq:identity).
\emph{Answer:} $-1=3^3$ in $\F_7$ is a cube and $v_3(6)=1$, so the minimum is $1$. The curve has $E[3]$ reducible (its torsion is $\Z/6$), so Theorem thm:main does not apply,
while \Cref{lem:blochkato:tamfactor} needs no irreducibility.
\end{exercise}

\begin{exercise}[the indeterminacy of $c'$]\label{ex:blochkato:ex-cprime}
In the situation of \Cref{conv:blochkato:tate-data}, show directly that changing the Frobenius lift by $t^a$ and the cocycle $\kappa$ by a coboundary changes $c'$ by an element of
$(n,q-1)\Zl$, and give an example with $\ell\mid n$, $\ell\mid q-1$ in which the change is not a multiple of $q-1$. Show that $\min(v_\ell(n),v_\ell(q-1),v_\ell(c'))$ does not change. Why does
``the image of $\kappa$ in $H^1(G/I,\Zl(1))$'' make sense for $\Q_q$ but not for a ramified extension of $\Q_q$ without a choice?
\emph{Hint:} the change under $t^a$ is $\varepsilon qan'$; take $v_\ell(n)<v_\ell(q-1)$. A uniformiser is needed to split $H^1(K,\Zl(1))=\widehat{K^\times}$ into a unit part and a valuation
part, and $q$ is a canonical uniformiser of $\Q_q$ only.
\end{exercise}

\begin{exercise}[$\Tam(T)$ and $c_q$]\label{ex:blochkato:ex-tam-T}
Check the last sentence of the proof of Lemma lem:tamfactor (``the same computation for $M=T$ gives \dots the $\ell$-part of the Tamagawa number $c_q(E)$'') for \curve{11a1} at $q=11$,
$\ell=5$ (split, $n=5$, $c_{11}=5$) and for \curve{14a1} at $q=2$, $\ell=3$ (nonsplit, $n=6$, $c_2=2$).
\emph{Answer:} $(\Z/5)\bar e_1$ with Frobenius acting by $\varepsilon=1$: invariants of order $5$; $(\Z/3)\bar e_1$ with Frobenius acting by $-1$: invariants $0$, and $3\nmid c_2$.
\end{exercise}

\begin{exercise}[$\Tam(A(1))=\Tam(A)$ by hand]\label{ex:blochkato:ex-tam-dual}
\sloppy Let $v\ge0$, $\gamma,\delta\in\Z/\ell^v$, and $m=\min(v,v_\ell(\gamma),v_\ell(\delta))$. Show that both $\{(x,y):\gamma y=\delta y=0\}$ and $\{(x,y):\gamma x-\delta y=0\}$ in $(\Z/\ell^v)^2$ are isomorphic to
$\Z/\ell^v\oplus\Z/\ell^m$. Deduce the structure statement of \Cref{prop:blochkato:tam-dual}(b), and show that the two invariant modules are in general different subgroups of
$(\Z/\ell^v)\bar e\oplus(\Z/\ell^v)\bar h$.
\emph{Hint:} reduce to $\gamma=\ell^m$ or $\delta=\ell^m$ by a unit change of variables (Smith form of a $1\times2$ matrix).
\end{exercise}

\begin{exercise}[type $\kod{IV}$ at $\ell=3$]\label{ex:blochkato:ex-c3}
Let $\sigma=\sm0{-1}1{-1}$. Verify the three coordinate vectors of \Cref{prop:blochkato:ad-free-c3}, and show that $H^1(\langle\sigma\rangle,T)\simeq\F_3$ while
$H^1(\langle\sigma\rangle,\ad T)=0$. The curve \curve{50a1} $=[1,0,1,-1,-2]$ has type $\kod{IV}$ at $5$ and $c_5=3$: deduce from \Cref{prop:blochkato:tam-of-E} that Frobenius acts
trivially on $H^1(I_5,T_3E)_{\mathrm{tor}}\simeq\F_3$, so $v_3(\Tam_5(T_3E))=1$, while $\Tam_5(\ad T_3E)=\Z_3$.
\end{exercise}

\begin{exercise}[the refined bound for \curve{791a1}]\label{ex:blochkato:ex-791}
The curve \curve{791a1} $=[1,0,1,-31,117]$ has conductor $7\cdot113$, type $\kod{I}_3$ at $7$ and $\kod{I}_2$ at $113$, and $\degphi=144$. Show that the unit part of its Tate parameter
at $7$ reduces to a cube in $\F_7^\times$, so $v_3(\Tam_7(\ad T))=2$; that the bound (eq:thm) at $\ell=3$ is $1$ and the refined bound is $2=v_3(144)$; and conclude, granting
Theorem thm:main, that $H^1_f(\Q,\ad T_3E\otimes\Q_3/\Z_3)=0$.
\emph{Hint:} the unit part is the reduction of $j_E^{-1}7^3$; it is $6=-1$.
\end{exercise}

\begin{exercise}[an Euler term]\label{ex:blochkato:ex-121}
For \curve{121a1} at $p=11$ (type $\kod{II}$) and $\ell=3$, check the paper's statement in Lemma lem:fpr that $L_p(A_f,1)^{-1}$ is an $\ell$-adic integer of valuation $e_\ell(E,p)$.
\emph{Answer:} $\chi_{-3}(11)=-1$, so $L_{11}(A_f,1)^{-1}=1+\tfrac1{11}=\tfrac{12}{11}$, of $3$-adic valuation $1$.
\end{exercise}

\begin{exercise}[the right end of Lemma lem:dfg21]\label{ex:blochkato:ex-surj}
Prove that the last map of the sequence of \Cref{lem:blochkato:dfg21} is surjective, using the abstract sequence of \Cref{ch:galcoh} and the injectivity of
$H^0(\Q_p,W^D)\to H^1(\Q_p,T^D)$ for $p\in\Sigma$. Then consider $\ell=3$ and $T=\Z_3(\chi_{-3})$ (the character $\chi_{-3}$ on $\Z_3$): show $H^0(\Q,W^D)\simeq\Z/3$ for $W^D=\mu_{3^\infty}\otimes\chi_{-3}$, and
explain why the sequence cannot be exact with $\Sigma=\emptyset$ and the Bloch-Kato conditions.
\emph{Hint:} $\chi_3^{\mathrm{cyc}}\chi_{-3}\equiv1\pmod3$ but not modulo $9$; with $\Sigma=\emptyset$ the middle map would be the identity of $H^1_f(\Q,W)$, and the sequence would force $H^0(\Q,W^D)=0$.
\end{exercise}

\begin{exercise}[$h_\ell(E)$]\label{ex:blochkato:ex-h0}
Show by the determinant argument that $H^0(\Q,W^D)=0$ for $\ell\ge5$ and has order at most $3$ for $\ell=3$, for every $E$ with $\bar\rho$ irreducible. For \curve{726a1}, use the factorisation
of its $3$-division polynomial to show $h_3=1$. Why is the original argument, which used $a_p=0$ at the primes inert in $\Q(\sqrt{\ell^*})$, not the right one?
\emph{Hint:} $\rho\simeq\rho\otimes\chi_\ell$ would give $pa_p=a_p$ at good $p$; the relevant twist is by the cyclotomic character, not by the quadratic character $\chi_{\ell^*}$.
\end{exercise}

\begin{exercise}[the Euler characteristic]\label{ex:blochkato:ex-euler}
Verify the four rows of the table of \Cref{prop:blochkato:euler-characteristic} and deduce $H^1_f(\Q,A_f)=0\iff H^1_f(\Q,V^D)=0$. Then apply the general formula to $V^D$ instead of $A_f$
and check that it gives the same equivalence. Where is the sign convention for Hodge-Tate weights used?
\emph{Hint:} for $V^D$ the term at $\ell$ is $\dim H^1_f(\Q_\ell,V^D)-0=2$ and the term at $\infty$ is $-\dim(V^D)^{c=1}=-2$ ($c$ acts on $\Zl(1)$ by $-1$).
\end{exercise}

\begin{exercise}[Section sec:tam on a small range]\label{ex:blochkato:ex-small-range}
Using the expansion $q_E=j^{-1}+744j^{-2}+\cdots$ to obtain the unit parts, reproduce the test of \Cref{ex:blochkato:sec-tam} for conductors at most $2{,}000$ and $\ell=3$: among the
$2{,}158$ optimal non-CM curves with $3\nmid N$, $E[3]$ irreducible and image not \texttt{3Ns}, find the $128$ with a potentially multiplicative prime $q\equiv1\pmod3$, $3\mid\nu_q$; check that
equality rises from $87$ to $102$ curves, and that $19$ have a positive Tamagawa excess.
\end{exercise}

\begin{exercise}[the conjecture and the identity]\label{ex:blochkato:ex-bsd}
(Harder.) Show that \Cref{conj:blochkato:tnc} for $M=h^1(E)(1)$ reproduces the local and torsion terms of the Birch and Swinnerton-Dyer formula in rank $0$, and that for $M=\ad M_f(1)$
each term on the right corresponds to a term of (eq:identity) or of Remark rem:ns. Where does $h_\ell(E)$ appear, and why does $H^0(\Q,W)$ not appear in (eq:identity)?
\emph{Hint:} \Cref{ex:blochkato:bsd}; \Cref{prop:blochkato:global-invariants}(i).
\end{exercise}

\section{Notes and references}\label{sec:blochkato:notes}

\paragraph{Sources followed.} The Bloch-Kato local conditions, Selmer groups and the conjecture are from Bloch and Kato \cite{BlochKato} (Sections 3, 5 and 6) and
the reformulation of Fontaine and Perrin-Riou \cite{FPR94}; the propagation of local conditions and the comparison sequences follow Rubin \cite{Rubin} (Chapter 1) and
are proved, for this book, in \Cref{ch:galcoh}; local and global duality are in Milne \cite{MilneADT}, Neukirch, Schmidt and Wingberg \cite{NSW}, Tate \cite{TateDuality},
Greenberg \cite{Greenberg89} and Wiles \cite{Wiles95}. $p$-adic Hodge theory is quoted from Fontaine \cite{Fontaine}, Fontaine and Laffaille \cite{FontaineLaffaille},
Coleman and Iovita \cite{ColemanIovita99} and Berger's survey \cite{Berger}. The adjoint motive and the theorems on it are those of Diamond, Flach and Guo \cite{DFG04} and Flach
\cite{Flach92}; critical values are Deligne's \cite{Deligne79}; the symmetric square is due to Shimura \cite{Shimura76} and Gelbart and Jacquet \cite{GelbartJacquet}. The
structure of $E(\Q_p^{\ur})$ used in \Cref{prop:blochkato:tam-of-E} is in Silverman \cite{SilvermanATAEC} (Chapter IV) and Grothendieck \cite{SGA7}; the images of inertia
are due to Serre and Tate \cite{SerreTate68} and Kraus \cite{Kraus90}.

\paragraph{Theorem numbers.} We cite theorem numbers only where the paper cites them: \cite[Theorems 0.1, 2.7, 3.7, Lemma 2.1, (31), (57), p.~708]{DFG04} and
\cite[I.4.2.2]{FPR94}. All other references are to sections, and the section numbers given for Bloch and Kato (Sections 3 to 6), Rubin (Sections 1.3 and 1.7), Berger
(Section 2) and Silverman (Chapter IV) are pointers we have not checked against the printed texts. In particular we have not verified the numbers under which the
dimension formula of \Cref{prop:blochkato:dimension-formula}, the duality of \Cref{prop:blochkato:local-duality-f} at $\ell$ and the identification of the Kummer image at
$\ell$ (\Cref{ex:blochkato:classical-selmer}) appear in \cite{BlochKato}, nor the exact hypotheses of \cite{Flach92} (\Cref{thm:blochkato:flach} is stated in outline for that reason).

\paragraph{History.} Bloch and Kato formulated their conjecture in 1990 \cite{BlochKato}; Fontaine and Perrin-Riou recast it in terms of Galois cohomology and defined
the local Tamagawa factors used here \cite{FPR94}. Flach proved the first finiteness theorem for the Selmer group of the symmetric square in 1992 \cite{Flach92}, and
Diamond, Flach and Guo proved the $\ell$-part of the conjecture for adjoint motives of modular forms in 2004 \cite{DFG04}, using the Taylor-Wiles method.

\paragraph{What the book adds.} \Cref{prop:blochkato:tam-dual} proves $\Tam_p(T^D)=\Tam_p(A)$ for the adjoint of an elliptic curve in every case, where the paper cites \cite[(57)]{DFG04};
\Cref{prop:blochkato:tam-of-E} gives a proof that $\Tam_p(T_\ell E)$ is the $\ell$-part of $c_p$ at every $p\ne\ell$; \Cref{prop:blochkato:selmer-relations}(iii) and
\Cref{prop:blochkato:euler-characteristic} make explicit the step from Theorem thm:dfg27 to the finiteness of $H^1_f(\Q,W)$ and the equivalence used in Remark rem:ns. Lemma
lem:tamadd and Corollary cor:explicit, which answer the first half of the original Open question 3, are in the revised paper; \Cref{thm:blochkato:tam-additive} proves them
with the same argument.

\paragraph{Computations.} All numbers were computed with PARI/GP 2.17.2 \cite{PARI} through \texttt{cypari2}, on Cremona's files \texttt{allcurves}, \texttt{degphi} and
\texttt{galrep} \cite{ecdata} for conductors below $10{,}000$. Curves with CM were recognised by their $j$-invariant among the $13$ rational CM $j$-invariants and removed. The
optimal curve of a class is the one listed in \texttt{degphi}; irreducibility of $E[\ell]$ means that \texttt{galrep} lists no label $\ell\mathrm B$ or $\ell\mathrm{Cs}$; the image \texttt{3Ns} is
read from \texttt{galrep}. The terms $e_\ell(E,q)$ were computed from \texttt{lfuneuler(lfunsympow(E,2),q)} exactly as in the paper's function \texttt{rule}; the unit parts of the
Tate parameters as the reduction of $j_E^{-1}q^{v_q(j_E)}$, checked against \texttt{ellinit(E,O(q\^{}20)).tate} on \curve{11a1} and \curve{791a1}; $v_\ell(c(u))$ by testing
$\bar u^{(q-1)/\ell^k}=1$. The counts are those of \Cref{rem:blochkato:open-question-3} ($13{,}203$, $23{,}361$, $28{,}134$ curves), \Cref{ex:blochkato:sec-tam} and
\Cref{ex:blochkato:ex-small-range}; the check of \Cref{prop:blochkato:tam-dual} ($26{,}448$ parameter choices) and of $v_\ell(c_q)$ ($101{,}336$ pairs on the $17{,}834$ curves of conductor
below $3{,}000$) are described after \Cref{prop:blochkato:tam-dual} and \Cref{prop:blochkato:tam-of-E}; the coordinates of \Cref{prop:blochkato:ad-free-c3} and the outputs of
\texttt{rule} for \curve{121a1}, \curve{726a1}, \curve{791a1} ($[1,1]$, $[2,2]$, $[2,1]$) were recomputed. Nothing beyond conductor $10{,}000$ was checked.
